\documentclass[11pt]{ctexart}
\usepackage[a4paper,margin=1in]{geometry}
\usepackage{graphicx}
\usepackage{xcolor}
\usepackage{hyperref}
\definecolor{refgray}{gray}{0.45}
\title{\textbf{市场最新观点汇总}\\[0.4em]{\large 通胀不降息、AI资本开支超互联网泡沫、能源供给硬约束——市场定价逻辑正从基本面转向供给节奏与融资结构}}
\author{KC桌面 自动生成}
\date{260610}
\begin{document}
\maketitle
\tableofcontents
\newpage
\section{一页摘要}
\begin{itemize}
  \item 市场低估的不是通胀本身，而是“通胀但不加息”的定价逻辑，这对企业盈利和周期股构成正面推动。
  \item AI资本开支已超越互联网泡沫峰值，但市场尚未定价折旧的滞后冲击和资产负债表外承诺的风险。
  \item 全球能源与大宗商品市场面临供给硬约束，霍尔木兹海峡的例外通行无法改变LNG和原油的结构性缺口。
  \item 中国地产成交量回升但结构性分化加剧，二手房成为更可靠的市场温度计，新房反弹依赖供给节奏。
  \item AI产业增长引擎正从B2C转向B2B，企业级Agentic AI需求加速爆发，中国模型在成本效率上形成优势。
  \item 新兴市场基本面改善，但美元融资单一化成为最大约束，多元化融资货币可降低波动。
  \item 全球资金流向显示AI乐观情绪驱动美国权益流入，但新兴市场资金撤退信号值得警惕。
  \item 中国尿素出口政策从配额转向价格下限，全球化肥定价权正发生结构性转移。
  \item 存储行业正经历结构性重定价，NAND和HDD的后半场机会被低估，AI驱动需求阶梯式放大。
  \item CLO市场超额回报正从beta转向alpha，管理人选股能力成为首要变量。
\end{itemize}
\section{宏观与利率：通胀但不加息，市场定价逻辑的重构}
\textbf{当前市场对通胀的担忧过度聚焦于负面情景，而忽略了通胀在美联储不收紧的条件下对企业盈利和股票估值的正面推动。劳动力市场结构性再平衡、薪资增长放缓，为美联储按兵不动提供了空间。}

\begin{itemize}
  \item MS指出，非农就业三个月均值188k，失业率稳定在4.3\%，但薪资增长低于去年同期，单位劳动力成本压力减轻，这是典型的软着陆信号而非再通胀。
  \item MS认为，只要美联储不因通胀而加息，通胀本身对企业盈利反而是好事；PPI领先标普500营收增长约4个月，历史上盈利高增长+美联储中性环境下，标普500的12个月回报率中位数为14\%。
  \item NOM强调，当前通胀并非需求过热驱动，而是能源、芯片和食品三重供给冲击叠加，传统“通胀高则加息”的逻辑已失效，需区分谁有资格加息。
  \item GS分析中国5月通胀数据，PPI同比升至3.9\%，但82\%的贡献来自上游行业，核心CPI反而走软至1.1\%，表明成本挤压而非需求复苏。
  \item GS指出，日本5月CGPI环比增速从4月的2.8\%回落至0.9\%，四月跳升是一次性定价调整，五月才是判断通胀持续性的真实基线。
  \item BARC预测，即便霍尔木兹海峡恢复通行，2027年布伦特原油均价仍将达88美元/桶，高出远期曲线7美元，市场系统性低估了供给侧再定价。
  \item JPM认为，关税政策的真正转折不在税率高低，而在法律框架从IEEPA/Section 122向Section 301切换，实际税率仅6.7\%，豁免清单才是博弈场。
\end{itemize}
\begin{figure}[htbp]
\centering
\includegraphics[width=0.88\linewidth]{figures/fig_001.jpg}
\caption{Figure 1 - Figure 1: Global commodity flows recovered over the week, owing to metals-led inflows, partially offset by outflows from crude oil USD billion, cumulative flows across tracked commodity markets}
\end{figure}
\begin{figure}[htbp]
\centering
\includegraphics[width=0.88\linewidth]{figures/fig_240.jpg}
\caption{Exhibit 4 - Exhibit 1: Year-over-year PPI inflation increased further from April to May}
\end{figure}
\begin{figure}[htbp]
\centering
\includegraphics[width=0.88\linewidth]{figures/fig_243.jpg}
\caption{Exhibit 4 - Exhibit 4: PPI inflation increased further from April to May, mainly due to higher energy and related chemical prices Contribution to year-over-year PPI inflation}
\end{figure}
\begin{figure}[htbp]
\centering
\includegraphics[width=0.88\linewidth]{figures/fig_276.jpg}
\caption{Exhibit 1 - Exhibit 1: Inflation Indices and USD/JPY Rate}
\end{figure}
\begin{figure}[htbp]
\centering
\includegraphics[width=0.88\linewidth]{figures/fig_408.jpg}
\caption{FIGURE 1 - BCI, US FIGURE 1. Managed money positioning in oil is back to pre-war levels...}
\end{figure}
\begin{figure}[htbp]
\centering
\includegraphics[width=0.88\linewidth]{figures/fig_409.jpg}
\caption{FIGURE 2 - Note: Net speculative positioning in Brent and WTI futures and options combined, percentile rank based on data since 2014. FIGURE 2. ...so is excess implied vol in oil markets}
\end{figure}
{\small\color{refgray}\textbf{References:} [R001] 摩根斯坦利：市场真正低估的不是通胀本身，而是“通胀但不加息”的定价逻辑; [R004] NOM：全球市场正在为一场“滞胀式供给冲击”重新定价，但真正的分歧在于谁有资格加息; [R007] GS：中国通胀数据的真正信号不在 CPI 本身，而在 PPI 的结构性集中; [R013] GS：日本通胀的“二阶导”拐点比绝对水平更重要; [R026] BARC：市场低估了供给侧的持续收紧，而非单纯押注地缘风险; [R009] JPM：关税政策的真正转折不在税率高低，而在法律框架的切换}

\section{AI/算力：资本开支超越互联网泡沫，折旧冲击与融资结构成新焦点}
\textbf{AI基础设施投资规模不仅在总量上创纪录，更在资本密集度、融资结构和资产负债表外承诺三个维度超越市场认知。市场当前定价由供给预期驱动，折旧的滞后冲击和融资节奏的波动将成为下一阶段核心叙事。}

\begin{itemize}
  \item MS预测，2026-2028年超大规模云服务商资本开支与销售比将达36\%-44\%，超过2000年互联网泡沫时期32\%的峰值，且增量资本开支常伴随资本募集或创新融资安排。
  \item MS指出，主要AI参与者的长期采购承诺接近1万亿美元，未开始的租赁承诺超过8000亿美元，RPO近三年涨了3倍，但折旧将是下一个关注点，微软、甲骨文、Meta、谷歌未来三年累计折旧预计超5200亿美元。
  \item MS预计2026年全球AI相关债务发行量达5700亿美元，较2025年增长162\%，四月单月发行超740亿美元，但五月美国市场降温后，头部公司转向欧元、加元等非美元市场发行约240亿美元。
  \item BARC分析，AI资本开支正快速吞噬经营现金流，2028年五大云厂商总资本支出预计达1.16万亿美元，现金流覆盖比例可能逼近97\%，可转债和强制可转债成为战略融资工具。
  \item HSBC将2026-2030年全球AI行业累计收入预期上调38\%至2.55万亿美元，但B2B收入预期上调74\%，B2C下调20\%，企业级Agentic AI是核心驱动力。
  \item JEF数据显示，中国模型在OpenRouter平台周Token消耗首次超越美国模型，DeepSeek V4 Flash单周消耗3.69万亿Token，中国模型API成本仅为美国几分之一，成本效率优势改写竞争格局。
  \item MS报告指出，GB200/300机柜5月出货7700台环比下降7\%，但全年预计超7-8万台同比翻倍以上，月度波动是供应链节奏问题而非需求拐点。
\end{itemize}
\begin{figure}[htbp]
\centering
\includegraphics[width=0.88\linewidth]{figures/fig_397.jpg}
\caption{Exhibit 1 - Exhibit 1: MSe expects Hyperscaler capex/sales to surpass dot-com peaks reaching 44\% in '27}
\end{figure}
\begin{figure}[htbp]
\centering
\includegraphics[width=0.88\linewidth]{figures/fig_399.jpg}
\caption{Exhibit 3 - Exhibit 3: Leases not yet started have hit}
\end{figure}
\begin{figure}[htbp]
\centering
\includegraphics[width=0.88\linewidth]{figures/fig_400.jpg}
\caption{Exhibit 4 - Exhibit 4: ...while purchase commitments reach close to \textbackslash{}\$1tn}
\end{figure}
\begin{figure}[htbp]
\centering
\includegraphics[width=0.88\linewidth]{figures/fig_590.jpg}
\caption{FIGURE 1 - \#\# The AI Capex Cycle Begins to Overwhelm Operating Cash Flows The AI investment cycle continues to defy traditional capital intensity frameworks. Following the latest spending commitments, we believe peak capex has been deferred further out, reflecting both p}
\end{figure}
\begin{figure}[htbp]
\centering
\includegraphics[width=0.88\linewidth]{figures/fig_593.jpg}
\caption{FIGURE 4 - BARC Estimates Estimates from BARC' Internet research team. Note that for AMZN, capex is AWS only, while OCF is total AMZN. FIGURE 4. ...which is not baked into consensus estimates Capex/OCF}
\end{figure}
\begin{figure}[htbp]
\centering
\includegraphics[width=0.88\linewidth]{figures/fig_344.jpg}
\caption{Exhibit 1 - Exhibit 1 - Weekly token consumption over the past 12 months as of 8 Jun (tn)}
\end{figure}
\begin{figure}[htbp]
\centering
\includegraphics[width=0.88\linewidth]{figures/fig_403.jpg}
\caption{Exhibit 1 - Exhibit 1: Industry-wide GB200/300 NVL72-equivalent monthly rack output, by major ODMs (2025) GB200/300 NVL72 racks by major ODMs (000s)}
\end{figure}
{\small\color{refgray}\textbf{References:} [R023] 摩根斯坦利：AI资本开支已超越互联网泡沫，但市场尚未定价折旧的滞后冲击; [R044] 摩根斯坦利：AI融资的“夏季攻势”已提前到来，市场定价权正从基本面转向供给节奏; [R040] 0005-41-BARC-U.S. Equity-Linked Strategies AI Capex Funding- Why Equity- Why Converts--260609; [R016] HSBC：AI产业真正被低估的不是消费者市场，而是企业级需求的加速爆发; [R022] JEF：AI竞争已从模型性能转向Agent生态与成本效率的再定价; [R024] 摩根斯坦利：GB200/300机柜出货量环比下降7\%，但这轮调整的真正含义是供应链议价权的再分配}

\section{能源与大宗：供给硬约束重塑定价权，中国出口成金属市场泄压阀}
\textbf{全球能源和大宗商品市场正经历供给侧结构性收紧，霍尔木兹海峡的例外通行无法改变LNG和原油的结构性缺口。中国从需求方转变为供给调节器，尿素出口政策从配额转向价格下限，金属出口成为全球短缺的泄压阀。}

\begin{itemize}
  \item JPM指出，霍尔木兹海峡内仅有9艘可用LNG船，卡塔尔2026年全年LNG产量将从2025年的1.05亿吨腰斩至0.6亿吨，JKM对TTF的溢价将持续存在。
  \item BARC预测，即便海峡恢复通行，2026年布伦特原油均价100美元/桶，全球库存以450万桶/天速度下降，美国商业原油库存逼近2008年低点，库欣储油设施利用率降至2010年以来最低。
  \item BARC提出，中国4月石油需求同比下降12\%可能并非价格弹性，而是主动抑制采购、暗中储备的战略行为，可见库存自冲突以来基本未变。
  \item MS分析，中国尿素出口从配额管理转向价格下限，大颗粒尿素不低于500美元/吨FOB，这是从量控到价控的范式转换，全球化肥定价权正发生结构性转移。
  \item JPM报告显示，中国铜表观消费4月同比增长9\%，但终端消费加权指标同比下降4\%，背离源于库存去化和电网投资支撑，真实需求并未如价格反映的那么强劲。
  \item JPM指出，中国铝出口套利窗口打开，半成品出口加速，LME铝库存降至33万吨而中国库存仍有140万吨，中国出口成为全球金属市场结构性短缺的泄压阀。
  \item MS认为，全球矿业正经历主权供需的结构性重塑，中国CMRG集中采购与印尼DSI集中出口同步崛起，澳大利亚矿企面临竞争法约束，若建立国家出口平台将彻底改变定价模型。
\end{itemize}
\begin{figure}[htbp]
\centering
\includegraphics[width=0.88\linewidth]{figures/fig_278.jpg}
\caption{Figure 1 - Figure 1: Recent LNG tanker movements across the Strait of Hormuz Vessels All Types Vessel History AL DRAYEN}
\end{figure}
\begin{figure}[htbp]
\centering
\includegraphics[width=0.88\linewidth]{figures/fig_279.jpg}
\caption{Figure 3 - Figure 3: USGC LNG exports to Europe/Mediterranean}
\end{figure}
\begin{figure}[htbp]
\centering
\includegraphics[width=0.88\linewidth]{figures/fig_410.jpg}
\caption{Figure 3 - Note: OVX minus VIX in standard deviation from mean since 2020 We reiterate that this is not a new equilibrium. The Strait cannot remain closed in perpetuity with oil prices at \textbackslash{}\$100/b. Oil inventories continue to draw down, although the pace has not accelerat}
\end{figure}
\begin{figure}[htbp]
\centering
\includegraphics[width=0.88\linewidth]{figures/fig_411.jpg}
\caption{FIGURE 4 - Note: Our weekly global total oil inventory indicator excluding the Middle East Gulf region minus the pre-pandemic seasonal average (2017-19), mb FIGURE 4. ...but the pace has not accelerated because of the pull back in Chinese demand}
\end{figure}
\begin{figure}[htbp]
\centering
\includegraphics[width=0.88\linewidth]{figures/fig_465.jpg}
\caption{Figure 53 - Figure 53: Global visible aluminium inventories}
\end{figure}
\begin{figure}[htbp]
\centering
\includegraphics[width=0.88\linewidth]{figures/fig_468.jpg}
\caption{Figure 59 - Figure 59: Chinese unwrought aluminium and aluminium product exports, 3-month rolling average}
\end{figure}
\begin{figure}[htbp]
\centering
\includegraphics[width=0.88\linewidth]{figures/fig_638.jpg}
\caption{Exhibit 1 - Exhibit 1: Australia Exports a Significant Share of Key Commodities Share of Global Seaborne Supply}
\end{figure}
{\small\color{refgray}\textbf{References:} [R014] JPM：霍尔木兹海峡的“例外通行”不会改变全球LNG市场的结构性缺口; [R026] BARC：市场低估了供给侧的持续收紧，而非单纯押注地缘风险; [R030] 摩根斯坦利：中国尿素出口底价的重新设定，是全球化肥定价逻辑的一次结构性重置; [R031] 摩根斯坦利：中国取消尿素出口底价，全球化肥市场的定价权正发生结构性转移; [R028] JPM：中国出口正在成为全球金属市场结构性短缺的“泄压阀”; [R042] 摩根斯坦利：全球矿业正在经历一场“主权供需”的结构性重塑}

\section{地产与消费：成交量回升但结构性分化，二手房成更可靠温度计}
\textbf{中国房地产市场成交量回升并非全面复苏，而是结构性分化加速——一线城市二手房成为资金避风港，新房反弹依赖供给节奏而非需求自发修复。香港房价已接近全年预测上限，但定价策略正从折让转向溢价。}

\begin{itemize}
  \item MS数据显示，截至6月7日当周，50城新房注册成交量同比增长23\%，10城二手房成交量同比增幅从9\%扩大至23\%，但年初至今新房累计仍下降12\%。
  \item MS强调，二手房成交量正成为比新房更可靠的市场温度计，一线城市二手房成交同比增幅从前一周的12\%跃升至34\%，二线城市从6\%提升至17\%。
  \item JPM报告显示，香港住宅市场年初至今已涨9.6\%，接近全年预测区间10\%-15\%的上沿，但一手楼定价从折让转向溢价，Pavilia Rosa均价高出二手市场41\%。
  \item JPM指出，香港高端地段一手楼敢于定出显著溢价，中端项目仍在谨慎试水，市场正经历K型复苏，与深圳高端住宅销售强劲、中端疲软的分化现象一致。
  \item JPM数据显示，5月一线城市二手房挂牌量环比降0.7\%，库存月数从18.8个月微降到18.6个月，北京降幅最大（-2\%），上海库存月数从12.4降到11.9。
  \item GS分析澳大利亚消费者信心连续三个月低迷，6月指数降至80.6，家庭财务状况感知全面恶化，住房价格预期三年来首次跌破长期均值，悉尼和墨尔本跌幅显著。
\end{itemize}
\begin{figure}[htbp]
\centering
\includegraphics[width=0.88\linewidth]{figures/fig_343.jpg}
\caption{Exhibit 1 - Exhibit 1: Primary weekly unit sales and four-week moving average – Total 50 cities}
\end{figure}
\begin{figure}[htbp]
\centering
\includegraphics[width=0.88\linewidth]{figures/fig_327.jpg}
\caption{Figure 5 - Figure 5: 12-city daily secondary sales registrations}
\end{figure}
\begin{figure}[htbp]
\centering
\includegraphics[width=0.88\linewidth]{figures/fig_333.jpg}
\caption{Figure 11 - Figure 11: Hong Kong secondary home prices (Centa-city Leading Index, or CCL)}
\end{figure}
\begin{figure}[htbp]
\centering
\includegraphics[width=0.88\linewidth]{figures/fig_323.jpg}
\caption{Figure 1 - Figure 1: Centraline secondary asking price index vs. NBS secondary home price index M/M in tier-1 cities}
\end{figure}
\begin{figure}[htbp]
\centering
\includegraphics[width=0.88\linewidth]{figures/fig_416.jpg}
\caption{Exhibit 1 - Exhibit 1: Consumer sentiment fell in June reversing the small rise in May}
\end{figure}
{\small\color{refgray}\textbf{References:} [R019] 摩根斯坦利：市场低估的并非政策力度，而是成交量回升背后的结构性分化; [R018] JPM：香港房价已接近全年预测上限，但真正的不确定性不在价格本身; [R027] GS：澳大利亚消费者信心连续三个月低迷，但市场真正该关注的不是情绪本身}

\section{区域市场：新兴市场融资货币单一化，亚洲资金撤退信号值得警惕}
\textbf{新兴市场基本面持续改善，但美元走强暴露了融资结构的脆弱性。全球资金流向显示AI乐观情绪驱动美国权益流入，但亚洲市场的资金撤退可能才刚刚开始，人民币中间价模型释放稳定信号。}

\begin{itemize}
  \item MS提出反直觉结论：新兴市场问题不在于是否持有资产，而在于用什么货币融资；以加元或G3篮子货币融资可在不牺牲收益的情况下显著降低波动。
  \item NOM数据显示，截至6月5日的五个交易日内，外国投资者向美国权益基金净投入25.6亿美元，5月总额达87.8亿美元创今年1月以来新高，但新兴市场ETF同期净流出2.5亿美元。
  \item NOM指出，AI乐观情绪是驱动美国权益流入的唯一支柱，一旦叙事出现裂痕，资金流出速度可能比流入更快；韩国散户在连续两个月净卖出后开始回补，台湾投资者加速抛售美元债券。
  \item NOM模型显示，人民币兑美元中间价预测值下移至6.7775，较前次大幅下移372个基点，逆周期因子在平滑波动，定价机制本身释放稳定性信号。
  \item MS分析，中国新ODI框架对全球银行财富管理业务的冲击被低估，中国内地居民定义比预期宽泛，持有香港身份证但未注销内地户籍的高净值人士将被纳入监管。
  \item JPM指出，中国新ODI框架下，境外IPO套现、境外子公司收入可能纳入监管，HSBC和渣打的财富收入2025年同比增长24\%，合规成本系统性上升。
\end{itemize}
\begin{figure}[htbp]
\centering
\includegraphics[width=0.88\linewidth]{figures/fig_301.jpg}
\caption{Exhibit 1 - Exhibit 1: EM gains vs G3 expected}
\end{figure}
\begin{figure}[htbp]
\centering
\includegraphics[width=0.88\linewidth]{figures/fig_302.jpg}
\caption{Exhibit 2 - Exhibit 2: Diversified funding reduces risk Expected Total Returns / Vol}
\end{figure}
\begin{figure}[htbp]
\centering
\includegraphics[width=0.88\linewidth]{figures/fig_275.jpg}
\caption{Figure 3 - - Foreign inflows into US-focused equity ETFs and mutual funds[1] totaled USD2.6bn between 1 and 5 June (Figure 3), and compared with inflows of USD2.9bn over the previous week. In May, inflows into these funds totaled USD8.8bn, the largest inflows since USD9.}
\end{figure}
\begin{figure}[htbp]
\centering
\includegraphics[width=0.88\linewidth]{figures/fig_306.jpg}
\caption{Exhibit 7 - Exhibit 7: Debt flows have somewhat stabilized Net Monthly FPI in Debt (US\$ bn)}
\end{figure}
{\small\color{refgray}\textbf{References:} [R015] 摩根斯坦利：新兴市场真正的风险不是基本面，而是融资货币的单一化; [R011] NOM：全球资金正在重新定价“美国例外论”，但亚洲的撤退信号更值得警惕; [R025] NOM：人民币中间价模型正在发出一个被市场忽视的稳定信号; [R012] JPM：中国新ODI框架对全球银行的冲击比预期更广，财富管理业务的增量摩擦才刚刚开始}

\section{AI产业链：从GPU短缺到存储与计算架构再定价，企业级需求加速爆发}
\textbf{AI基础设施需求正从GPU短缺转向存储与计算架构的再定价。推理和Agentic AI正在创造传统计算的新需求周期，存储行业结构性重定价远未完成，NAND和HDD的后半场机会被低估。}

\begin{itemize}
  \item BofA指出，推理和Agentic AI正在将部分计算从并行GPU处理转移到顺序CPU处理，CPU密集型服务器和企业基础设施需求被重新激活，Dell和HPE报告传统计算需求显著增长。
  \item BofA上调IBM、NetApp、希捷和西部数据目标价，逻辑从卖硬件转向卖架构与锁定收益，存储成了新瓶颈，WDC预计HDD收入未来3-5年增速可能超过25\%。
  \item Bernstein认为，AI对存储的需求是逐级放大的阶梯，从训练到Agentic AI，NAND和HDD需求将从高升至极端，市场对供给侧的长期结构性约束定价不足。
  \item Bernstein指出，HDD行业首次同时实现整合与增长，长协条款越来越友好，供应商在定价、预付款和照付不议条款上获得更多保障，周期性被系统性削弱。
  \item JPM分析，SOCAMM从192GB降至96GB并非需求萎缩，而是供应紧张下的被动重组，NVIDIA与SK海力士签署多年技术合作协议，存储器产业从周期博弈进入结构锁定。
  \item BARC报告显示，企业级服务器收入2026年预计增长33\%（此前预期仅2\%），戴尔传统服务器业务增长超60\%，慧与订单三位数增长，但2027年可能只剩1\%增长。
  \item MS指出，Vera CPU单颗搭载1.5TB LPDDR5X，是上一代的三倍，以200-400万颗产量估算，将消耗全球DRAM bit总需求的10\%-16\%，存储设备订单可预见性远高于市场共识。
\end{itemize}
\begin{figure}[htbp]
\centering
\includegraphics[width=0.88\linewidth]{figures/fig_575.jpg}
\caption{EXHIBIT 3 - EXHIBIT 3: Memory/Storage Hierarchy for AI At the heart of AI is data, which resides in Memory and Storage}
\end{figure}
\begin{figure}[htbp]
\centering
\includegraphics[width=0.88\linewidth]{figures/fig_576.jpg}
\caption{EXHIBIT 4 - EXHIBIT 4: AI Demand for Memory \& Storage AI is driving unprecedented demand surge in Memory \& Storage}
\end{figure}
\begin{figure}[htbp]
\centering
\includegraphics[width=0.88\linewidth]{figures/fig_577.jpg}
\caption{EXHIBIT 5 - EXHIBIT 5: Data Center Mix for NAND, DRAM, and HDD \#\# Less price sensitive Data Center customers now dominate demand NAND}
\end{figure}
\begin{figure}[htbp]
\centering
\includegraphics[width=0.88\linewidth]{figures/fig_539.jpg}
\caption{FIGURE 1 - - Our new Cloud server revenue estimates are \textbackslash{}\$567Bn/\textbackslash{}\$885Bn in CY26/27 vs our prior estimates of \textbackslash{}\$585Bn/\textbackslash{}\$817Bn, respectively. We model growth of 77\% and 56\% in 2026/2027 vs our prior estimates of 84\% and 40\%. xPU server estimates and updated segmentation fr}
\end{figure}
\begin{figure}[htbp]
\centering
\includegraphics[width=0.88\linewidth]{figures/fig_540.jpg}
\caption{FIGURE 2 - We expect growth from Cloud and Enterprise servers moving through CY26-27. We find the Cloud growth impressive, given the large base it's growing from. Additionally, our Enterprise server growth this year accounts for the ASP increases along with continued vol}
\end{figure}
\begin{figure}[htbp]
\centering
\includegraphics[width=0.88\linewidth]{figures/fig_833.jpg}
\caption{Exhibit 1 - Exhibit 1: Vera CPU alone could account for more than \$10\textbackslash{}\%\$ of total global DRAM bit demand for 2025 LPDDR5X bit demand for Vera CPU relative to 2025 total DRAM bit demand (=38TB)}
\end{figure}
{\small\color{refgray}\textbf{References:} [R006] 美国银行：AI基础设施需求正从“GPU短缺”转向“存储与计算架构的再定价”; [R038] Bernstein：存储行业的结构性重定价远未完成，真正的机会在NAND和HDD的后半场; [R034] JPM：市场误读了SOCAMM降配，真正的机会藏在供给硬约束和NVIDIA的生态锁定里; [R032] BARC：企业级服务器的意外走强，正在改写服务器市场的增长叙事; [R048] 摩根斯坦利：市场低估的不是半导体设备周期性，而是结构性瓶颈的再定价}

\section{新能源与汽车：中国车企欧洲优势不在价格，储能市场供给端结构性收紧}
\textbf{中国车企在欧洲的份额增长已非简单的便宜BEV出口故事，零售执行与生态系统成为护城河。储能市场真正的超预期不在需求，而在供给端产能利用率从69\%跃升至94\%带来的定价权转移。}

\begin{itemize}
  \item JPM实地调研显示，2026年4月中国品牌在欧洲新能源车市场份额达18.3\%，是去年同期的三倍以上，PHEV在BYD门店订单中占比高达50\%-60\%。
  \item JPM指出，中国车企竞争优势来自多动力总成组合+品牌阶梯的产品矩阵，比亚迪计划到2026年底在欧洲建3000个超快充桩，充电生态成为下一个护城河。
  \item Citi分析，5月新能源乘用车批发135万辆同比+10.6\%，但零售仅95万辆同比-7.5\%，批发与零售40万辆的差距暴露渠道库存深层博弈，新能源出口占比首次突破54\%。
  \item Bernstein报告显示，4月中国储能电池出货量93.1GWh，同比增长92\%，产能利用率从69\%跃升至94\%，行业从产能过剩进入供需紧平衡，定价逻辑从量的增长切换到利润率的修复。
  \item JEF预计2026年全球储能装机超150GWh同比增长80\%，美国AI数据中心储能需求2026年翻倍至10GWh，到2030年将增长至70-80GWh，储能正从光伏附属品变为独立增长极。
  \item JPM上调比亚迪和吉利2026年海外销量预测约15\%和30\%，基于产品组合策略、零售执行能力和充电生态建设三个被低估的结构性变量。
\end{itemize}
\begin{figure}[htbp]
\centering
\includegraphics[width=0.88\linewidth]{figures/fig_317.jpg}
\caption{Figure 1 - Figure 1: Chinese OEMs' annual market share in the NEV space in Europe}
\end{figure}
\begin{figure}[htbp]
\centering
\includegraphics[width=0.88\linewidth]{figures/fig_318.jpg}
\caption{Figure 3 - Figure 3: BYD market share in Europe's NEV market (8\% in Apr-26)}
\end{figure}
\begin{figure}[htbp]
\centering
\includegraphics[width=0.88\linewidth]{figures/fig_244.jpg}
\caption{Exhibit 2 - EXHIBIT 2: Monthly China ESS batteries shipments 93.1GWh (92\% YoY, 14\% MoM) China ESS shipment (GWh)}
\end{figure}
\begin{figure}[htbp]
\centering
\includegraphics[width=0.88\linewidth]{figures/fig_248.jpg}
\caption{EXHIBIT 8 - EXHIBIT 8: Monthly China ESS manufacturing capacity utilization rate China Manufacturing Capacity Utilization Rate}
\end{figure}
\begin{figure}[htbp]
\centering
\includegraphics[width=0.88\linewidth]{figures/fig_251.jpg}
\caption{Exhibit 13 - EXHIBIT 13: Monthly China ESS tender (GWh)}
\end{figure}
{\small\color{refgray}\textbf{References:} [R017] JPM：中国车企在欧洲的真正优势不是价格，而是零售执行与生态系统的护城河; [R033] Citi：市场低估的不是新能源渗透率，而是出口定价权的结构性重构; [R008] Bernstein：储能市场真正的超预期不在需求，而在供给端正在发生的结构性收紧; [R021] JEF：储能市场的真正拐点不是增长，而是增长的结构性分化}

\section{风险偏好：CLO市场从beta转向alpha，军工复合体面临工业基础结构性过时}
\textbf{信用周期进入后半段，CLO市场超额回报正从市场方向转向管理人能力。美国军工复合体面临的不是预算问题，而是工业基础的结构性过时，低成本可消耗系统正在改变战场规则。}

\begin{itemize}
  \item MS指出，当前信用利差在多数层级偏紧，AI颠覆风险导致抵押品表现离散度加大，投资者越来越关注管理人特有的alpha而非广泛的市场beta。
  \item MS报告显示，5月CLO新发行量从4月的60亿美元反弹至约170亿美元，再融资与重置活动达390亿美元，但投资者叙事已从买什么区域转向买谁管理的资产。
  \item JPM分析，美国国防工业基础面临四个结构性漏洞：供应链太浅、采购太慢、行业太集中、供应链太脆弱，关键材料和元器件对中国依赖度高。
  \item JPM指出，冷战后的旗舰崇拜让美国军工体系失去规模弹性，主承包商从50多家缩到5家，控制60\%以上合同额，而乌克兰和伊朗战争暴露了低成本可消耗系统的重要性。
  \item MS认为，欧洲猎头行业临时招聘回暖不能掩盖永久招聘的结构性恶化，Page Group有72\%的毛利来自永久招聘，是同行中暴露度最高的。
  \item MS数据显示，Adecco的永久招聘占毛利润比例约15-17\%，Randstad约16\%，Hays达38\%，Page Group达72\%，永久招聘萎缩正在侵蚀毛利率和经营利润。
\end{itemize}
\begin{figure}[htbp]
\centering
\includegraphics[width=0.88\linewidth]{figures/fig_597.jpg}
\caption{Exhibit 1 - Exhibit 1: PC vs BSL AAAs widened to 32bp over the 3 months ending in May}
\end{figure}
\begin{figure}[htbp]
\centering
\includegraphics[width=0.88\linewidth]{figures/fig_602.jpg}
\caption{Exhibit 11 - Exhibit 11: Monthly CLO Creation (New Issue and Refis/Resets)}
\end{figure}
\begin{figure}[htbp]
\centering
\includegraphics[width=0.88\linewidth]{figures/fig_509.jpg}
\caption{Figure 1 - Figure 1: State-based conflicts by region}
\end{figure}
\begin{figure}[htbp]
\centering
\includegraphics[width=0.88\linewidth]{figures/fig_510.jpg}
\caption{Figure 2 - Figure 2: Global defense spending}
\end{figure}
\begin{figure}[htbp]
\centering
\includegraphics[width=0.88\linewidth]{figures/fig_836.jpg}
\caption{Exhibit 1 - Exhibit 1: Staffers' organic growth has recovered, mainly driven by temp...}
\end{figure}
\begin{figure}[htbp]
\centering
\includegraphics[width=0.88\linewidth]{figures/fig_837.jpg}
\caption{Exhibit 3 - Exhibit 3: Indeed job postings have continued to deteriorate... Exhibit 5: Staffing exposure to permanent recruitment (\% of gross profit)}
\end{figure}
{\small\color{refgray}\textbf{References:} [R041] 摩根斯坦利：CLO市场真正的分化不在区域，而在管理人的议价权; [R029] JPM：美国军工复合体的真正危机不在于预算，而在于工业基础的结构性过时; [R049] 摩根斯坦利：欧洲猎头行业的真正风险不在需求，而在永久招聘的结构性萎缩}

\section{AI竞争格局：从模型性能转向Agent生态与成本效率，中国软件更具韧性}
\textbf{AI竞争已从模型性能转向Agent生态与成本效率的再定价。中国模型在API成本上仅为美国几分之一，Agent-to-Agent生态化协作正在替代单一模型的能力竞赛。市场对AI颠覆软件的恐惧过度，中国软件公司反而更具韧性。}

\begin{itemize}
  \item JEF数据显示，截至6月1日当周，OpenRouter平台周Token消耗环比增长13.5\%至36.1万亿，中国模型周Token消耗首次超越美国模型，DeepSeek V4 Flash单周消耗3.69万亿Token。
  \item JEF指出，腾讯正搭建Agent-to-Agent生态，元宝App已接入美团AI助手，手机厂商OPPO、小米、华为、vivo、荣耀也已接入，用户通过语音助手可直接发微信或启动音视频通话。
  \item HSBC报告显示，Anthropic在B2B领域的ARR已达470亿美元，超过OpenAI的300亿，OpenAI砍掉Sora等副线项目全力押注编程和企业用户。
  \item JEF认为，AI对软件行业是自然选择而非末日，中国软件公司因按项目或模块收费、国企客户偏好定制化+本地部署，对AI冲击的免疫力更强。
  \item JEF指出，中国软件板块年初至今下跌13\%，主因是盈利预期下修和估值双杀而非AI颠覆本身，工业软件、金融IT、ERP三类软件最不容易被AI替代。
  \item MS分析，北美互联网板块估值离散度急剧扩大，电商板块市值加权平均下跌8.9\%，数字广告仅下跌3.8\%，市场在独立评估每个公司AI投资的回报周期。
\end{itemize}
\begin{figure}[htbp]
\centering
\includegraphics[width=0.88\linewidth]{figures/fig_344.jpg}
\caption{Exhibit 1 - Exhibit 1 - Weekly token consumption over the past 12 months as of 8 Jun (tn)}
\end{figure}
\begin{figure}[htbp]
\centering
\includegraphics[width=0.88\linewidth]{figures/fig_346.jpg}
\caption{Exhibit 4 - Exhibit 4 - Market share: weekly token consumption by company as of 8 Jun}
\end{figure}
\begin{figure}[htbp]
\centering
\includegraphics[width=0.88\linewidth]{figures/fig_389.jpg}
\caption{Exhibit 79 - Exhibit 79 - ARR of Anthropic}
\end{figure}
\begin{figure}[htbp]
\centering
\includegraphics[width=0.88\linewidth]{figures/fig_545.jpg}
\caption{Exhibit 1 - Exhibit 1 - SW Baskets' EV/S - US vs CN}
\end{figure}
\begin{figure}[htbp]
\centering
\includegraphics[width=0.88\linewidth]{figures/fig_547.jpg}
\caption{Exhibit 3 - Exhibit 3 - Ranking of China SW AI Defensibility}
\end{figure}
\begin{figure}[htbp]
\centering
\includegraphics[width=0.88\linewidth]{figures/fig_568.jpg}
\caption{Exhibit 4 - Exhibit 4: NTM EV/EBITDA for AMZN/GOOGL/META vs. Historical Averages Exhibit 5: Internet names fell -5\% last week (SPX/NDX -3/-5\%)}
\end{figure}
{\small\color{refgray}\textbf{References:} [R022] JEF：AI竞争已从模型性能转向Agent生态与成本效率的再定价; [R016] HSBC：AI产业真正被低估的不是消费者市场，而是企业级需求的加速爆发; [R035] JEF：市场对AI颠覆软件的恐惧过度，中国软件公司反而更具韧性; [R036] 摩根斯坦利：市场正在用不同估值逻辑为互联网公司定价，但真正的分歧不在AI硬件与软件之间}

\section{特高压与电网设备：订单翻倍背后，竞争从订单量转向议价权}
\textbf{国家电网2026年第二批特高压设备招标总金额94亿元，是第一批的两倍有余，已接近2025年全年总额的43\%。随着招标量放大、铜价上行、交付节奏加快，真正赢家是能在规模扩张中守住毛利率的企业。}

\begin{itemize}
  \item Citi数据显示，2026年第二批特高压设备招标94亿元，上半年累计已接近去年全年，电网投资从脉冲式放量进入持续高强度建设周期。
  \item Citi指出，平高电气以20.92亿元中标额排名第一，占22.3\%，中国西电18.99亿元占20.2\%，特变电工4.73亿元占5.0\%，前三名占了近58\%，集中度很高。
  \item Citi分析，1000kV GIS占招标总额的57.4\%，是特高压核心设备，技术壁垒高、毛利率可观，平高预计2026年还将有至少2个新的交流特高压项目获批。
  \item Citi援引数据，2026年一季度全国电网资本开支1675亿元，同比增40\%，增量主要来自土建工程，政府需要拉动固投转正。
  \item MS分析，中国数据中心市场真正被低估的不是2万亿投资规模，而是电力需求的非线性跃迁，NEA预计到2030年数据中心用电量达800TWh，占全国总用电量的6\%。
  \item GS指出，中国数据中心用电量将以36\%的年复合增速增长，超过美国预测的726TWh，中国从跟随者变为引领者，电力需求非线性增长逻辑尚未被充分定价。
\end{itemize}
\begin{figure}[htbp]
\centering
\includegraphics[width=0.88\linewidth]{figures/fig_589.jpg}
\caption{Figure 1 - Not for distribution in the People's Republic of China, excluding the Hong Kong Special Administrative Region and Qualified Foreign Institutional Investors. comprising Rmb647m won by Pinggao which ranked No. 2 among all companies participating in the biddings.}
\end{figure}
\begin{figure}[htbp]
\centering
\includegraphics[width=0.88\linewidth]{figures/fig_106.jpg}
\caption{Exhibit 18 - Exhibit 18: China's Electricity Output \& Consumption}
\end{figure}
{\small\color{refgray}\textbf{References:} [R039] Citi：特高压订单翻倍背后的关键信号——电网设备竞争已从“订单量”转向“议价权”; [R020] GS：中国数据中心市场真正被低估的不是投资规模，而是电力需求的非线性跃迁}

\section{化工与材料：合成橡胶供给缺口重塑定价权，松下投资者日揭示AI电源机遇}
\textbf{合成橡胶市场的超额利润并非来自需求反弹，而是原料石脑油/LPG短缺引发的供给刚性收缩。AI基础设施的物理层正经历从够用到必须重构的转折，松下等传统电子巨头在电源与材料环节占据被低估的位置。}

\begin{itemize}
  \item JPM报告显示，韩国NBL价格5月攀升至1633美元/吨，4-5月均价创五年新高，NBL加工利润自2022年以来首次超越SBR，达到916美元/吨。
  \item JPM指出，原料短缺是主因，中东局势导致石脑油供应紧张，锦湖石化二季度NBL产量预计环比下降超20\%，马来西亚手套厂因原料短缺减产。
  \item JPM分析，松下对数据中心储能系统目标明确，FY2028年实现约1万亿日元销售收入，FY2026指引仅为5500亿，电容器备用单元市场预计FY2030到1000亿日元。
  \item JPM指出，松下行业板块AI相关业务目标FY2028达4300亿日元，核心产品是先进基板材料和导电聚合物电容，客户已在询问M9/M10级别产品。
  \item MS分析，中国尿素出口底价取消暴露了出口策略的真正优先级，维持出口通道畅通、消化国内过剩产能比短期价格谈判权更重要。
  \item MS指出，中国化肥出口政策从开关式管控进入精细化、选择性、与地缘风险高度耦合的新阶段，中国正在成为供给侧的主动定价变量。
\end{itemize}
\begin{figure}[htbp]
\centering
\includegraphics[width=0.88\linewidth]{figures/fig_830.jpg}
\caption{Figure 1 - Figure 1: JPM NBL price forecast (\textbackslash{}\$/t) Figure 2: Korea NBL price (\textbackslash{}\$/t)}
\end{figure}
\begin{figure}[htbp]
\centering
\includegraphics[width=0.88\linewidth]{figures/fig_831.jpg}
\caption{Figure 3 - Figure 3: Korea NBL export volume (kt)}
\end{figure}
{\small\color{refgray}\textbf{References:} [R043] JPM：合成橡胶的供给缺口正在重塑化工定价权，而市场仍在关注需求侧; [R045] JPM：AI基础设施的下一轮投资机会，藏在松下投资者日被低估的细节里; [R031] 摩根斯坦利：中国取消尿素出口底价，全球化肥市场的定价权正发生结构性转移}

\section{黄金与矿业：黄金股回调是表象，矿业主权供需重塑竞争格局}
\textbf{黄金股当前价格错误地定价了公司未来两年持续改善的自由现金流，而非金价短期波动。全球矿业供需两端正被主权力量重新组织，国家正在成为最重要的交易对手。}

\begin{itemize}
  \item Citi分析，金价从1月高点回落约11\%，但黄金股修正高达25\%，市场在为金价继续下跌支付过高保险溢价，而Q1运营数据稳定，现金流依然强劲。
  \item Citi指出，AngloGold自由现金流收益率10.5\%，估值4.7倍EBITDA，Gold Fields自由现金流收益率12\%，估值3.6倍EBITDA，均低于全球同行和历史均值。
  \item MS报告显示，中国CMRG集中采购与印尼DSI集中出口同步崛起，澳大利亚占全球铁矿石海运出口59\%、焦煤38\%，若建立国家出口平台将彻底改变定价模型。
  \item MS指出，全球关键矿产供需两端被主权力量重新组织，过去市场习惯将矿业视为自由市场定价典范，但未来国家正在成为最重要的交易对手。
  \item JPM分析，美国军工复合体工业基础结构性过时，冷战后的旗舰崇拜让体系失去规模弹性，低成本可消耗系统和实时软件整合正在改变战场规则。
  \item JPM指出，美国国防工业基础供应链太浅，二级三级供应商减少成了瓶颈，采购太慢，重大武器项目动辄8年起步，行业太集中，关键材料对中国依赖度高。
\end{itemize}
\begin{figure}[htbp]
\centering
\includegraphics[width=0.88\linewidth]{figures/fig_638.jpg}
\caption{Exhibit 1 - Exhibit 1: Australia Exports a Significant Share of Key Commodities Share of Global Seaborne Supply}
\end{figure}
\begin{figure}[htbp]
\centering
\includegraphics[width=0.88\linewidth]{figures/fig_639.jpg}
\caption{Exhibit 5 - Exhibit 5: EBITDA margins of ASX miners in our coverage CY2026 EBITDA Margin}
\end{figure}
\begin{figure}[htbp]
\centering
\includegraphics[width=0.88\linewidth]{figures/fig_509.jpg}
\caption{Figure 1 - Figure 1: State-based conflicts by region}
\end{figure}
\begin{figure}[htbp]
\centering
\includegraphics[width=0.88\linewidth]{figures/fig_510.jpg}
\caption{Figure 2 - Figure 2: Global defense spending}
\end{figure}
{\small\color{refgray}\textbf{References:} [R010] Citi：黄金股回调是表象，真正的机会藏在现金流与远期定价的错位中; [R042] 摩根斯坦利：全球矿业正在经历一场“主权供需”的结构性重塑; [R029] JPM：美国军工复合体的真正危机不在于预算，而在于工业基础的结构性过时}

\section{结语}
综合来看，市场正从单一的通胀叙事转向多维的结构性重定价：AI资本开支的融资节奏与折旧冲击、能源供给的硬约束、中国出口角色的转换、以及新兴市场融资货币的单一化，共同构成了当前宏观策略的核心分歧。理解这些结构性变化，比追逐短期数据波动更为关键。
\newpage
\section{报告覆盖清单}
\begin{itemize}
  \item [R001] 宏观与利率：通胀但不加息，市场定价逻辑的重构 - 摩根斯坦利：市场真正低估的不是通胀本身，而是“通胀但不加息”的定价逻辑
  \item [R002] 未被主题摘要引用 - JPM：大宗商品持仓的收缩不是趋势反转，而是结构性再平衡的开始
  \item [R003] 未被主题摘要引用 - 摩根斯坦利：市场低估的不是美联储降息时点，而是中国资本管制收紧的资产重定价含义
  \item [R004] 宏观与利率：通胀但不加息，市场定价逻辑的重构 - NOM：全球市场正在为一场“滞胀式供给冲击”重新定价，但真正的分歧在于谁有资格加息
  \item [R005] 未被主题摘要引用 - 摩根斯坦利：铜关税决策将成为2026年下半年铜价最重要的单一路径依赖
  \item [R006] AI产业链：从GPU短缺到存储与计算架构再定价，企业级需求加速爆发 - 美国银行：AI基础设施需求正从“GPU短缺”转向“存储与计算架构的再定价”
  \item [R007] 宏观与利率：通胀但不加息，市场定价逻辑的重构 - GS：中国通胀数据的真正信号不在 CPI 本身，而在 PPI 的结构性集中
  \item [R008] 新能源与汽车：中国车企欧洲优势不在价格，储能市场供给端结构性收紧 - Bernstein：储能市场真正的超预期不在需求，而在供给端正在发生的结构性收紧
  \item [R009] 宏观与利率：通胀但不加息，市场定价逻辑的重构 - JPM：关税政策的真正转折不在税率高低，而在法律框架的切换
  \item [R010] 黄金与矿业：黄金股回调是表象，矿业主权供需重塑竞争格局 - Citi：黄金股回调是表象，真正的机会藏在现金流与远期定价的错位中
  \item [R011] 区域市场：新兴市场融资货币单一化，亚洲资金撤退信号值得警惕 - NOM：全球资金正在重新定价“美国例外论”，但亚洲的撤退信号更值得警惕
  \item [R012] 区域市场：新兴市场融资货币单一化，亚洲资金撤退信号值得警惕 - JPM：中国新ODI框架对全球银行的冲击比预期更广，财富管理业务的增量摩擦才刚刚开始
  \item [R013] 宏观与利率：通胀但不加息，市场定价逻辑的重构 - GS：日本通胀的“二阶导”拐点比绝对水平更重要
  \item [R014] 能源与大宗：供给硬约束重塑定价权，中国出口成金属市场泄压阀 - JPM：霍尔木兹海峡的“例外通行”不会改变全球LNG市场的结构性缺口
  \item [R015] 区域市场：新兴市场融资货币单一化，亚洲资金撤退信号值得警惕 - 摩根斯坦利：新兴市场真正的风险不是基本面，而是融资货币的单一化
  \item [R016] AI/算力：资本开支超越互联网泡沫，折旧冲击与融资结构成新焦点 / AI竞争格局：从模型性能转向Agent生态与成本效率，中国软件更具韧性 - HSBC：AI产业真正被低估的不是消费者市场，而是企业级需求的加速爆发
  \item [R017] 新能源与汽车：中国车企欧洲优势不在价格，储能市场供给端结构性收紧 - JPM：中国车企在欧洲的真正优势不是价格，而是零售执行与生态系统的护城河
  \item [R018] 地产与消费：成交量回升但结构性分化，二手房成更可靠温度计 - JPM：香港房价已接近全年预测上限，但真正的不确定性不在价格本身
  \item [R019] 地产与消费：成交量回升但结构性分化，二手房成更可靠温度计 - 摩根斯坦利：市场低估的并非政策力度，而是成交量回升背后的结构性分化
  \item [R020] 特高压与电网设备：订单翻倍背后，竞争从订单量转向议价权 - GS：中国数据中心市场真正被低估的不是投资规模，而是电力需求的非线性跃迁
  \item [R021] 新能源与汽车：中国车企欧洲优势不在价格，储能市场供给端结构性收紧 - JEF：储能市场的真正拐点不是增长，而是增长的结构性分化
  \item [R022] AI/算力：资本开支超越互联网泡沫，折旧冲击与融资结构成新焦点 / AI竞争格局：从模型性能转向Agent生态与成本效率，中国软件更具韧性 - JEF：AI竞争已从模型性能转向Agent生态与成本效率的再定价
  \item [R023] AI/算力：资本开支超越互联网泡沫，折旧冲击与融资结构成新焦点 - 摩根斯坦利：AI资本开支已超越互联网泡沫，但市场尚未定价折旧的滞后冲击
  \item [R024] AI/算力：资本开支超越互联网泡沫，折旧冲击与融资结构成新焦点 - 摩根斯坦利：GB200/300机柜出货量环比下降7\%，但这轮调整的真正含义是供应链议价权的再分配
  \item [R025] 区域市场：新兴市场融资货币单一化，亚洲资金撤退信号值得警惕 - NOM：人民币中间价模型正在发出一个被市场忽视的稳定信号
  \item [R026] 宏观与利率：通胀但不加息，市场定价逻辑的重构 / 能源与大宗：供给硬约束重塑定价权，中国出口成金属市场泄压阀 - BARC：市场低估了供给侧的持续收紧，而非单纯押注地缘风险
  \item [R027] 地产与消费：成交量回升但结构性分化，二手房成更可靠温度计 - GS：澳大利亚消费者信心连续三个月低迷，但市场真正该关注的不是情绪本身
  \item [R028] 能源与大宗：供给硬约束重塑定价权，中国出口成金属市场泄压阀 - JPM：中国出口正在成为全球金属市场结构性短缺的“泄压阀”
  \item [R029] 风险偏好：CLO市场从beta转向alpha，军工复合体面临工业基础结构性过时 / 黄金与矿业：黄金股回调是表象，矿业主权供需重塑竞争格局 - JPM：美国军工复合体的真正危机不在于预算，而在于工业基础的结构性过时
  \item [R030] 能源与大宗：供给硬约束重塑定价权，中国出口成金属市场泄压阀 - 摩根斯坦利：中国尿素出口底价的重新设定，是全球化肥定价逻辑的一次结构性重置
  \item [R031] 能源与大宗：供给硬约束重塑定价权，中国出口成金属市场泄压阀 / 化工与材料：合成橡胶供给缺口重塑定价权，松下投资者日揭示AI电源机遇 - 摩根斯坦利：中国取消尿素出口底价，全球化肥市场的定价权正发生结构性转移
  \item [R032] AI产业链：从GPU短缺到存储与计算架构再定价，企业级需求加速爆发 - BARC：企业级服务器的意外走强，正在改写服务器市场的增长叙事
  \item [R033] 新能源与汽车：中国车企欧洲优势不在价格，储能市场供给端结构性收紧 - Citi：市场低估的不是新能源渗透率，而是出口定价权的结构性重构
  \item [R034] AI产业链：从GPU短缺到存储与计算架构再定价，企业级需求加速爆发 - JPM：市场误读了SOCAMM降配，真正的机会藏在供给硬约束和NVIDIA的生态锁定里
  \item [R035] AI竞争格局：从模型性能转向Agent生态与成本效率，中国软件更具韧性 - JEF：市场对AI颠覆软件的恐惧过度，中国软件公司反而更具韧性
  \item [R036] AI竞争格局：从模型性能转向Agent生态与成本效率，中国软件更具韧性 - 摩根斯坦利：市场正在用不同估值逻辑为互联网公司定价，但真正的分歧不在AI硬件与软件之间
  \item [R037] 未被主题摘要引用 - NOM：2万亿投资计划不仅是短期催化剂，更是中国AI供应链国产化进程的加速器
  \item [R038] AI产业链：从GPU短缺到存储与计算架构再定价，企业级需求加速爆发 - Bernstein：存储行业的结构性重定价远未完成，真正的机会在NAND和HDD的后半场
  \item [R039] 特高压与电网设备：订单翻倍背后，竞争从订单量转向议价权 - Citi：特高压订单翻倍背后的关键信号——电网设备竞争已从“订单量”转向“议价权”
  \item [R040] AI/算力：资本开支超越互联网泡沫，折旧冲击与融资结构成新焦点 - 0005-41-BARC-U.S. Equity-Linked Strategies AI Capex Funding- Why Equity- Why Converts--260609
  \item [R041] 风险偏好：CLO市场从beta转向alpha，军工复合体面临工业基础结构性过时 - 摩根斯坦利：CLO市场真正的分化不在区域，而在管理人的议价权
  \item [R042] 能源与大宗：供给硬约束重塑定价权，中国出口成金属市场泄压阀 / 黄金与矿业：黄金股回调是表象，矿业主权供需重塑竞争格局 - 摩根斯坦利：全球矿业正在经历一场“主权供需”的结构性重塑
  \item [R043] 化工与材料：合成橡胶供给缺口重塑定价权，松下投资者日揭示AI电源机遇 - JPM：合成橡胶的供给缺口正在重塑化工定价权，而市场仍在关注需求侧
  \item [R044] AI/算力：资本开支超越互联网泡沫，折旧冲击与融资结构成新焦点 - 摩根斯坦利：AI融资的“夏季攻势”已提前到来，市场定价权正从基本面转向供给节奏
  \item [R045] 化工与材料：合成橡胶供给缺口重塑定价权，松下投资者日揭示AI电源机遇 - JPM：AI基础设施的下一轮投资机会，藏在松下投资者日被低估的细节里
  \item [R046] 未被主题摘要引用 - JPM：市场误读了英伟达进入基站射频单元的影响
  \item [R047] 未被主题摘要引用 - 摩根斯坦利：移动广告技术市场最被低估的增长引擎，不是需求，而是转化效率的杠杆效应
  \item [R048] AI产业链：从GPU短缺到存储与计算架构再定价，企业级需求加速爆发 - 摩根斯坦利：市场低估的不是半导体设备周期性，而是结构性瓶颈的再定价
  \item [R049] 风险偏好：CLO市场从beta转向alpha，军工复合体面临工业基础结构性过时 - 摩根斯坦利：欧洲猎头行业的真正风险不在需求，而在永久招聘的结构性萎缩
\end{itemize}
\section{逐篇报告摘录}
\subsection{[R001] 摩根斯坦利：市场真正低估的不是通胀本身，而是“通胀但不加息”的定价逻辑}
{\small\color{refgray}归类：宏观与利率：通胀但不加息，市场定价逻辑的重构}

摩根斯坦利：市场真正低估的不是通胀本身，而是“通胀但不加息”的定价逻辑

这份报告来自摩根斯坦利2026年6月8日的全球宏观论坛，由跨资产策略主管Serena Tang、美国经济学家Diego Anzoategui、全球宏观策略主管Matthew Hornbach以及首席美国股票策略师Michael Wilson联合撰写。它讨论的不是单一资产类别，而是试图回答一个根本问题：当劳动力市场强劲、通胀回升，但美联储保持中立时，资产价格应该如何重新定价？

报告的核心判断清晰且反直觉：当前市场对通胀的担忧被过度聚焦于负面情景，而忽略了通胀在“美联储不收紧”的条件下，对企业盈利和股票估值可能产生的正面推动。这不是一份看空通胀的报告，而是一份重新定义通胀含义的报告。

报告首先给出了一个看似矛盾的画面。非农就业新增172k，前值上修93k，三个月移动平均达到188k。失业率稳定在4.3\%。从表面看，这是一个典型的“过热”数据。但摩根斯坦利的经济团队提醒读者关注细节。

第一，就业增长的广度并不均匀。报告没有直接给出行业细分，但从其对“能源拖累”尚未显现的判断可以推断，当前的增长更多来自服务业的韧性，而非制造业或能源驱动的扩张。这意味着就业增长的可持续性可能被高估。

第二，薪资增长正在冷却。报告提到“薪资增长低于去年同期”，这被报告视为关键信号。当薪资增长放缓，而就业人数增长稳定时，单位劳动力成本的压力就会减轻。这是一个典型的“软着陆”信号，而不是“再通胀”信号。

第三，报告的隐含判断是，就业市场的韧性恰恰给了美联储“不行动”的空间。如果就业数据疲软，市场会预期降息，但也会引发衰退担忧。如果就业数据过热，市场会预期加息，压制估值。而当前的数据恰好落在了“足够强劲以避免衰退，但不够热以引发紧缩”的区间。

对读者的含义：劳动力市场不再是决定利率方向的单一变量。真正重要的是就业与薪资的组合。如果未来几个月薪资增速继续放缓，美联储的中立立场将更加稳固。这是当前市场定价中尚未充分反映的。

报告对CPI的拆解是整份文件中最具技术含量的部分。5月核心CPI环比仅0.22\%，同比2.8\%。但核心通胀的构成正在发生结构性变化。

最大的贡献来自租金，环比贡献0.11个百分点。而核心服务（除租金）贡献0.11个百分点，核心商品贡献0.02个百分点。对比4月的数据：租金贡献0.22，核心服务贡献0.17，核心商品贡献0.38。5月核心通胀的降温，主要来自租金和商品通胀的同步回落。

报告特别指出“关税影响继续消退”。4月核心商品通胀的跳升（0.39个百分点）被报告视为一次性冲击，而非趋势。5月的数据证实了这一点。

但报告也诚实地点出了风险：核心通胀的降温主要依赖服务业的放缓，而服务业通胀的粘性远高于商品。如果服务业薪资重新加速，核心通胀可能会再次抬头。

对读者的含义：市场对通胀的叙事过于二元——要么通胀失控，要么通缩。而报告展示的画面是：核心通胀正在温和降温，但过程不会平坦。关键在于，只要核心通胀不加速，美联储就没有理由改变其中立立场。这意味着债券收益率的上行空间可能比市场预期的要有限。

这是报告最核心的洞察，也是市场讨论最不充分的部分。报告引用了一张历史图表，显示PPI同比增速领先标普500销售增速约4个月。当通胀温和上升时，企业的定价权增强，名义收入增长加速。

更重要的是，报告提供了一个统计框架：在“盈利增长高于中位数”且“美联储政策中立”的历史时期，标普500指数未来12个月的平均回报率为16.3\%，中位数14.0\%，正回报概率高达92\%。

这个框架的逻辑链条是：通胀上升意味着名义GDP增长加速，企业收入增长加快；而美联储不收紧意味着估值不受压制。两者的结合创造了盈利增长与估值扩张的“双引擎”。

报告同时指出，上周的调整

[... middle omitted ...]

利率的相关性会从正相关切换为负相关。这意味着，一旦收益率突破这个水平，利率上行就会对股票估值形成直接的压制。

报告展示的相关系数图清晰地显示了这种非线性关系。在2025年的大部分时间里，股票与利率的正相关程度在0.4至0.6之间，这意味着利率上行伴随着股票上涨（因为增长预期改善）。但进入2026年，尤其是在5月，相关性已经转为负值，达到-0.8。这标志着市场已经进入了“利率是风险”的区间。

更值得关注的是债券波动率。报告将其描述为“是否导致更大幅度回撤的触发因素”。当利率波动率（MOVE指数）上升时，股票的远期市盈率会系统性收缩。报告展示的图表显示，MOVE指数与标普500远期市盈率之间存在显著的负相关关系。

对读者的含义：当前市场面临的核心矛盾是：通胀温和、美联储中立、盈利增长强劲，这些都对股票有利。但利率水平已经触及了4.5\%这个敏感阈值，利率波动率也在上升。如果利率继续上行，或者波动率进一步放大，股票的回调可能会从“技术性修正”演变为“基本面驱动的调整”。

\subsection{[R002] JPM：大宗商品持仓的收缩不是趋势反转，而是结构性再平衡的开始}
{\small\color{refgray}归类：未被主题摘要引用}

JPM：大宗商品持仓的收缩不是趋势反转，而是结构性再平衡的开始

这份JPM截至6月5日的最新商品市场持仓与资金流报告，给出了一个值得产业决策者和投资者认真对待的信号：全球商品期货市场的未平仓合约总值已连续三周下降，累计减少约800亿美元，至1.8万亿美元。表面看，这是市场情绪转冷的证据——农产品和贵金属价格大幅回调，持仓同步萎缩。但深入拆解资金流和仓位结构后，结论远比“避险情绪上升”复杂。

真正重要的不是持仓总量的收缩，而是收缩的内部结构。资金并未逃离商品资产类别，而是在不同板块之间剧烈再分配：贵金属和基本金属获得了净资金流入，而能源和农产品则遭遇了净流出。这种分化背后，是市场正在重新定价三个核心变量——美联储政策路径、中东供应风险的真实持续性、以及中国需求的结构性变化。

这份报告的价值不在于告诉你市场在跌，而在于告诉你，下跌中谁在买入、谁在卖出、以及这些行为背后的一致性逻辑。以下是我们从这份研报中提炼出的五个核心洞察。

1. 贵金属的持仓收缩是价格下跌的结果，而非资金逃离的原因，但空头信号正在累积

贵金属板块未平仓合约总值单周下降72亿美元至2570亿美元，连续多周下滑。但同期净合约流入达到54亿美元，其中黄金和白银分别吸收了24亿和27亿美元。这意味着持仓价值的下降几乎完全由价格下跌驱动，而非投资者主动平仓。

JPM的数据显示，截至6月2日，COMEX黄金期货的投机性净多头头寸仍高达11.2万手，较前一周增加1.47万手。这说明市场在价格下跌中反而在加仓。然而，报告同时指出，自中东冲突爆发以来，金价已下跌18\%，且短期价格动量信号已触及极端负值阈值，提示趋势可能衰竭。

这里存在一个关键张力：如果美联储因强劲就业数据而进一步转鹰，金价的下行风险将持续存在。当前持仓结构中，多头拥挤程度仍然较高，一旦宏观预期出现方向性变化，这部分头寸的平仓压力可能被低估。对于持有黄金资产的配置者而言，当前的核心问题不是“黄金是否还有避险价值”，而是“当前价格是否已经充分定价了加息路径的不确定性”。报告没有给出明确答案，但数据暗示，市场可能仍处于“定价调整”而非“趋势逆转”的阶段。

2. 基本金属的资金流入背后是关税政策的二阶效应，而非单纯的需求故事

基本金属板块净合约流入20亿美元，主要集中在镍（12亿美元）和铜（9亿美元）。表面看，这是对中国经济复苏预期的反映。但JPM的分析指向了一个更具体的驱动因素：美国铜关税审查的临近。

报告明确指出，6月30日美国商务部长将向总统提交关于是否需要进一步修改232条款铜关税的评估。JPM的观点是，特朗普政府更可能实施去年7月行政令中设定的阶梯式关税——2027年15\%，2028年升至30\%。这一政策框架的目的不是限制进口，而是确保已经流入美国的铜留在境内，并作为关键战略储备。

这会产生一个重要的二阶效应：在关税正式实施前的几个季度，美国将继续出现强劲的进口拉动，从而持续挤压非美地区的阴极铜供应，尤其是在中国进口套利窗口同时打开的时候。这意味着，当前铜价的支撑力量并非来自全球需求的同步复苏，而是来自政策预期驱动的区域套利行为。一旦关税政策落地或预期发生变化，这部分价格溢价可能迅速消退。

对于产业决策者而言，这提示了一个关键判断：铜的全球库存分布正在发生结构性变化——美国库存增加的同时，非美地区的可贸易库存正在收紧。这不是一个短期的供需错配，而是政策驱动的库存重配置。

能源板块未平仓合约总值基本持平于8190亿美元，但内部结构出现了值得注意的分化：原油市场净合约流出100亿美元，而天然气市场则录得23亿美元的净流入。

原油的流出集中在布伦特、迪拜和TTF天然气基准，这反映了投资者对中东供应风险持续性的怀疑。JPM的基本假设是霍尔木兹海峡将在6月重新开放，维持布伦特原

[... middle omitted ...]

移动，但目前海峡内仅有9艘液化天然气船，出口能力仍受严重限制。这意味着天然气市场的供应风险溢价尚未完全消退，而当前价格和持仓数据可能低估了这一尾部风险。

对于能源交易者和下游用户来说，报告的隐含判断是清晰的：市场正在定价“正常化”，但现实数据并不支持这一假设。原油和天然气的持仓分化，本质上是市场在“相信供应恢复”和“看到供应受限”之间的犹豫。这种犹豫本身就是一种风险。

4. 农产品板块的持仓崩溃是多重因素共振的结果，但软商品提供了结构性机会

农产品板块是本周持仓收缩最严重的领域，未平仓合约总值下降3\%至3820亿美元，净持仓价值单周减少约60亿美元。谷物和油籽是主要拖累，玉米、大豆和小麦的净多头头寸均出现显著下降。

但报告揭示了一个重要的结构性分化：尽管谷物和油籽遭遇了16亿美元的净流出，软商品板块却获得了48亿美元的净流入。可可、咖啡、糖等品种的资金流入，部分抵消了农产品的整体下滑。这一分化提示，农产品市场的下跌并非系统性恐慌，而是板块内部基本面差异的体现。

\subsection{[R003] 摩根斯坦利：市场低估的不是美联储降息时点，而是中国资本管制收紧的资产重定价含义}
{\small\color{refgray}归类：未被主题摘要引用}

摩根斯坦利：市场低估的不是美联储降息时点，而是中国资本管制收紧的资产重定价含义

这份来自摩根斯坦利亚洲团队的研报，表面上讨论的是美联储政策立场对中国的溢出效应，但真正值得产业决策者和资产配置者关注的，不是美联储何时降息，而是中国正在主动收紧资本外流渠道，并且这一动作正在改变人民币资产的定价逻辑。

报告的时间节点值得注意——2026年6月。此时市场主流叙事仍聚焦于美联储降息路径的不确定性，但摩根斯坦利把分析重心放在了另一个维度：当美国经济从消费驱动转向资本支出驱动，当中国在资本账户管理上从“被动应对”转向“主动调控”，两个经济体之间的资金流向和资产定价关系正在发生结构性的重置。

这份报告的核心价值在于，它提供了一套完整的分析框架，把美联储政策、中国经济内循环、人民币汇率和资本流动整合到一个逻辑链条中。以下是我们提炼出的五个关键洞察。

1. 美国经济正在经历增长引擎的结构性切换，消费不再是唯一主角

摩根斯坦利对 年美国实际GDP增长的预测，表面上是一个温和加速的数字——从2025年的1.8\%逐步升至2027年的2.5\%——但真正重要的不是总量，而是结构。

报告明确指出，增长贡献中最大的变化来自企业固定投资。2025年企业固定投资对GDP的贡献为0.8个百分点，2026年预计升至0.9，2027年进一步升至1.0。与此同时，消费的贡献虽然也在上升，但幅度明显更小——从2025年的1.0个百分点升至2027年的1.3。

这意味着什么？美国经济正在从“消费者花钱”驱动，转向“企业花钱”驱动。而企业投资的核心驱动力，报告毫不含糊地指向了AI相关支出。根据报告的估算，2025年和2026年，AI相关投资对非住宅固定投资的贡献均高达1.6个百分点，而非AI部分则是负贡献——分别为-0.5和-0.4个百分点。

这是一个典型的“一枝独秀”格局。AI投资不仅支撑了整体资本支出，还在对冲传统领域的投资下滑。对于中国资产而言，这一结构性变化的含义是双重的：一方面，美国经济韧性超预期的来源从消费转向投资，这意味着利率敏感度可能发生变化；另一方面，AI投资的持续高增长意味着全球科技产业链的重组仍在加速，中国在这一链条中的位置需要重新审视。

2. 美联储降息的条件比市场想象的更苛刻，能源价格是关键变量

市场对美联储降息的预期已经反复调整，但摩根斯坦利提供的新信息在于，它明确指出了两个制约降息的结构性因素。

第一个是能源价格对消费的挤压。报告做了一个非常具体的测算：如果今年汽油价格上涨15\%，每个家庭平均增加的汽油支出约为375美元，而同期退税增加的平均幅度只有325美元。这意味着能源价格上涨将完全抵消财政刺激带来的消费增量。这还不是最坏的情况——如果汽油价格上涨20\%，家庭支出将增加500美元，比退税多出175美元。

第二个是劳动力市场的韧性。报告数据显示，非农就业人数虽然从2023年的月均45万下降，但目前的水平仍然高于摩根斯坦利估算的“盈亏平衡点”——即维持失业率稳定所需的月均就业增量，目前约为5万人。只要就业数据持续高于这个水平，美联储就没有动力因为“担忧就业下行风险”而提前降息。

从这些数据可以推导出一个判断：美联储降息的时间窗口，不是由通胀回落的速度单独决定的，而是由“能源价格不超预期上涨”和“就业数据不超预期恶化”两个条件同时满足来决定的。而这两个条件在当前的地缘政治环境下，都面临相当大的不确定性。

3. 中国收紧资本管制的真实意图不是应对贬值压力，而是主动塑造资金流向

这可能是这份报告最容易被市场忽视但最具前瞻性的部分。摩根斯坦利明确指出，中国近期收紧对外投资管控，并非因为人民币贬值压力加剧——恰恰相反，人民币贬值预期实际上已经有所缓解。

报告提供了两个关键数据来支持这一判断。第一，出口商的结汇比

[... middle omitted ...]

的回报率在2025年约为9.4\%，而直接投资的回报率仅为4.4\%。收紧资本管制意味着国内投资者将失去一部分海外高收益机会，这可能会改变国内资产的风险定价结构——更多的资金追逐有限的国内优质资产，可能推动部分资产类别出现估值溢价。

4. 人民币汇率定价权正在从外部转向内部，USDCNY的波动逻辑变了

摩根斯坦利对人民币汇率的预测框架值得仔细拆解。报告认为，CFETS人民币篮子指数将出现温和升值，但USDCNY的走势本质上取决于美元自身的强弱。这一判断本身并不新鲜，但报告提供的历史数据揭示了一个重要的结构性变化。

回顾过去几轮美元周期中人民币的反应：2018年美元走强时，CFETS指数下跌4.5\%，USDCNY贬值9.5\%，人民币的贬值幅度几乎是美元升值的两倍；而在 年美元走弱时，CFETS指数基本持平，但USDCNY升值9.8\%，人民币的升值幅度几乎与美元贬值幅度相当。这表明，在过去，人民币汇率主要由外部因素驱动，中国央行的干预更多是“平滑波动”而非“设定方向”。

\subsection{[R004] NOM：全球市场正在为一场“滞胀式供给冲击”重新定价，但真正的分歧在于谁有资格加息}
{\small\color{refgray}归类：宏观与利率：通胀但不加息，市场定价逻辑的重构}

NOM：全球市场正在为一场“滞胀式供给冲击”重新定价，但真正的分歧在于谁有资格加息

2026年6月，NOM在新加坡举办了年度旗舰亚洲投资论坛。在为期四天的会议中，全球首席经济学家团队给出了一个高度一致的判断：当前全球经济正处于一场“供给冲击复合体”之中，其结构性影响远超市场定价所反映的程度。这份研报的核心价值不在于预测了某个具体利率路径，而在于它提供了一套区分“谁在真正面临通胀压力”与“谁只是在承受增长放缓”的分析框架。

市场目前最大的误判，是认为所有经济体的通胀压力是同质的。NOM经济学家们明确指出：当前的通胀并非需求过热驱动，而是由能源、芯片和食品三重供给冲击叠加而成。这意味着，传统“通胀高则加息”的简单逻辑已经失效。真正值得关注的问题是：哪些经济体有资格、有能力、有必要加息？哪些经济体应当忍受通胀、优先保增长？

1. 供给冲击的广度被严重低估，全球供应链压力指数已回到疫情外历史最高

NOM全球首席经济学家Rob Subbaraman在论坛上指出，即使霍尔木兹海峡很快重新开放，全球供给侧的中断惯性仍将显著存在。纽约联储的全球供应链压力指数（GSCPI）目前已经攀升至1998年以来的最高水平——如果不算疫情期间的极端值。历史经验表明，2021年12月GSCPI见顶后，美国CPI通胀在六个月后才达到峰值。

这意味着什么？当前这轮供给冲击的传导尚未完全反映在终端价格中。市场对通胀“见顶回落”的乐观预期，可能忽视了供应链压力向消费端传导的滞后性。更关键的是，NOM强调这并非单一的油价冲击，而是一个“能源-芯片-食品”三重价格冲击的组合。芯片价格飙升正在通过AI基础设施投资传导至软件和服务价格，而食品价格冲击则与地缘政治和气候因素叠加。

这一判断对资产定价的含义是：债券市场可能低估了通胀持续性的尾部风险。NOM展示了从2021年12月到2026年5月的四个时间点的主权债券收益率曲线快照，清晰地呈现出收益率上升和曲线陡峭化的趋势。该机构认为，这反映了三个结构性变化：通胀上升与央行独立性侵蚀、持续的大规模财政赤字与公共债务比率上升、以及巨型科技公司也开始发债时面临的巨量政府债券供给。

2. 美国拥有“加息特权”，但美联储实际可能按兵不动——AI繁荣创造了一种新的通胀粘性

美国首席经济学家David Seif的论述揭示了美国经济的特殊性。美国正以高于趋势的速度增长，且增速快于其他发达经济体。AI基础设施投资正将美国GDP增速拉高整整一个百分点。美国在AI投资中占据了“狮子的份额”，而伊朗战争对美国的负面影响小于大多数发达经济体，因为北美拥有独立的天然气市场且美国是石油净出口国。

但通胀问题并不简单。NOM指出，美国核心通胀上升的一个新来源是“软件商品”，这与AI资本支出繁荣直接相关。更值得注意的是，美联储主席Warsh偏好的“截尾均值PCE”指标在历史上对通胀趋势变化的反应滞后于核心PCE。而Warsh本人比大多数FOMC成员更偏鸽派，这可能限制他对货币政策的掌控力。

NOM的最终判断是：美联储至少到2027年底都将维持利率不变，加息和降息的概率大致相等。这个看似平淡的结论背后，是一个重要的结构性叙事——美国经济已经形成了“AI繁荣支撑增长、供给冲击维持通胀、但政策制定者不愿收紧”的微妙平衡。对于投资者而言，这意味着美国利率的“长期高位”可能比市场预期的更持久，而AI相关资产的定价逻辑需要纳入更高的利率贴现因子。

3. 欧元区的“滞胀”比2022年更危险，但央行只能选择“有限收紧”

NOM欧洲首席经济学家George Buckley的论述可能是这份研报中最具警示性的部分。自伊朗战争爆发以来，欧元区是共识预期被“滞胀式下修”幅度最大的发达经济体——增长预期更弱，通胀预期更强。但关键区别在于，当前的政策

[... middle omitted ...]

晰信号：市场参与者已经接受了“滞胀优先保增长”的思维框架。对于欧元区资产而言，这意味着利率峰值将低于市场此前定价，但增长前景的持续疲软将压制风险资产的表现。

4. 日本是唯一一个真正需要“正常化加息”的发达经济体——但市场仍认为它会落后于曲线

NOM日本首席经济学家Kyohei Morita的论述在整份报告中显得尤为独特。他假设中东紧张局势不会持续到2026年下半年、霍尔木兹海峡逐步重新开放，在此前提下，日本经济在二季度收缩后将重回复苏轨道。中东局势引发的上游通胀将以约六个月的滞后传导至CPI，核心CPI预计在2027年一季度达到3.6\%的峰值。

但真正支撑日本加息逻辑的是劳动力市场结构变化。2026年春季工资谈判（春斗）正在确认工资上涨变得更加协调和可持续。加上企业通过节省劳动力的投资来降低劳动依赖度，工资上涨的可持续性进一步增强。NOM预计，在日本央行将政策利率上调至1.5\%（中性利率估计区间内）的过程中，2026年6月和12月、2027年6月将各加息一次。

\subsection{[R005] 摩根斯坦利：铜关税决策将成为2026年下半年铜价最重要的单一路径依赖}
{\small\color{refgray}归类：未被主题摘要引用}

摩根斯坦利：铜关税决策将成为2026年下半年铜价最重要的单一路径依赖

市场参与者正在为一个关键的政策节点做准备。美国对精炼铜进口是否征收15\%关税的最终决定，预计将在2026年下半年落地。这并非一个普通的贸易政策事件。它直接关系到过去半年多时间里，全球约2.6\%的铜需求被美国囤货行为所吸收的异常状态能否持续。

摩根斯坦利在最新发布的金属市场报告中，系统拆解了这一决策的三种可能情景及其对铜价的含义。报告的核心判断是：市场当前定价隐含了约43\%的概率认为15\%关税将在2027年1月前生效。但真正重要的不是这个概率本身，而是无论最终结果如何，铜市场的定价逻辑都将经历一次结构性重置。当前COMEX相对LME约6\%的溢价，以及美国已超过一年正常进口量的库存水平，都预示着市场已经提前在为这个决策定价。接下来的问题不是铜价会不会波动，而是波动的方向、幅度和持续时间将如何重塑全球铜的贸易流和库存分布。

摩根斯坦利分析认为，如果美国在2026年7月左右提前宣布从2027年1月起实施15\%的关税，这将是最为看涨的情景。逻辑链条清晰：提前宣布意味着市场有充足的时间窗口进行套利操作，将铜加速运往美国。这会进一步收紧美国以外市场的铜供应，推高LME和COMEX两个基准价格。具体而言，LME的远期曲线前端会进入现货溢价状态，而COMEX相对于LME的溢价将向15\%靠拢。

这里需要特别留意报告中的一个隐含判断：宣布时间越早，看涨效应越强。如果关税宣布时间临近实施日期，或者实施日期被提前，市场就没有足够时间进行物流操作，看涨效应将大打折扣。这意味着，政策制定者的“通知期”本身就是价格变量的组成部分。投资者不能简单地将“关税=看涨”作为交易依据，必须同时评估关税宣布与实施之间的时间差。这个时间差决定了套利空间的实际可操作性。

最看跌的情景是关税被完全排除。摩根斯坦利指出，一旦如此，美国当前每年约260千吨（占全球需求2.6\%）的超额进口将大概率停止。这将直接对LME和COMEX铜价形成下行压力。COMEX相对LME的溢价将迅速收窄至平价甚至小幅贴水，LME远期曲线也会趋于宽松。

然而，报告同时提供了一个重要的结构性缓冲判断：铜价在每吨12000美元附近存在坚实支撑。这个支撑并非来自需求端的刺激，而是来自供给端的长期结构性约束——持续的供应短缺和预期的市场缺口。换言之，即便关税带来的囤货需求消失，铜市场的基本面并不支持价格无底线下行。这里需要读者注意的是，摩根斯坦利给出的12000美元支撑位是基于当时（2026年4月）的预测框架。实际价格运行中，支撑位会随着基本面数据变化而动态调整，但供给约束作为长期定价锚点的逻辑不会改变。

第三种情景是决策被推迟到未来某个日期。摩根斯坦利将其定性为“中性偏温和负面”。表面上看，这延续了现状，关税选项仍然在桌面上。但报告敏锐地指出，如果市场参与者将推迟解读为“关税可能性下降”的信号，铜价可能面临温和下行压力。

这个判断的微妙之处在于，它揭示了市场对政策不确定性的定价方式。当决策被推迟时，不确定性本身并未消失，但市场会倾向于重新评估概率。如果推迟被视为一种隐性的否决，那么当前已经定价的部分关税溢价就可能被压缩。这意味着，即使没有明确的利好或利空消息，仅仅因为时间线的拉长，市场预期也会发生偏移。对于投资者而言，需要关注的不仅是决策本身，还有市场对决策时间线变化的反应函数。

4. 报告尚未完全回答的关键问题：关税实施后，囤货行为会如何演变？

摩根斯坦利对2027年关税实施后的情景进行了推演，但留下了一个值得深入探讨的问题。报告指出，一旦15\%关税生效，COMEX-LME价差将反映这一关税水平，意味着COMEX仍有相对LME的上涨空间。但与此同时，美国超额进口的动力将消失，因为一旦需要实际支付关税，套

[... middle omitted ...]

战略缓冲。投资者需要密切跟踪美国终端铜消费数据以及COMEX库存的变动趋势，才能对这个未解问题形成自己的判断。

5. 宏观因子正在成为铜价的双刃剑，COMEX的多头仓位已达历史极值

报告在最后部分提醒了一个容易被关税叙事掩盖的重要变量：宏观因素。摩根斯坦利观察到，铜价对美联储加息预期的变化反应强烈。当市场定价显示2026年可能加息时，铜价出现负面反应。COMEX铜在风险偏好日表现优于LME，但在风险规避日表现更差。同时，COMEX的净多头持仓已经达到历史最高水平。

这个观察的So what含义是：铜价当前处于一个多重因素叠加的脆弱均衡状态。关税预期提供了看涨支撑，但宏观紧缩预期构成上行阻力。更关键的是，多头持仓的极端集中意味着，一旦关税预期出现任何负面变化，或者宏观数据意外转鹰，拥挤的多头交易可能迅速平仓，放大价格下行幅度。投资者在关注关税决策的同时，必须将COMEX持仓数据和美国宏观数据作为同等重要的观察变量。关税看涨逻辑的有效性，部分取决于宏观环境是否配合。

\subsection{[R006] 美国银行：AI基础设施需求正从“GPU短缺”转向“存储与计算架构的再定价”}
{\small\color{refgray}归类：AI产业链：从GPU短缺到存储与计算架构再定价，企业级需求加速爆发}

美国银行：AI基础设施需求正从“GPU短缺”转向“存储与计算架构的再定价”

市场正在经历一个关键的认知转变。过去两年，投资者对AI基础设施的关注几乎全部集中于GPU的供应与需求，以及围绕英伟达的资本开支叙事。但美国银行在2026年全球科技大会上释放出的信号表明，真正结构性的变化正在更底层发生：AI的推理和智能体工作负载正在重塑整个IT硬件价值链，从存储到传统服务器，从分销商到企业级软件，都在经历一轮需求与定价权的同步重构。

这份报告的核心判断不是“AI需求依然强劲”这样广为人知的结论，而是一个更微妙但更具投资含义的洞察：需求强劲已是共识，但市场低估的是供给侧的结构性受限如何改变了定价模式与竞争格局。 美国银行因此上调了IBM、NetApp、希捷和西部数据的目标价，其逻辑并非简单的需求拉动，而是这些公司正在从“卖硬件”转向“卖架构与锁定收益”。

以下是我们从这份研报中提炼出的五个关键洞察，它们共同指向一个结论：AI基础设施投资的下一阶段，赢家将是那些能够将规模转化为议价权、将技术转化为长期合同的企业。

1. 推理与智能体AI正在创造传统计算的新需求周期，这比市场预期的更持久

大多数投资者仍将AI服务器需求视为围绕GPU的单一叙事。但美国银行在会议中发现，推理和智能体AI正在成为传统计算的新驱动力。这一判断的关键在于，智能体工作负载将部分计算从并行GPU处理转移到了顺序CPU处理。这意味着，CPU密集型服务器和企业基础设施的需求正在被重新激活。

Dell和HPE都报告了传统计算需求的显著增长，而IBM则展示了推理和现代化工作负载如何在其非GPU平台（如大型机）上创造新的消费增长。这不仅仅是“AI服务器之外还有传统服务器”的简单补充，而是意味着整个企业IT基础设施的升级周期正在被AI的部署节奏所驱动。当企业开始大规模部署智能体时，它们需要的不仅是算力，更是数据存储、查询、保护与回馈的完整闭环。这为整个IT硬件生态创造了比单纯GPU采购更持久、更广泛的需求基础。

2. 存储正从“配角”变为AI架构的战略瓶颈，定价权正在转移

这是整份研报中最具颠覆性的判断之一。美国银行指出，存储正在成为AI基础设施中新的战略瓶颈。训练、推理和智能体工作流都在增加需要存储、保护、查询并反馈回模型的数据量。西部数据将HDD的EB增长预期为未来3-5年可能超过25\%，而希捷则明确指出需求持续超过供应，订单已覆盖未来4-5个季度，客户正在提前数年规划EB分配。

更值得关注的是定价模式的变化。报告特别提到，SanDisk展示了其新的商业模式，包括数量承诺以及合同初期固定定价、后期可变定价的组合。这标志着存储行业正在从“按需购买”转向“长期锁定”。对于投资者而言，这意味着存储公司的收入可见性和利润稳定性正在结构性提升，而不仅仅是周期性的改善。Dell ISG总裁Arthur Lewis也强调，智能体AI的数据增长意味着“冷数据”或“暗数据”不再存在，所有数据都变得具有战略价值。这进一步强化了存储作为AI基础设施核心资产的逻辑。

报告揭示了一个被市场忽视的趋势：组件环境的紧张正在重塑分销商和电子制造服务商的角色。内存、CPU、HDD、NAND都显示出定价或成本通胀的迹象。在这种环境下，分销商（如Arrow、Avnet）和EMS（如Flex、CLS、SANM）的价值不再仅仅是“搬运工”，而是成为了客户在复杂供应链中确保供应、管理长周期交货的关键合作伙伴。

Arrow的临时CEO Bill Austen明确指出，当前的复苏是单位驱动的，而非价格驱动的，这与疫情短缺周期或对双重下单的担忧有本质区别。这种区别意味着，分销商的利润率改善是更高质量、更可持续的。Avnet的CEO Phil Gallagher也强调了复杂性对Avnet的正面

[... middle omitted ...]

 Assistant for Z）并未削弱大型机的粘性，反而扩大了使用。使用该工具的客户MIPs利用率增长是基础客户群的2-3倍，而Linux、容器及相邻工作负载的专业MIPs增长速度是核心MIPs的3-4倍。这意味着，IBM正在将AI转化为一种锁定机制，而非迁移工具。同时，组件成本的上升正在创造对IBM更高效率系统的增量需求，因为其垂直整合的架构在定价控制和利润保护方面具有天然优势。这解释了美国银行为何将IBM的目标价从300美元上调至315美元，并采用了更高的估值倍数。

5. 报告尚未完全回答的关键问题：需求可持续性的“质量”与定价权的“持久性”

尽管美国银行的会议纪要给市场注入了极强的信心，但报告本身也留下了一些值得深思的未解问题。首先，需求的可持续性虽然被广泛讨论，但其“质量”仍有待验证。例如，Avnet的复苏中，内存定价从占业务的5-7\%跃升至10-15\%，主要是因为价格翻倍而非单位增长。这种由价格驱动的增长是否可持续，取决于终端需求能否消化更高的成本。

\subsection{[R007] GS：中国通胀数据的真正信号不在 CPI 本身，而在 PPI 的结构性集中}
{\small\color{refgray}归类：宏观与利率：通胀但不加息，市场定价逻辑的重构}

GS：中国通胀数据的真正信号不在 CPI 本身，而在 PPI 的结构性集中

五月的中国通胀数据，表面上看是一个“温和”的故事。CPI 同比持平于 1.2\%，PPI 从 2.8\% 升至 3.9\%。市场大多会将其解读为“需求端没有过热风险，供给端成本压力正在传导”。但GS这份最新报告揭示了一个更值得关注的判断：真正重要的不是通胀的绝对水平，而是 PPI 反弹的结构——超过 80\% 的贡献来自上游行业，而核心 CPI 却在走软。

这意味着什么？意味着当前的通胀格局不是需求驱动的全面复苏，而是一场由能源和化工价格主导的“成本冲击”，并且这个冲击正在挤压下游利润空间。对于投资者和产业决策者而言，忽视这个结构性差异，可能会对宏观环境做出误判。

这份报告提供了一个关键的分析框架：把通胀数据拆解为“上游成本推动”和“下游需求拉动”两个维度。五月的数据清晰表明，前者在加速，后者在减速。核心 CPI 从 1.2\% 回落至 1.1\%，虽然幅度不大，但方向明确——旅游相关服务价格的走软，暗示消费端的恢复可能比预期更脆弱。

1. 上游 PPI 的反弹高度集中，这意味着成本压力并非均匀分布

GS报告中最引人注目的数字是：上游行业贡献了 PPI 同比反弹的 82\%。具体来看，化学品、石油天然气和煤炭价格分别贡献了 0.3、0.2 和 0.2 个百分点，推动 PPI 同比从 4 月的 2.8\% 升至 5 月的 3.9\%。

这个集中度意味着什么？它意味着成本压力并非广泛扩散，而是集中在少数几个关键上游环节。对于下游制造业来说，这是一个典型的“成本挤压”场景——上游涨价，但下游消费品 PPI 同比仅为 -0.8\%，虽然比 4 月的 -1.0\% 略有改善，但仍在负值区间。

从产业链角度看，这种结构意味着：拥有上游资源或定价权的企业能够转嫁成本甚至受益，而处于中下游、竞争激烈的行业将面临利润率的双重压力——原材料成本上升，但产成品价格无法同步上涨。这一判断对于研究化工、钢铁、建材等行业的企业盈利至关重要。

五月的核心 CPI 同比从 1.2\% 降至 1.1\%，虽然变化幅度不大，但方向值得警惕。GS明确指出，这“可能反映了旅游相关服务价格的走软”。具体来看，交通服务通胀从 5.1\% 放缓至 4.7\%。

这个细节是理解当前消费复苏质量的关键。核心 CPI 剔除了食品和能源的波动，是衡量内生需求温度的更准确指标。它的走软，与 PPI 的加速形成鲜明对比——上游在“热”，下游在“凉”。

对于投资者而言，这意味着需要重新审视对消费复苏的预期。如果核心 CPI 继续走软，那么依赖消费拉动增长的故事可能需要调整。报告没有直接给出结论，但数据隐含的逻辑是：服务消费的恢复可能已经过了最强劲的阶段，后续动能需要新的催化剂。

五月的食品 CPI 同比为 -1.7\%，较 4 月的 -1.6\% 进一步走低。其中猪肉价格同比下跌 16.1\%，较 4 月的 -15.2\% 跌幅扩大。鲜果价格也从 -1.0\% 进一步降至 -2.2\%。

GS在报告中详细拆解了食品项的贡献结构。猪肉价格的持续低迷，反映的是供给端过剩的问题——生猪产能去化速度慢于预期。这个问题在 2023 年就已经显现，但持续到 2024 年五月仍未解决，说明市场出清的时间可能比预想更长。

对于农业和食品加工行业，这意味着：上游养殖企业的亏损周期可能延长，而下游食品加工企业则受益于较低的原料成本。但后者能否将成本优势转化为利润，取决于终端需求能否支撑产品定价。

4. 能源价格是唯一支撑非食品 CPI 的变量，但它的持续性存疑

非食品 CPI 同比从 1.8\% 微升至 1.9\%，但GS指出，这完全由能源价格上涨驱动。燃料成本同比涨幅从 17.4\% 扩大至 21.1\%，贡献了非食品 CPI 的全

[... middle omitted ...]

导。

5. 报告尚未完全回答的关键问题：PPI 向 CPI 的传导何时会真正启动

GS这份报告提供了详尽的数据拆解，但它留下了一个悬而未决的问题：上游 PPI 的上涨何时会传导到下游 CPI？从数据看，消费品 PPI 仍在负值区间（-0.8\%），说明传导尚未发生。

这是一个经典的宏观经济学问题——成本推动型通胀的传导需要条件：要么下游需求足够强劲，让企业能够涨价；要么上游涨价持续足够长，迫使企业不得不转嫁成本。目前来看，这两个条件都不完全具备。

报告没有给出预测，但数据暗示了两种可能的情景：一是如果需求端持续疲软，那么 PPI 的上涨只会压缩企业利润，不会引发 CPI 上行；二是如果上游涨价持续超预期，那么最终会通过成本压力倒逼下游涨价，但时间滞后可能长达 2-3 个季度。哪种情景更可能发生，取决于未来几个月的需求数据。

基于这份报告，一个更实用的分析框架是：不要只看 CPI 和 PPI 的绝对水平，而要关注它们的内部结构。具体来说，有三个维度值得持续追踪：

\subsection{[R008] Bernstein：储能市场真正的超预期不在需求，而在供给端正在发生的结构性收紧}
{\small\color{refgray}归类：新能源与汽车：中国车企欧洲优势不在价格，储能市场供给端结构性收紧}

Bernstein：储能市场真正的超预期不在需求，而在供给端正在发生的结构性收紧

储能行业正在经历一个罕见的时刻：需求端的高增长已经被市场充分定价，但供给端的变化——产能利用率从69\%跃升至94\%、头部企业份额的微妙再分配、以及海外市场的结构性启动——才是这份Bernstein最新全球储能追踪报告最值得关注的信号。

2026年4月，中国储能电池出货量达到93.1GWh，同比增长92\%，环比增长14\%。这个数字本身并不令人意外——过去一年中国储能市场一直保持着高速增长。真正值得深究的是隐藏在总量背后的三个结构性变化：产能利用率已经逼近极限、竞争格局正在从“百花齐放”转向“强者分化”、以及海外市场终于开始贡献有意义的增量。

Bernstein在这份报告中给出了一个清晰的判断：中国储能市场已经从“产能过剩”阶段进入“供需紧平衡”阶段，而这一点尚未被市场充分定价。对于关注电池产业链的投资者来说，这意味着定价逻辑需要从“量的增长”切换到“利润率的修复”。

1. 产能利用率从69\%到94\%的变化，意味着行业定价权正在发生转移

这是这份报告中最容易被忽视但最重要的数据点。2025年4月，中国储能电池工厂的产能利用率仅为69\%。一年后，这个数字跃升至94\%。在同一时期，月度产能从67.1GWh增长到98.0GWh，增幅为46\%——产能确实在扩张，但出货量的增长速度（92\%）远远超过了产能扩张的速度。

94\%的产能利用率意味着什么？在制造业中，当产能利用率超过85\%，行业就进入了“卖方市场”。当利用率超过90\%，任何意外的需求波动都会导致供应紧张和价格上行。Bernstein的报告明确指出，这“预示着近期市场条件正在收紧，因为需求增长正在赶上行业产能”。

这对投资框架的含义是直接的：过去两年，储能电池价格持续下跌的主要驱动力是产能过剩。一旦供需关系逆转，价格下跌的速度将显著放缓，甚至可能出现阶段性企稳。那些拥有成本优势和规模优势的头部企业，其单位利润将不再被压缩，反而可能开始修复。

报告没有直接展开的是：94\%的产能利用率是在什么价格水平下实现的？如果价格继续下行，是否会有产能主动退出？这些问题需要结合后续的月度数据来验证。

2026年前四个月，中国储能出货量达到309.1GWh，同比增长109\%。这个增速不仅超过了Bernstein年初对全球储能市场86\%的全年增长预期，更重要的是，其驱动力正在从政策驱动转向经济性驱动。

报告提到，中国需求强劲的背后是“电池成本下降和政策环境支持使得储能经济性更具吸引力”。这是一个重要的定性判断。当储能项目的内部收益率（IRR）不再依赖补贴，而是可以独立站住脚时，需求增长的可持续性就大大增强了。

招标数据进一步验证了这一判断。4月份中国储能招标电池容量达到55.6GWh，同比增长44\%。2026年前四个月累计招标量达到224.6GWh，同比增长190\%。招标量是出货量的领先指标，190\%的增速意味着未来几个月的出货量仍有强劲支撑。

但这里有一个需要警惕的差异：招标量的增速（190\%）远高于出货量的增速（109\%）。这通常意味着两种可能：一是项目从招标到实际出货的时间在拉长，二是存在一定程度的“重复招标”或“虚标”。Bernstein的报告没有对此展开分析，但这正是需要持续跟踪的关键变量。

CATL在4月份的储能电池出货量为19.8GWh，同比增长72\%，市场份额为21.3\%。2026年前四个月累计市场份额为22.5\%。这个数字看起来并不算特别高——毕竟CATL在动力电池领域的份额超过40\%。但仔细看竞争格局的变化，才能理解Bernstein为什么认为CATL是“参与这一主题的最佳方式”。

第二梯队的格局正在发生微妙的变化。EVE和Hithium分别以9.6\%和9.5

[... middle omitted ...]

市场地位仍然稳固，但并没有形成绝对的垄断优势。另一方面，第二梯队的内部正在发生分化：一些企业正在快速追赶，而另一些则开始掉队。这种分化格局对投资者来说既是机会也是风险——选对标的的回报可能远超行业平均水平，但选错的代价也会更大。

Bernstein对CATL A股给出“跑赢大盘”评级，目标价800元人民币，对H股给出“同步大市”评级。这个评级差异本身就值得思考：同一家公司，不同市场的定价为什么会出现如此大的分歧？

中国储能市场的绝对规模已经不需要过多强调，但海外市场的变化可能才是未来12-18个月最大的增量来源。

德国市场的表现令人瞩目。4月份储能电池安装量达到0.8GWh，同比增长74\%，其中大型储能贡献了0.3GWh，同比增长833\%。这个增速背后反映的是欧洲能源转型的加速，以及大型储能项目从规划到落地的周期正在缩短。2026年前四个月德国累计安装量2.8GWh，同比增长39\%，增速虽然不如中国，但考虑到德国市场的基数效应和政策的确定性，这个增速是可持续的。

\subsection{[R009] JPM：关税政策的真正转折不在税率高低，而在法律框架的切换}
{\small\color{refgray}归类：宏观与利率：通胀但不加息，市场定价逻辑的重构}

市场正在低估一个关键变化：美国关税政策的核心博弈点，已经从“加多少”转向了“用什么法律工具来加”。JPM在最新发布的研报中明确指出，随着IEEPA关税被法院推翻、Section 122关税即将于7月24日到期，美国政府正在系统性地将关税权力基础从临时性的紧急授权，转向更具持久性的Section 301框架。这个切换，表面上是法律技术的调整，实质上意味着关税政策的可预期性进一步降低、诉讼风险显著上升，而企业需要面对的，将是一轮比2018年中美贸易摩擦更广泛、更碎片化的政策不确定性。

这不是一个关于关税税率升降的故事。这是一个关于政策工具选择如何重塑全球供应链决策逻辑的故事。

1. 关税的“名义税率”和“实际税率”之间存在巨大鸿沟，而豁免清单才是真正的博弈场

JPM的测算揭示了一个被广泛忽视的事实：美国对全球征收的平均关税税率，在静态权重下约为11\%，但实际征收中观察到的税率仅有6.7\%。这4.3个百分点的差距，几乎全部来自豁免条款。能源、药品、电子产品、USMCA合规贸易——这些领域的大规模豁免，使得实际税率远低于政策声明的水平。

这意味着，企业在评估关税风险时，不能只看白宫发布的税率数字，而必须拆解每一类产品的豁免状态。当前Section 301的强制劳动调查中，美国贸易代表办公室提议对60个经济体征收10\%-12.5\%的关税，但豁免范围比之前的IEEPA和Section 122更为宽泛——新增了飞机、部分机械设备等品类。对于法国这样的航空出口大国，由于飞机出口占其对美贸易的相当比重，实际有效税率反而可能下降。

这里的核心洞察是：关税政策的真实影响，从来不取决于最高税率，而取决于谁在豁免清单上、谁不在。企业需要建立的不是单一的关税预测模型，而是动态的“豁免-税率”组合评估框架。

2. 巴西成为本轮关税调整的最大“意外受害者”，但经济冲击可能有限

在JPM的国别分析中，巴西的关税变动最为剧烈。美国对巴西单独启动了Section 301调查，理由是“不合理”的贸易实践，提议将巴西对美出口的基准关税提高至25\%。即便考虑咖啡、肉类、稀土、橙汁、飞机等产品的豁免，加上钢铁50\%的关税，巴西的加权平均关税仍将从当前的10.5\%跃升至19\%，升幅达8个百分点。

但JPM同时指出，这一冲击对巴西经济的实际影响可能有限。原因有二：其一，巴西对美国的直接贸易敞口本身不高；其二，巴西对美出口结构中，石油等大宗商品占比上升，而这些产品受关税影响较小。这是一个典型的“政策冲击大、经济影响小”的案例，但它的信号意义不容忽视——美国政府正在将贸易工具从“针对中国”扩展到“针对特定国家”，哪怕这些国家并非传统意义上的贸易对手。

对于投资者而言，巴西关税案例提供了一个重要的分析模板：当关税冲击与贸易敞口不匹配时，真正的风险不在于直接贸易损失，而在于政策不确定性对投资信心的侵蚀。

3. 新关税框架的法律基础可能比IEEPA更脆弱，诉讼将成为常态

JPM在报告中点出了一个关键但容易被忽略的风险点：Section 301关税的法律基础，可能比已经被法院推翻的IEEPA更加脆弱。法院在IEEPA案中的裁决逻辑强调“成文法文本优先于行政解释”，这意味着，如果Section 301的适用被认定为超出了该条款的历史使用范围，它同样可能面临司法挑战。

历史数据显示，Section 301通常用于针对特定国家的特定贸易行为，而非像当前这样，对60个经济体施加近乎普遍性的关税。美国政府试图用这个工具来重建IEEPA时期的关税体系，但这一做法是否合法，最终将由法院裁决。JPM的判断是：无论结果如何，诉讼过程本身将显著推高贸易政策的不确定性。

这里的含义是：关税政策的不确定性，正在从“会不会加”演变为“加了之后会不会被法院推翻”。企业不能再

[... middle omitted ...]

任何关税税率。美国贸易代表Greer表示，调查预计在“数周内”完成。关键在于，这项调查的关税是否会与现有关税叠加。如果答案是肯定的，那么对于某些行业——尤其是受政府补贴影响较大的制造业——将面临关税成本的指数级上升。

这份报告没有完全展开的是：结构性产能过剩调查的关税设计逻辑，可能与传统关税完全不同。它可能不是简单的从价税，而是针对特定产品、特定企业的惩罚性税率。这种“定制化”关税的执行难度和合规成本，将远高于当前的普遍性关税。

JPM的这份报告在关税的“第一层影响”——即直接税率变化——上提供了极为详尽的分析，但有两个问题没有得到充分回答。

第一，关税的传导效应。当前的实际税率远低于名义税率，一个重要原因是企业通过调整进口来源、产品组合和定价策略来规避关税。但随着Section 301框架的全面铺开，这种套利空间正在收窄。当更多的经济体被纳入关税体系，企业的“避税”选择将急剧减少，实际税率可能加速向名义税率收敛。这一传导过程的速度和幅度，报告并未给出明确判断。

\subsection{[R010] Citi：黄金股回调是表象，真正的机会藏在现金流与远期定价的错位中}
{\small\color{refgray}归类：黄金与矿业：黄金股回调是表象，矿业主权供需重塑竞争格局}

Citi：黄金股回调是表象，真正的机会藏在现金流与远期定价的错位中

当金价从2026年1月的高点回落约11\%，黄金股却经历了高达25\%的修正。市场直观地将此解读为“金价跌了，黄金股自然要跌”。但Citi这份最新研报揭示了一个被忽视的核心判断：黄金股当前的价格，并非在定价金价的短期波动，而是在错误地定价了这些公司未来两年持续改善的自由现金流。 真正值得关注的，不是金价还能不能回到高点，而是这些公司能否在当前的运营效率下，将现金流转化为估值重估的杠杆。

这份发布于2026年6月8日的报告，覆盖了AngloGold Ashanti和Gold Fields两家标的，均维持“买入”评级。但报告最引人深思的，并非简单的评级，而是Citi在调降目标估值倍数（从8倍降至6-7倍）的同时，仍大幅上调了远期EBITDA预期，并给出了60\%以上的预期总回报率。这种“降倍提价”的微妙组合，恰恰点出了当前市场定价中最核心的错位。

市场对黄金股的担忧，往往集中在金价下跌对盈利的直接侵蚀。Citi的数据显示，自1月高点以来，金价从约5000美元/盎司回落至当前约4350美元，跌幅约13\%。但同期两家公司的股价跌幅却远超金价，达到20\%-25\%。这种非对称下跌，意味着市场正在为“金价继续下跌”这一情景支付过高的保险溢价。

然而，Citi的Q1 2026运营数据给出了相反的信号。两家公司的运营表现均“稳定”，这并非一句套话。在成本端，能源价格高企带来的通胀压力尚未实质性地传导至矿山运营成本；在产量端，Obuasi矿山的爬坡和Nevada项目的推进，正在为AngloGold提供内生增长。Gold Fields的Windfall项目同样在按计划推进。

这意味着，当前股价中隐含的“盈利崩塌”假设，与正在发生的“运营稳定”现实之间存在显著背离。Citi的测算显示，即便在当前的黄金现货价格下，两家公司的自由现金流收益率分别达到10.5\%（AngloGold）和12\%（Gold Fields），这已经显著高于矿业股历史上8\%的平均水平。换句话说，当前股价反映的不是“金价下跌”，而是“市场不相信运营能持续”。 而这份报告的核心洞察在于：运营韧性正在被低估。

2. 真正的估值锚不在当前金价，而在2027年的5000美元

Citi的研报中有一个容易被忽视的结构性判断：Citi house views on gold prices are constructive for \textbackslash{}\$5,000/oz in 2027。 这是一个极为关键的假设，因为它决定了估值的时间维度。

当前市场更多在交易“金价还会跌多少”，而Citi在交易“2027年金价能达到多少”。这两者的区别，决定了估值方法的选择。如果只看当前金价，AngloGold 4.7倍和Gold Fields 3.6倍的一年前瞻EV/EBITDA似乎合理，甚至略贵。但如果将视角拉长至2027年，情况截然不同。

Citi的模型更新显示，由于上调了 年的金价假设（分别上调8\%/25\%/13\%），两家公司的EBITDA在2027年将出现24\%的同比跃升。这意味着，当前的低估值倍数并非“贵”，而是对远期盈利增长缺乏定价。Citi虽然将目标估值倍数从8倍下调至6-7倍，但这一倍数仍然高于各自的历史长期均值（AngloGold 4.7倍，Gold Fields 5.1倍）。这种“高于历史均值”的定价，所对应的正是2027年5000美元金价下的盈利水平。

*所以，Citi的投资逻辑并非“金价涨了所以买”，而是“当前市场定价的是金价下跌的尾部风险，而非金价上涨的确定性收益”。** 这种不对称性，是当前黄金股最吸引人的地方。

报告中反复出现一个数字：自由现金流收益率。AngloGold约10.5\%，Gold Fiel

[... middle omitted ...]

予更高的估值倍数。 Citi的研报明确指出，两家公司的资本开支峰值“很可能已经过去”。这正是估值重估的催化剂。

但这里有一个尚未完全展开的问题：超额现金的分配优先级是什么？ 是优先投入增长项目，还是优先回报股东？报告没有给出明确的指引。这取决于管理层对资本配置的决策，而这一决策将直接影响股价重估的节奏和幅度。

Citi在研报中做了一个看似矛盾的操作：上调了目标价（AngloGold从120美元升至130美元，Gold Fields从65美元降至58美元），但下调了目标估值倍数（从8倍降至6-7倍）。对于Gold Fields，目标价甚至出现了下调。

这种“降倍提价”的背后，反映的是Citi对短期风险的谨慎与对长期价值的乐观之间的平衡。报告明确提到，下调倍数“反映了对黄金的短期谨慎看法”。但问题在于：如果2027年金价真的达到5000美元，当前6-7倍的估值倍数是否仍然合理？ 如果市场届时对黄金股的信心恢复，估值倍数回到8倍甚至更高，那么当前的目标价可能仍然保守。

\subsection{[R011] NOM：全球资金正在重新定价“美国例外论”，但亚洲的撤退信号更值得警惕}
{\small\color{refgray}归类：区域市场：新兴市场融资货币单一化，亚洲资金撤退信号值得警惕}

NOM：全球资金正在重新定价“美国例外论”，但亚洲的撤退信号更值得警惕

全球资金流向正在讲述一个比表面数据更复杂的故事。这份NOM在2026年6月9日发布的亚洲外汇周报，提供了一组高频资金流动数据，覆盖截至6月5日至6日的五个交易日。如果只扫一眼标题，你可能会得出一个简单的结论：资金仍在涌入美国，新兴市场正在失血。但真正重要的不是这个方向判断，而是流动的结构性变化——AI乐观情绪驱动的美国权益资金流入依然强劲，但债券市场几乎无人问津；新兴市场资金连续第二周转为净流出，且流出的速度在加快；韩国散户在连续两个月净卖出后开始回补，台湾投资者却在加速抛售美元债券。

这些信号叠加在一起，指向一个更深刻的判断：市场对“美国例外论”的定价正在从全面拥抱转向选择性下注，而亚洲市场的资金撤退可能才刚刚开始。这不是一次简单的风险偏好切换，而是一次资产定价逻辑的重新校准。

截至6月5日的五个交易日内，外国投资者向美国权益基金净投入25.6亿美元。这个数字虽然略低于前一周的29.1亿美元，但如果拉长到月度维度，5月美国权益基金的外资流入总额达到87.8亿美元，是自2026年1月（95.6亿美元）以来的最高单月水平。

表面上看，这是一个令人振奋的信号。但真正值得追问的是：这些钱到底流向了哪里？

NOM的数据没有拆分行业，但结合近期市场叙事，答案几乎是确定的——AI相关资产。从2025年底到2026年上半年，AI主题始终是驱动全球资金流向美国的核心引擎。问题在于，当AI乐观情绪成为唯一支撑美国权益流入的支柱时，这个市场就变得极其脆弱。一旦AI叙事出现任何裂痕——无论是监管收紧、技术突破放缓，还是商业化不及预期——这些资金流出的速度可能比流入时更快。

更值得警惕的是债券市场。同期美国债券基金的外资流入几乎为零，只有600万美元，而前一周还有2.09亿美元。5月全月债券流入也仅有7.76亿美元，远低于1月的26.7亿美元。权益和债券之间的资金流向出现了罕见的分化——投资者愿意为AI故事买单，但对美国利率环境和信用风险的态度却异常冷淡。

这意味着什么？市场对美国的定价已经从“全面看多”变成了“选择性下注”。投资者在押注AI，而不是押注美国。如果AI是唯一的理由，那么一旦这个理由动摇，整个资金流向结构就可能崩塌。

如果说美国权益市场还在AI的支撑下维持体面，那么新兴市场的情况则令人担忧。

截至6月5日的五个交易日内，新兴市场ETF的外资净流出达到2.5亿美元，而前一周还是净流入4.75亿美元。这个逆转本身并不罕见，但结合月度数据看，趋势更加清晰：5月新兴市场ETF的外资流入为24.7亿美元，远低于4月的62.9亿美元，更不用说2月的191亿美元和1月的255.5亿美元。

从1月的255.5亿美元到5月的24.7亿美元，新兴市场资金流入的萎缩幅度超过90\%。这不是一次正常的周期性波动，而是一次结构性的撤退。

更值得关注的是结构。NOM的数据显示，新兴市场资金流出主要集中在权益端，债券端仍有小幅净流入。这暗示投资者并不是全面看空新兴市场，而是在进行风险排序——先减持估值更高、波动更大的权益资产，保留相对安全的债券敞口。但这种“先撤权益”的顺序，通常是一轮更大规模撤退的前奏。如果美国利率继续维持高位，或者AI叙事出现扰动，债券端的流入也可能迅速转为流出。

对于亚洲市场而言，这个信号尤其重要。新兴市场资金撤退的下一步，往往就是亚洲货币和利率资产的重新定价。

韩国散户的数据提供了一组有趣的对照。在截至6月6日的五个交易日内，韩国散户净买入美国资产1.39亿美元，其中美国权益净买入2.35亿美元，美国债券净卖出0.95亿美元。这个数字看起来像是一个乐观信号——散户在连续两个月净卖出后开始回补。

但仔细看，这个“回补”的规模非常有限

[... middle omitted ...]

国散户的数据，重点不在于他们买了多少，而在于他们什么时候开始大规模卖出——那往往是市场情绪达到极端位置的信号。

截至6月5日的五个交易日内，台湾投资者净卖出美元计价债券基金1.96亿美元，高于前一周的1.49亿美元。拉长到月度维度，5月台湾投资者净买入1.83亿美元，但4月净卖出3.39亿美元。这个数据看起来波动不大，但结合历史趋势看，台湾投资者对美元债券的配置行为正在发生结构性变化。

回顾2024年，台湾投资者几乎每个月都在大举买入美元债券，2024年6月单月流入高达58亿美元。但进入2026年，这个趋势明显放缓，甚至在3月和4月出现了连续净卖出。台湾投资者是全球美元债券市场最大的零售力量之一，他们的行为变化往往领先于机构投资者的配置调整。

为什么台湾投资者在抛售？最可能的解释是，他们对美元利率走势的判断发生了变化。如果市场预期美联储将在更长时间内维持高利率，那么持有美元债券的吸引力就会下降——尤其是当新台币对美元汇率存在不确定性时，汇兑损失可能侵蚀债券收益。

\subsection{[R012] JPM：中国新ODI框架对全球银行的冲击比预期更广，财富管理业务的增量摩擦才刚刚开始}
{\small\color{refgray}归类：区域市场：新兴市场融资货币单一化，亚洲资金撤退信号值得警惕}

JPM：中国新ODI框架对全球银行的冲击比预期更广，财富管理业务的增量摩擦才刚刚开始

这份JPM最新研报揭示了一个被市场低估的关键信号：中国新的境外直接投资（ODI）监管框架，其影响范围远不止于跨境资金流动本身。该机构在与法律专家的深入讨论后，得出了一个比此前预期更为悲观的结论——无论是定义中国内地居民的范围，还是对存量海外资产的追溯，乃至对境外收入的覆盖，都意味着全球银行在亚洲财富管理业务上的增量摩擦可能才刚刚开始计价。

为什么这个判断现在值得关注？因为全球银行正处在财富管理业务贡献快速提升的关键阶段。JPM数据显示，HSBC和渣打的财富收入在2025年都实现了24\%的同比增长，而未来几年财富业务在集团收入中的占比预计将持续攀升。如果监管框架的执行力度超出预期，这些银行面临的不仅是新增客户的放缓，更是存量资产合规成本的系统性上升。

这份研报的核心价值在于，它通过法律专家的解读，将抽象的政策框架翻译成了具体的业务冲击路径。以下是我们从这份报告中提炼出的五个关键洞察，以及一个报告尚未完全回答的关键问题。

1. 中国内地居民的定义远比想象中宽泛，香港身份证不再是“豁免通行证”

市场此前普遍认为，持有香港身份证的个人可以被视为香港居民，从而规避ODI框架的约束。但JPM邀请的法律专家给出了截然不同的解读：只要个人仍持有中国内地户籍（户口），且每年在内地停留超过183天，即使持有香港身份证（包括永久居民身份），仍可能被认定为“中国内地居民”。

这个定义的改变意味着什么？它直接扩大了监管对象的人口基数。许多在香港工作、持有香港身份证但未注销内地户籍的高净值人士，将被纳入监管范围。对于银行而言，这意味着他们需要重新评估客户分类体系，而不是简单地把持有香港身份证的客户视为“非内地居民”。

从竞争格局角度看，那些在香港财富管理业务中高度依赖内地访客（MCV）的银行，受到的冲击将更为直接。JPM的测算显示，如果不再有新的内地访客贡献，HSBC和渣打的EPS将分别面临约4\%和5\%的下行风险。而考虑到定义范围的扩大，这个风险可能还未完全反映在股价中。

市场此前关注的焦点是新增资金流动的监管，但JPM的研报揭示了一个更令人担忧的可能性：现有海外金融资产也可能被纳入监管范围。法律专家预计，新的实施细则可能要求中国内地居民对现有海外金融资产进行备案披露，包括历史持仓和新增境外收入。

这意味着银行需要面对的不仅是新增客户的合规审查，还有存量客户的资产合规检查。对于HSBC这样拥有保险制造业务的银行，其存量保险资产的合规风险尤为突出。相比之下，渣打采用第三方分销模式，在保险资产上的暴露程度相对较低。

这里的关键洞察是：合规成本的上升不是一次性的，而是系统性的。银行需要建立持续的监控体系，对存量客户进行重新分类和审查，这将对成本收入比产生长期压力。JPM的分析师明确指出，直到监管细则明确，银行的股权成本将保持在较高水平。

3. 境外IPO退出和子公司收入成为新焦点，瑞士银行的超高净值业务面临更大不确定性

对于UBS和宝盛这样的瑞士银行，其核心业务是为超高净值客户提供跨境财富管理服务。这些客户的资产来源往往包括海外IPO退出所得、境外子公司利润等。根据JPM的法律专家解读，这些“合法生成的境外现金流”现在也可能被纳入ODI框架。

这个判断的影响是双重的。一方面，新增资金流入将面临更高的合规门槛，导致净新增资金的增速放缓。另一方面，如果存量资产也被纳入监管，银行需要投入大量资源进行合规审查，而这些成本最终可能转嫁给客户或侵蚀利润率。

JPM的分析师认为，相比HSBC和渣打，瑞士银行在信息披露方面更为有限，因此市场对这类风险的定价可能更加不充分。报告指出，超高净值客户的净新增资金面临的下行风险高于此前预期，这是一个需要投资

[... middle omitted ...]

1日后陆续出台，但时间表仍不明确。在此之前，银行面临的将是更高的尽职调查要求和合规成本。

对于投资者来说，这意味着需要区分“长期利好”和“短期压力”的节奏。那些能够迅速适应新监管环境、建立高效合规体系的银行，可能在长期竞争中占据优势。但在短期内，所有银行都需要面对成本上升和业务增长放缓的双重压力。

尽管这份研报提供了比市场预期更悲观的判断，但它也留下了一些关键问题没有完全回答。其中最重要的问题是：监管机构到底有多大的决心和执行力？

法律专家指出，中国监管机构已经通过CRS（共同申报准则）获得了相当广泛的海外资产信息，并且可能利用AI工具进行数据交叉比对。这意味着监管机构对存量资产的掌握程度可能超出市场想象。但问题在于，这些数据将如何被使用？是作为威慑工具，还是会真正启动大规模的合规审查？

另一个未解问题是“重大性门槛”和“祖父条款”的具体设定。法律专家预计，小额历史投资可能被豁免，但门槛的具体数值和过渡期安排仍然模糊。这些细节将直接决定银行存量资产的合规成本。

\subsection{[R013] GS：日本通胀的“二阶导”拐点比绝对水平更重要}
{\small\color{refgray}归类：宏观与利率：通胀但不加息，市场定价逻辑的重构}

五月日本国内企业商品价格指数（CGPI）环比上涨0.9\%，较四月2.8\%的历史性飙升大幅回落。这份来自GS日本经济团队的月度数据解读，看起来只是一次常规的通胀追踪，但真正值得关注的信号，隐藏在环比增速的减速幅度和品类扩散的边际变化中。

市场通常关注通胀的绝对水平——CPI同比是否破3\%，CGPI是否创历史新高。但GS这份报告揭示了一个更微妙的判断：日本通胀正在经历从“全面跳升”到“结构性分化”的转折。四月的价格跳升是一次性的定价调整冲击，五月的读数才是判断持续性通胀压力的真实起点。

对于任何关注日本资产定价、日元走势以及日本企业定价权可持续性的投资者来说，理解这个“二阶导”拐点，远比盯着同比数字的绝对值更重要。

四月CGPI环比飙升2.8\%，是1980年4月以来（剔除消费税上调影响）的最大单月涨幅。这个数字本身容易引发对通胀失控的担忧。但GS报告明确指出，五月环比增速回落至0.9\%，是因为“四月许多品类因价格修订出现的急剧上涨在五月暂时平息”。

这意味着四月的数据更多反映的是企业集中调价窗口期的一次性行为，而非需求驱动或成本持续攀升的线性外推。五月的0.9\%环比增速，虽然仍处于高位，但已经回到更接近趋势性的水平。真正的判断问题应该是：0.9\%的月环比增速是否可持续？它背后是成本推动还是需求拉动？

GS报告没有直接回答这个问题，但提供了关键线索：五月环比增速的减速是广泛性的，几乎所有主要品类都在减速，而非少数品类异常回调。这说明四月跳升后的“均值回复”是系统性的，而非结构性力量在减弱。

对于企业决策者而言，这意味着不应将四月的数据作为未来成本预测的基准。对于投资者而言，五月的读数才是评估日本通胀路径更可靠的起点。

2. 石油制品仍是核心变量，但“三个月定价周期”的机制值得密切关注

石油及煤炭制品环比上涨3.0\%，拉动整体指数0.21个百分点，但相较于四月11.8\%的环比涨幅和0.75个百分点的拉动，已显著减速。更关键的信息在于石脑油：五月石脑油价格环比持平，而四月曾暴涨83.2\%。

GS报告特别指出，石脑油价格通常每三个月重新定价一次，下一次定价调整预计在七月。这意味着七月可能再次出现一轮集中的价格上行压力。这个机制细节对于理解日本通胀节奏至关重要——它不是连续平滑的，而是脉冲式的，与特定的定价周期挂钩。

如果七月石脑油价格再次大幅上调，将直接传导至汽油、化工产品等下游品类，可能再次推高CGPI月度环比读数。但这次冲击的幅度取决于四月至七月期间国际原油市场的走势以及日元汇率的变动。GS报告没有给出预测，但点明了这个关键的时间节点和机制。

对于关注日本通胀路径的投资者，七月的数据将是验证“四月跳升是否只是孤立事件”的下一个关键观察窗口。

3. 进口价格减速但同比仍在加速，日元汇率是未被充分讨论的放大器

日元计价的进口价格五月环比上涨2.7\%，较四月8.7\%的环比涨幅显著减速。但同比增速从21.0\%加速至25.5\%。GS报告将此归因于石油相关产品价格的推动，但同时指出化工产品中的乙烯五月环比下跌4.1\%，而四月曾暴涨69.0\%。

这个数据组合揭示了一个重要但容易被忽视的结论：进口价格的压力正在从“全面扩散”转为“品类分化”。石油相关产品仍是主要推动力，但其他工业原料的价格压力正在缓解。如果这种分化趋势持续，意味着日本输入性通胀的广度正在收窄，但深度（石油品类）仍维持高位。

日元汇率在这一过程中的作用值得进一步拆解。GS报告没有展开讨论，但日元计价的进口价格不仅反映国际商品价格，还叠加了汇率变动。如果日元在此后出现进一步贬值，即使国际商品价格稳定，进口价格仍可能继续攀升。反之，如果日元升值，进口价格压力将显著缓解。这个汇率与通胀的互动关系，是GS这份月度数据报告没有完全展开、但实际影响巨大

[... middle omitted ...]

化成本。

电气电子设备出口价格同比上涨35.6\%，这个数字需要放在日元贬值的背景下理解。但即使剔除汇率因素，以美元计价的出口价格很可能也在上涨。这意味着日本企业在全球市场上的定价能力正在增强，尤其是在半导体设备、精密零部件等具有技术壁垒的领域。

这个趋势如果持续，对日本股市的盈利预期将产生深远影响。日本企业过去二十年最大的结构性问题是“赚了营收但赚不到利润”，而出口定价权的提升直接触及这个核心矛盾。GS报告没有展开讨论这个含义，但数据本身已经提供了足够强的信号。

5. 报告尚未完全回答的关键问题：通胀传导到消费端的时滞与力度

GS这份报告聚焦于企业层面的价格指数，但并未直接讨论CGPI向消费者价格指数（CPI）的传导。这是一个需要谨慎对待的缺口。

四月CGPI的跳升还没有完全反映到五月CPI数据中。企业定价调整与消费者价格之间的传导存在1-3个月的时滞。这意味着即使五月CGPI环比增速放缓，未来几个月的CPI数据仍可能因四月CGPI跳升的滞后影响而继续上行。

\subsection{[R014] JPM：霍尔木兹海峡的“例外通行”不会改变全球LNG市场的结构性缺口}
{\small\color{refgray}归类：能源与大宗：供给硬约束重塑定价权，中国出口成金属市场泄压阀}

JPM：霍尔木兹海峡的“例外通行”不会改变全球LNG市场的结构性缺口

这份JPM在6月8日发布的全球LNG供应与航运追踪报告，表面上在更新一艘船的航迹——ADNOC的Mubaraz号在4月26日成为冲突后首艘满载驶出霍尔木兹海峡的LNG船，如今又重新进入海峡、靠近Das Island准备装货。但真正值得关注的判断藏在数据背后：即便海峡出现零星通行，JPM仍将其定性为“例外而非新常态”，并维持卡塔尔LNG产能到9月才能恢复到83\%利用率的预测。这意味着2026年全年卡塔尔LNG产量将同比腰斩，从2025年的1.05亿吨降至仅0.6亿吨。市场可能正在低估这一供给缺口的持续时间和定价影响。

这份报告的核心贡献，不是告诉你霍尔木兹海峡有船在动，而是提供了一个框架：把“地缘政治事件”转化为“可量化的供给约束”，并由此推演出JKM（东北亚LNG现货价格）对TTF（欧洲天然气基准价）的溢价将持续存在，甚至需要更高的价格来抑制亚洲需求、支撑欧洲储气。对于产业决策者和资产配置者而言，真正需要关注的不再是“冲突何时结束”，而是“后冲突时代的LNG定价机制正在被重写”。

以下是我们从这份报告中提炼出的五个关键洞察，以及一个尚未被充分讨论的变量。

1. 霍尔木兹海峡的“例外通行”无法改变卡塔尔产能恢复的硬约束

JPM追踪到，截至6月8日，霍尔木兹海峡内有9艘可用的LNG船。Al Daayen号当天驶出海峡、目的地为中国，而ADNOC的Mubaraz号在6月3日重新进入海峡。这些信号容易被解读为“供给正在恢复正常”，但报告明确警告：这些通行是例外，而非新常态。

核心约束来自三个层面。第一，海峡内的LNG船数量有限，这天然限制了可立即装运的货物量。第二，即便海峡完全开放，卡塔尔的液化设施恢复也需要时间。JPM引用了卡塔尔能源CEO的最新表态：即使海峡完全开放，未受损的液化列车恢复出口“至少还需要两到三个月”。第三，报告维持其3月发布的判断——重启卡塔尔LNG设施将取决于复杂的运营、安全和战略考量，且需要两到三个月才能恢复全部产出，这还不包括两个受损的液化列车。

将这些约束量化后，JPM的模型假设卡塔尔能源到9月才能恢复到83\%的利用率，这意味着2026年全年利用率仅为57\%，较2025年的105\%近乎腰斩。这不是一个“短期扰动”的叙事，而是一个持续全年的结构性供给缺失。

2. 2026年全球LNG供给增量几乎被卡塔尔的缺口完全吞噬

这份报告最有力的分析之一，是将全球新增LNG项目与卡塔尔现有产能的萎缩放在同一张表格中比对。2026年，全球新增项目预计贡献约83.6亿立方米的增量，其中2025年启动的项目继续爬坡贡献72.5亿立方米，2026年新启动的项目贡献11.2亿立方米。但卡塔尔现有的14条液化列车，2026年产出预计仅为60.2亿立方米，较2025年的112.6亿立方米减少52.4亿立方米。

这意味着，全球所有新增供给的绝大部分，仅仅被用来填补卡塔尔的缺口。全球LNG总供给从2025年的598.9亿立方米微增至2026年的600.1亿立方米，同比增幅仅0.2\%。对于一个需求仍在增长的市场而言，这种供给紧平衡意味着价格必须上行以抑制需求。

这组数据回答了“所以呢”的问题：市场讨论的“新增产能大量投产”在总量层面几乎不产生净增量。卡塔尔缺口才是2026年全球LNG定价的核心变量。

3. 美国LNG流向亚洲的结构性变化正在重塑JKM-TTF价差

报告揭示了一个关键定价信号：JKM目前仍保持对TTF约2美元/百万英热单位的溢价。支撑因素包括即将到来的夏季降温需求，以及潜在的（超）厄尔尼诺天气事件——这通常意味着亚洲夏季气温高于正常水平。

但更值得关注的是美国LNG的流向变化。JPM的数据显示，美国墨西哥湾沿岸对欧

[... middle omitted ...]

格发现”的过程：当前2美元/百万英热单位的JKM-TTF溢价可能不足以完成这一任务，市场需要更高的价差来平衡两大消费区域的竞争性需求。

一个看似矛盾的数据是：LNG航运费率已从冲突第二周的高点暴跌超过60\%。西苏伊士运河区域的即期费率从20万美元/天降至8万美元/天，东苏伊士区域从11万美元/天降至4.5万美元/天。一年期期租费率评估维持在5.05万美元/天。

直观解读是“运力过剩、需求疲软”，但JPM的分析暗示了更复杂的结构。航运费率的下降，部分反映了卡塔尔LNG出口中断导致的长途运输需求减少，而非全球LNG贸易量的整体恢复。更重要的是，报告将航运费率列为JKM-TTF净回值（netback）的重要驱动因素——低航运费率降低了美国LNG运往亚洲的物流成本，这反而可能支撑JKM对TTF的溢价。

换句话说，航运费率的暴跌不是需求疲软的信号，而是地缘政治冲击导致贸易流重新配置的副产品。一旦卡塔尔出口恢复，航运费率可能迅速反弹，届时JKM-TTF价差的逻辑将再次变化。

\subsection{[R015] 摩根斯坦利：新兴市场真正的风险不是基本面，而是融资货币的单一化}
{\small\color{refgray}归类：区域市场：新兴市场融资货币单一化，亚洲资金撤退信号值得警惕}

摩根斯坦利：新兴市场真正的风险不是基本面，而是融资货币的单一化

全球经济正在经历一个罕见的错位阶段。美国增长强劲、利率预期重新定价、美元走强，而新兴市场的基本面却在持续改善——通胀回落、政策信誉修复、经常账户结构优化。这两个方向同时成立，却让大多数投资者的框架出现了盲区。

摩根斯坦利最新发布的全球新兴市场策略报告《Finding Your Funder》，给出了一个反直觉的结论：市场对新兴市场的担忧方向错了。问题不在于是否应该持有新兴市场资产，而在于用什么货币来为这一头寸融资。

这份报告的核心判断是：新兴市场资产本身的吸引力并未减弱，但传统的美元融资方式正在成为最大的约束。如果投资者继续以美元作为唯一的融资货币，即使新兴市场基本面继续向好，回报也可能被汇率波动侵蚀。而选择正确的融资货币组合——尤其是以加元或G3篮子货币融资——可以在不牺牲收益的情况下显著降低波动。

这不是一个关于“买不买”的讨论，而是一个关于“怎么买”的结构性建议。在当前美元强势的背景下，这个视角可能比大多数宏观判断更具实操价值。

1. 美元走强并非新兴市场的终结信号，而是暴露了融资结构的脆弱性

报告开篇就点明了当前市场的核心矛盾。摩根斯坦利团队在5月发布的中期展望中，对新兴市场本地货币资产的基本面持建设性态度，客户也普遍认同两个判断：新兴市场基本面在改善，且全球投资者对新兴市场本地货币资产的配置仍有提升空间。

但分歧出现在宏观层面。美国强劲的就业数据推动利率重新定价，市场对美联储年内降息的预期持续降温，美元指数走强。许多客户认为，美国股市的超额表现和增长前景将继续支撑美元——至少在DXY层面如此。

摩根斯坦利并没有改变其对美元的看空立场，但报告承认，美元走强对新兴市场确实构成显著压力，尤其是考虑到大多数投资者以美元衡量回报。然而，报告的核心洞察在于：美元走强不意味着新兴市场没有价值，而是意味着传统的美元融资方式正在成为回报的拖累。

这里的逻辑值得拆解。新兴市场本地货币资产的回报由两部分组成：资产本身的收益率，以及融资货币的汇率变动。当美元走强时，即使新兴市场资产本身的收益率不变，美元融资者的实际回报也会被侵蚀。这不是新兴市场的问题，而是融资结构的问题。

2. 加元融资的历史表现揭示了一个被低估的选择：它在提升回报的同时降低了尾部风险

报告对过去16年的历史数据进行了系统回测，考察了不同融资货币对新兴市场多头头寸的回报特征。结果中最引人注目的单一货币是加元。

自2010年以来，以加元融资的新兴市场多头头寸实现了年化3.6\%的总回报，波动率仅6.3\%，夏普比率0.58。这一夏普比率接近所有篮子策略中最好的结果。更重要的是，加元融资策略的最大回撤仅为-12.3\%，远低于美元融资的-28.8\%、日元的-25.2\%和瑞郎的-28.7\%。

加元融资的最差3个月和最差12个月表现也优于其他单一货币，分别为约-8.2\%和-8.4\%。这表明加元融资在历史上提供了更均衡的风险收益特征：它在提升回报的同时，没有引入日元或瑞郎那样的左尾风险。

报告指出，除最近一年外，加元融资在所有时间维度上都产生了高于其他单一货币的夏普比率。这一结论在3年和5年的子样本中同样成立，说明加元的优势并非仅仅受益于 年的美元牛市。

加元的这种特性可以从其经济结构找到解释。作为商品货币，加元与全球增长周期高度相关，但又与美国经济深度绑定。这种双重属性使其在新兴市场多头头寸中既提供了增长敞口，又保持了相对的稳定性。对于无法或不愿构建复杂货币篮子的投资者，加元可能是一个值得认真考虑的替代方案。

3. 篮子策略的历史回测证明：分散化融资不是理论，而是持续有效的风险管理工具

如果说加元是单一货币中的佼佼者，那么篮子策略就是整体框架中最稳健的选择。

报告测试了三种

[... middle omitted ...]

势的阶段，篮子融资仍然提供了更好的风险调整后回报。

报告谨慎地指出，这些回测未考虑交易成本，而篮子策略的执行成本确实可能高于单一货币。但即便如此，从风险管理角度看，分散化融资提供的尾部保护仍然具有显著价值。对于大型机构投资者而言，这些执行成本相对于尾部风险的降低而言，是可以接受的代价。

4. 报告尚未完全回答的关键问题：篮子策略的最优权重和动态调整机制

尽管篮子策略在历史回测中表现优异，但报告留下了一个重要问题：篮子权重应该如何设定？

目前报告使用的是等权策略，这虽然简单透明，但在实际操作中可能不是最优解。不同货币在不同宏观环境下的融资效率差异显著。例如，在风险偏好高涨时，商品货币可能更有效；而在避险情绪升温时，日元和瑞郎反而可能成为拖累。

报告没有深入探讨动态权重调整的可能性。一个更成熟的框架可能需要根据宏观经济状态、利率差异、风险偏好等指标，动态调整篮子中不同货币的权重。这种动态策略在理论上可能进一步提升夏普比率，但也引入了更复杂的执行问题和过拟合风险。

\subsection{[R016] HSBC：AI产业真正被低估的不是消费者市场，而是企业级需求的加速爆发}
{\small\color{refgray}归类：AI/算力：资本开支超越互联网泡沫，折旧冲击与融资结构成新焦点 / AI竞争格局：从模型性能转向Agent生态与成本效率，中国软件更具韧性}

HSBC：AI产业真正被低估的不是消费者市场，而是企业级需求的加速爆发

当市场还在争论ChatGPT的用户增长是否见顶、OpenAI能否守住领先地位时，一份来自HSBC的最新研报给出了一个更值得关注的信号：全球AI产业的可寻址市场（TAM）正在被系统性重估，而驱动这一重估的核心力量，并非来自消费者端，而是来自企业级市场的结构性提速。

这份发布于2026年6月的报告，将2026年至2030年全球AI行业累计收入预期上调了38\%，从1.84万亿美元提升至2.55万亿美元。但真正值得注意的不是这个数字本身，而是数字背后的分化：B2B收入预期被大幅上调74\%，而B2C收入预期反而被下调了20\%。

这意味着什么？意味着AI产业的增长引擎正在发生一次静悄悄的切换。过去三年，市场习惯用ChatGPT的月活增长、OpenAI的融资故事、消费者对聊天机器人的好奇心去定义这个行业的想象力。但HSBC的报告清楚地告诉我们：接下来的五年，真正决定行业走向的，是企业客户是否愿意为AI代理、代码工具和行业解决方案买单，而不是消费者是否愿意为聊天机器人订阅费升级。

这不是一个关于“AI是否过热”的判断，而是一个关于“AI的钱到底从哪里来”的结构性重估。

1. B2B市场正在经历一次由Agentic AI驱动的非线性加速

HSBC之所以将B2B收入预期大幅上调74\%，核心依据是Agentic AI（代理型AI）的涌现正在改变企业采用AI的方式。与传统的对话式AI不同，Agentic AI能够自主执行多步骤任务、调用外部工具、与企业现有系统深度集成。这不再是一个“问答工具”，而是一个“数字员工”。

报告特别指出，Anthropic在这一领域的领先地位已经确立。截至2026年5月，Anthropic的年化收入运行率（ARR）已达到470亿美元，而仅仅在2025年底这一数字还是90亿美元。六个月增长超过4倍，这在任何企业软件市场都是罕见的。

更关键的是，Anthropic几乎完全专注于B2B市场，而OpenAI则更依赖B2C收入。这种战略重心的差异，正在转化为市场份额的实质性变化。HSBC估计，在基础模型的企业级细分市场中，OpenAI正在丢失份额，而Anthropic正在接管。

这意味着，我们正在见证AI产业从“技术演示”到“企业工具”的转折点。Agentic AI带来的不是渐进式改进，而是使用场景的质变——当AI可以从“回答问题”进化为“完成任务”，企业客户愿意支付的价格和部署速度都会发生阶跃变化。

与B2B的乐观形成鲜明对比的是，HSBC将B2C市场的收入预期下调了20\%。但“下调”不等于“否定”，这里的核心问题不是用户不再使用AI，而是用户使用AI的方式与企业的变现预期之间存在错位。

报告提到了一个非常具体的现象：OpenAI推出的低价订阅层ChatGPT Go（每月8美元）导致了大量用户从每月20美元的ChatGPT Plus降级。虽然Go套餐包含广告收入，但整体来看，用户的ARPU（每用户平均收入）正在被稀释。

这揭示了一个被市场低估的挑战：消费者对AI聊天机器人的付费意愿并非线性增长。用户愿意为更好的体验付费，但“更好”的定义权不在企业手里，而在用户手里。当企业试图通过分层定价来扩大用户基数时，往往会发现自己正在用低价套餐侵蚀自己的高端用户群。

更根本的问题是，消费者AI产品目前仍缺乏“非用不可”的刚性场景。聊天机器人是“nice to have”，不是“must have”。这与企业级AI形成了鲜明对比——当AI能够替代一个员工的工作流时，企业客户的留存率和付费意愿会完全不同。

HSBC在报告中提出了一个值得深思的判断：全球AI市场正在走向一个由西方大型语言模型（LLM）主导的寡头格局。驱动这一格局的核心

[... middle omitted ...]

新追赶的代价将指数级上升。HSBC认为，只有少数几家公司能够最终留在牌桌上：OpenAI、Anthropic、Gemini，以及从远处观望的Meta。

但报告也指出，欧洲和亚洲正在形成不同的市场结构。欧洲的Mistral AI选择了一条更务实的路径——不追求AGI，而是专注于构建精准的商业用例和深度客户集成。亚洲市场则更加碎片化，独立实验室和科技巨头并行发展，且由于电力供应相对充足，AI采用速度可能更快，但市场集中度可能更低。

对于投资者而言，这意味着“谁能在算力竞赛中持续融资”正在成为比“谁的模型更聪明”更重要的竞争维度。算力壁垒正在重塑AI产业的竞争规则：从技术竞赛转向资本竞赛。

HSBC将OpenAI在2026年至2030年的融资缺口预测从1540亿美元大幅下调至770亿美元，降幅达到50\%。这一调整来自三个因素：自由现金流赤字低于预期（减少570亿美元）、新增股权融资120亿美元、以及AMD股价上涨使得OpenAI可出售的AMD股票价值增加了80亿美元。

\subsection{[R017] JPM：中国车企在欧洲的真正优势不是价格，而是零售执行与生态系统的护城河}
{\small\color{refgray}归类：新能源与汽车：中国车企欧洲优势不在价格，储能市场供给端结构性收紧}

JPM：中国车企在欧洲的真正优势不是价格，而是零售执行与生态系统的护城河

市场对中国汽车出海的讨论，长期集中在两个叙事上：一是关税壁垒，二是低价倾销。但JPM最近一份基于欧洲实地调研的报告，提供了一个截然不同的判断框架。

这份报告的核心发现是：中国车企在欧洲的份额增长，已经不是简单的“便宜BEV出口”故事。2026年4月，中国品牌在欧洲新能源车市场的份额达到18.3\%，是去年同期的三倍以上。更关键的是，这一增长背后，是三个被市场低估的结构性变量在同时起作用——产品组合策略、零售执行能力和充电生态建设。

报告分析师团队在伦敦和巴黎的BYD门店进行了实地走访，并与欧洲汽车大会的行业参与者进行了深度交流。基于这些第一手信息，JPM将比亚迪和吉利2026年的海外销量预测分别上调了约15\%和30\%。

以下是我们从这份报告中提炼出的四个关键洞察，以及一个报告尚未完全回答的关键问题。

1. 欧洲不是一个纯BEV市场，中国车企正在用“混合动力+品牌阶梯”组合赢得过渡期

市场普遍认为中国电动车出海就是卖纯电，但JPM的门店调研揭示了一个截然不同的现实：在欧洲这个“混合动力过渡市场”，PHEV的贡献正在快速上升。

在巴黎和伦敦的BYD门店，销售人员反馈，订单中PHEV占比高达50\%到60\%。最畅销的PHEV车型Seal U DMi起售价约4万欧元，月供仅600-700欧元，等待时间约一个月。而纯电车型Sealion 7虽然也被定位为“高性价比”产品，但在巴黎门店的等待时间长达三到四个月，表明供应端存在明显瓶颈。

这意味着什么？中国车企的竞争优势并非来自单一价格战，而是来自“多动力总成组合+品牌阶梯”的产品矩阵。比亚迪从主流品牌到高端品牌Denza的布局，吉利通过极氪等品牌的多层次覆盖，都在执行同样的逻辑。

对行业竞争格局的含义：传统欧洲车企面临的不再是“一个便宜BEV对手”，而是一个能同时提供BEV、PHEV和HEV的完整产品组合，并且通过品牌阶梯不断向上延伸的竞争者。这种竞争压力是全方位的，而不是局部的。

2. 零售执行正在成为中国车企最深的护城河：渠道密度、试驾转化率和残值管理

JPM报告中一个容易被忽视但极具洞察力的发现是：中国车企正在把“零售执行”变成核心竞争壁垒，而不是简单的渠道铺设。

BYD在欧洲的扩张策略强调物理门店密度和试驾转化率。报告提到，BYD的试驾转化率高达45\%到50\%，这意味着每两个进店试驾的消费者中，就有一个最终下单。同时，BYD强调残值管理和月供竞争力，而不是依赖MSRP折扣竞争。门店调研显示，BYD大多数车型的折扣仅在5\%到7\%之间，销售人员解释这是“营销惯例”，而非价格战。

这种“快零售+服务信用”的打法，与传统出口分销模式形成了鲜明对比。传统车企的经销商网络往往更注重库存管理和价格维护，而中国车企正在把零售端的每一个环节都变成可量化的增长引擎。

对竞争格局的含义：传统欧洲车企如果要应对这种竞争，不仅需要调整产品，还需要重新配置经销商网络的经济模型和库存策略。这比单纯降价要困难得多，也慢得多。

报告中一个令人意外的发现是，BYD正在将充电基础设施（包括闪充技术）作为整体拥有成本体验的一部分来构建，而不仅仅是卖车。

BYD管理层表示，计划到2026年底在欧洲部署3,000个闪充桩，其中英国300个。更值得注意的是，在巴黎门店，销售人员反馈，尽管Denza Z9 GT尚未在欧洲正式发布，仅有一个约11.5万欧元的初步定价，但已经收到了大量订单，等待时间长达六个月。这款车将搭载BYD最新的闪充解决方案，SOC从10\%充到97\%仅需9分钟。

这意味着什么？如果BYD的充电网络能够按计划落地，它将从根本上改变消费者的购买决策逻辑。买家将不再仅仅比较车辆价格和配置，而是比较“

[... middle omitted ...]

KD模式，与马来西亚宝腾、西班牙福特等合作。两种模式各有优劣，但共同点是都在加速本地化布局。

对投资者的含义：本地化能力将成为区分中国车企长期竞争力的核心指标。能够有效规避关税和非关税壁垒、同时保持成本优势的企业，将在海外市场获得更可持续的增长。

报告尚未完全回答的关键问题：充电网络执行能否真正改变品牌溢价？

JPM对BYD充电网络的描述令人振奋，但报告本身也承认，执行将是关键。3,000个闪充桩在欧洲的部署，涉及土地获取、电网接入、当地政策审批等一系列复杂问题。BYD能否在2026年底前完成这一目标，目前仍是一个需要持续验证的假设。

更重要的是，充电网络的部署能否真正转化为品牌溢价？Denza在巴黎门店的初步订单数据令人印象深刻，但这批早期购买者可能是品牌先行者或技术尝鲜者。当Denza正式发布并进入大众市场后，消费者是否愿意为“闪充”支付11.5万欧元的高价，与保时捷Taycan直接竞争？这个问题，报告没有给出答案，但值得所有关注中国车企出海的人持续跟踪。

\subsection{[R018] JPM：香港房价已接近全年预测上限，但真正的不确定性不在价格本身}
{\small\color{refgray}归类：地产与消费：成交量回升但结构性分化，二手房成更可靠温度计}

JPM：香港房价已接近全年预测上限，但真正的不确定性不在价格本身

这份JPM最新发布的Property Data Monitor，数据本身并不复杂。但如果我们只把它当作周度数据读，就错过了报告真正想传递的信号。

核心判断是：香港住宅市场正在快速接近JPM2026年全年房价预测区间10\%-15\%的上沿——年初至今已涨9.6\%，距离上限仅0.4个百分点。这意味着，如果当前趋势延续，JPM对全年的基准预测可能在未来几周内被触及甚至突破。但报告真正值得关注的，不是这个数字本身，而是它背后揭示的市场结构：一手楼定价策略正在发生微妙但重要的转变，而内地市场的库存消化节奏，正在从“量”的改善转向“价”的博弈。

以下是我们从这份报告中提炼的五个关键洞察，以及一个报告尚未完全回答的问题。

1. 香港一手楼定价正在从“折让”转向“溢价”，这可能是市场信心分层的信号

过去两年，香港开发商普遍采用“低开”策略——一手楼定价低于或持平于二手市场，以加速去化。但这份报告显示，上周推出的三个新盘中，有两个定价显著高于周边二手物业：Pavilia Rosa（九龙塘）首批均价达33,800港元/平方呎，高出二手市场41\%，创过去五年一手新盘定价最高纪录。即使是定价相对保守的Headland Residences（柴湾）和One Victoria Cove Ph4（红磡），仍分别高出二手市场15\%和18\%。

这组数据意味着什么？不是开发商集体乐观，而是市场正在出现明显的“价格分层”。高端地段（如九龙塘）的一手楼开始敢于定出显著溢价，而中端项目仍在谨慎试水。报告没有直接展开这个判断，但结合其后续提到的深圳高端住宅销售强劲、中端疲软的分化现象，可以合理推断：香港住宅市场正在经历类似的“K型复苏”——高净值买家的购买力依然充沛，但中产阶层的承接力正在减弱。

对投资者的含义：如果这一趋势持续，未来几个月香港住宅市场可能出现“量缩价稳”甚至“量缩价升”的格局——交易量下降，但价格指数因结构性偏向高端成交而继续上行。这会让表面上的价格数据更具误导性。

2. 内地二手房挂牌量下降，但库存月数仍高，真正的拐点取决于“卖方退场速度”是否快于“买方入场速度”

报告最值得关注的内地数据，不是18\%的同比销售增长，而是两个领先指标：一线城市二手房挂牌量环比下降0.7\%，库存月数从18.8个月微降至18.6个月。北京降幅最大（-2\%），上海次之（-1\%）。

从绝对水平看，18.6个月的库存月数仍然较高——行业一般认为12个月是供需平衡线。但边际变化的方向更重要：挂牌量下降意味着卖方在减少，而非买方在增加。这是一个典型的“供给收缩型改善”。报告中的中原经理信心指数持平于54，也印证了市场情绪并未显著升温。

这意味着当前内地房地产市场的改善，更多来自“卖房的人少了”，而不是“买房的人多了”。这种改善的可持续性，取决于卖方退场速度能否持续快于买方入场速度。如果未来几个月政策刺激或经济预期改善带动需求回升，库存月数有望加速下降；但如果需求维持疲弱，供给收缩的边际效应会递减。

对投资者的含义：当前不宜将销售数据改善简单解读为“市场见底”。更值得跟踪的是库存月数的下降斜率，以及中原挂牌价格指数是否会从当前持平状态转向上行。

3. 深圳住宅市场正在出现“AI财富效应”驱动的结构性分化，但这一逻辑能否外推至其他城市仍存疑

报告在Credit View部分引用了JPM在深圳的专家会议信息，这是一个值得单独拆解的洞察。专家指出，深圳住宅市场呈现明显的两极分化：高端住宅销售强劲，背后是AI/股市带来的正向财富效应；中端市场疲弱，低端市场虽有部分成交量恢复，但整体仍不乐观。南山区和福田区表现优于其他区域。

这一观察的重要性在于：它首次明确将“AI财富效应”与住宅需求

[... middle omitted ...]

那么深圳高端住宅开发商和持有者可能受益最直接。但这一判断的验证，需要更多城市层面的数据支撑。

4. 内地资本外流管控收紧对香港楼市的实际影响，可能低于市场担忧

报告提到，上周香港地产板块下跌6\%，跑输恒指，主要原因是市场担忧内地资本外流管控收紧和潜在加息。但JPM的深圳专家会议给出了一个反直觉的判断：专家认为，近期资本管控收紧对内地居民购买香港住宅的实质影响有限。

这一判断的逻辑基础是：内地居民购买香港房产的资金流动，大多数已经通过合规渠道完成，增量管控的边际效应可能不大。此外，香港楼市的本地需求（换楼、投资）仍然存在，并非完全依赖内地买家。

但报告没有完全展开的是：如果管控收紧叠加香港本地加息，两者的叠加效应可能被低估。尤其是当前香港住宅价格已接近全年预测上限，任何负面冲击都可能触发获利回吐。

对投资者的含义：市场对政策风险的定价可能已经过度，但真正的风险在于“政策叠加”而非单一因素。建议关注香港银行体系流动性变化和按揭利率走势，这些才是更直接的领先指标。

\subsection{[R019] 摩根斯坦利：市场低估的并非政策力度，而是成交量回升背后的结构性分化}
{\small\color{refgray}归类：地产与消费：成交量回升但结构性分化，二手房成更可靠温度计}

摩根斯坦利：市场低估的并非政策力度，而是成交量回升背后的结构性分化

过去一周，中国房地产市场的周度数据出现了一组值得注意的信号。摩根斯坦利最新发布的周度数据库追踪报告显示，截至6月7日当周，50城新房注册成交量同比增长23\%，较前一周14\%的同比增速明显抬升。10城二手房成交量同比增幅则从9\%扩大至23\%。

单看这些数字，很容易被解读为“政策见效、市场回暖”。但真正重要的判断藏在数字背后：这轮成交量回升并非全面复苏，而是一场结构性分化正在加速——一线城市二手房正在成为资金真正的“避风港”，而新房市场的反弹高度依赖供给端的节奏控制，并非需求端的自发修复。

这份报告的价值不在于告诉你“市场变好了”，而在于它提供了拆解“变好”背后真实驱动力的数据框架。以下是我们基于报告核心数据的四点解读，以及一个报告尚未完全回答的关键问题。

报告最核心的信号来自二手房市场。10城二手房周度成交量同比增长23\%，年初至今累计增幅已达6\%。更值得关注的是结构：一线城市二手房成交同比增幅从前一周的12\%跃升至34\%，二线城市也从6\%提升至17\%。

这意味着什么？二手房市场正在承接新房市场的部分需求外溢。当购房者对期房交付信心不足、对开发商信用风险仍有顾虑时，二手房“所见即所得”的属性正在转化为实实在在的成交优势。这个趋势在年初至今的数据中已经形成了持续轨迹，而非单周波动。

对产业决策者而言，这一信号意味着：未来判断市场真实温度，二手房成交量可能比新房数据更具先行指标意义。如果二手房成交量能够持续维持在正增长区间，才意味着需求端真正出现了底部企稳的可能。而新房成交的波动，更多反映的是供给节奏和开发商推盘意愿。

新房市场看似亮眼的数据需要谨慎解读。50城新房成交量同比增长23\%，但报告明确提示，这在一定程度上受到去年端午节假期基数效应的影响。更关键的是结构：一线城市新房成交同比仅增长1\%，远低于二线城市的30\%和三线城市的22\%。

这种分化并非偶然。一线城市新房供应本就稀缺，而报告显示当周上海的新房整体去化率达到了100\%。这听起来是好事，但背后的逻辑是“只有少量新盘入市”——当周一线城市去化率从前一周的63\%跃升至100\%，二线城市甚至没有新盘推出。换句话说，新房成交的“好数据”很大程度上是供给端主动控制的结果，而非需求端爆发。

这引出一个更本质的问题：当前新房市场的成交量反弹，有多少是真实需求驱动，有多少是开发商“控盘”的结果？如果供给端放量，成交量能否维持？这份报告没有直接回答，但数据暗示了一个方向——一旦推盘节奏恢复正常，去化率很可能回落，新房成交的同比增速也会随之收窄。

报告中有两个值得追踪的微观指标：中央指数六城二手房报价指数维持在17.5\%的水平，与上周持平；一线城市房产中介指数则从53.8微升至54.4。

报价指数持平意味着卖房者的价格预期没有进一步下修，这可以解读为市场情绪企稳的信号。但问题在于，这个指数仍然处于历史低位区间，说明“以价换量”的格局没有改变。成交量在增长，但成交价格仍在承压。中介指数的微幅回升则反映了交易活跃度改善对中介业务的正面影响，但这个升幅是否可持续，取决于成交量能否在未来几周继续维持。

这两个指标结合在一起，指向一个尚未被充分定价的风险：成交量的回暖可能只是价格持续下修后的“自然出清”，而非价格企稳后的“价值发现”。如果成交量增长的同时价格指数没有同步回升，那么这轮反弹的持续性就要打上问号。

4. 报告尚未完全回答的关键问题：成交量能否在价格企稳前持续？

这是整份报告最核心的未解问题，也是投资者最需要自己判断的变量。

当前的数据格局是：成交量在恢复，但价格仍在低位。历史上，房地产市场的底部通常经历三个阶段：成交量先企稳，然后价格止跌，最后成交量进一步放大。我们现在处于第一

[... middle omitted ...]

价指数能否从“持平”转向“回升”。如果成交量持续增长而报价指数仍然横盘，说明卖房者仍然缺乏提价信心，市场本质上仍是买方市场。

这三个条件中，任何一个未能满足，当前这轮成交量反弹都可能只是“脉冲式”的，而非趋势性的。报告提供了数据，但并未给出判断——这正是需要读者自己深入思考的地方。

基于这份报告的数据结构，我们建议建立一个更系统的观察框架，核心指标可以概括为“三率”：

成交量增速的持续性。关注周度同比增速能否连续四周维持在正增长区间，而非仅看单周数据。当前新房和二手房都只满足单周条件，尚未形成连续趋势。

去化率的供给弹性。关注当推盘量恢复正常后，去化率会回落多少。如果去化率回落幅度超过10个百分点，说明需求端支撑不足。

报价指数的方向变化。关注中央指数是否从“持平”转向“上行”。这是判断市场是否从“以价换量”转向“量价齐升”的关键信号。

这三个指标相互验证，比单一数据点更能反映市场的真实状态。当前的数据在三项指标中只部分满足了第一项，第二和第三项均未确认。

\subsection{[R020] GS：中国数据中心市场真正被低估的不是投资规模，而是电力需求的非线性跃迁}
{\small\color{refgray}归类：特高压与电网设备：订单翻倍背后，竞争从订单量转向议价权}

GS：中国数据中心市场真正被低估的不是投资规模，而是电力需求的非线性跃迁

当市场还在热议“2万亿投资计划”能否落地时，一份来自GS的研报揭示了一个更值得关注的信号：中国国家能源局（NEA）预计，到2030年数据中心电力消费将达到800 TWh。这个数字意味着什么？它意味着中国数据中心用电量将以36\%的年复合增速增长，到2030年将占全国总用电量的6\%，远超当前1.6\%的水平。

更关键的是，这一预期已经超过了GS自己对美国数据中心市场的预测——726 TWh。换句话说，中国正在从“跟随者”变为“引领者”，而市场对这一结构性变化的定价可能还远远不够。

这份报告的核心判断并非“2万亿投资”本身，而是这背后隐含的电力需求非线性增长逻辑。当多数投资者还在纠结于政策能否兑现、建设节奏会不会低于预期时，真正需要追问的问题应该是：如果需求增速真的达到NEA预测的水平，哪些环节会被迫重新定价？

1. 2万亿投资计划不是新闻，真正的新信息是NEA的需求预测比GS自己的模型更乐观

2025年1月，国家发改委、国家数据局、工信部联合发布的《国家数据基础设施建设指引》中，国家数据局就已提出未来五年数据基础设施将吸引约2万亿直接投资。这一数字并非新消息，但市场此前并未充分消化其隐含的需求端含义。

GS研报的核心贡献在于，它将NEA在2026年5月发布的数据中心电力需求预测与自己的模型进行了对比。NEA预计2030年数据中心用电量达800 TWh，而GS自己此前的预测是IT电力需求增速约为20\%的年复合增长率。两者之间的差距，正是市场认知的盲区。

这里的逻辑链条值得拆解：如果数据中心用电量真的以36\%的增速增长，考虑到PUE（电能利用效率）的持续改善，IT设备本身的电力需求增速将更快。这意味着，市场此前关于“数据中心建设过剩”的担忧可能需要重新审视。需求侧的实际增长，可能远超供给侧规划时的预期。

2. 中国数据中心用电量增速已现加速迹象，2026年开局数据提供了早期验证

研报引用了NEA公布的2026年1-4月数据：全国互联网数据服务业用电量同比增长44\%。这个数字本身已经超过了NEA对全年的预期增速，更不用说GS自己的模型。

从历史经验看，这种早期数据的信号意义不容忽视。2024年全年中国数据中心用电量约为170 TWh，如果2026年前四个月已经实现44\%的同比增长，那么全年数据大概率会超出年初的预期。这种“超预期”一旦形成趋势，就会倒逼供给侧加速扩张，进而影响设备采购、电力基础设施投资、乃至电价政策。

对于投资者而言，这意味着关注焦点应从“政策说了什么”转向“数据证明了什么”。当用电量增速持续高于预期时，此前关于“产能过剩”的担忧就会自然瓦解，而相关企业的收入可见度将显著提升。

GS研报提供了一个极为重要的对比维度：到2030年，中国数据中心用电量预计达到800 TWh，而美国为726 TWh。这意味着中国将在这一领域超越美国，成为全球最大的数据中心电力消费国。

这一判断的底层逻辑值得深挖。美国数据中心需求主要受云服务商和AI大模型训练驱动，而中国的驱动力则更为多元：除了AI训练和推理，还有产业数字化、智慧城市、工业互联网等场景。更重要的是，中国的电力基础设施建设和电网扩容能力，在政策推动下具有更强的执行力。

对于全球科技投资者而言，这意味着中国数据中心产业链的估值逻辑正在发生变化。过去，市场习惯于用美国市场的增速和利润率来给中国公司定价，认为中国公司存在“折扣”。但如果中国市场的增速开始超过美国，这种定价逻辑本身就值得质疑。

4. 报告尚未完全回答的关键问题：PUE改善的“二阶效应”可能被低估

GS研报提到“考虑到PUE的改善，这可能会转化为更快的IT电力需求增长”，但没有展开计算。这正是报告留给读者

[... middle omitted ...]

电源模块，从智能运维到新型芯片架构，这些细分赛道的价值可能需要重新评估。

5. 对投资者的观察框架：用三个维度重新审视中国数据中心标的

基于GS报告的逻辑，我们建议投资者从以下三个维度构建自己的分析框架：

第一，需求增速的验证窗口。2026年全年用电量数据将在2027年初公布，在此之前，每个季度的用电量增速都将是重要的先行指标。如果增速持续高于预期，那么整个行业的收入预期就需要上修。

第二，政策执行力的观察点。2万亿投资计划能否落地，关键在于地方政府和电网公司的配合程度。可以关注各省份的数据中心专项规划、电力审批进度、以及土地和能耗指标的分配情况。

第三，竞争格局的分化信号。GS在报告中给出了对具体公司的评级：买入万国数据、世纪互联、朗源股份，中性上海数据港，卖出光环新网。这种分化本身说明，即使行业整体增速超预期，不同公司在订单获取、成本控制、客户质量上的差异也会导致业绩分化。尤其是那些能够获得优质电力资源、与头部云厂商深度绑定的公司，可能拥有更强的定价权。

\subsection{[R021] JEF：储能市场的真正拐点不是增长，而是增长的结构性分化}
{\small\color{refgray}归类：新能源与汽车：中国车企欧洲优势不在价格，储能市场供给端结构性收紧}

这份研报最值得看的判断，并非2026年全球储能装机将增长80\%——这个数字已经在预期之内。真正重要的信号是：储能市场的增长正在从“总量驱动”转向“场景驱动”，而不同场景的竞争壁垒、利润结构和投资逻辑正在快速分化。那些能在AI数据中心储能、工商业储能和户用储能三个赛道同时建立优势的企业，将享受估值溢价；而只依赖单一市场或单一产品的公司，可能面临增长见顶后的残酷洗牌。

JEF在SNEC 2026首日与华为储能专家、SMM专家、远景能源及多家头部企业高管的交流，揭示了一个正在发生的结构性变化：储能不再是光伏的附属品，而是正在成为独立于光伏的、由AI算力需求驱动的全新增长极。这一判断如果成立，将重塑整个产业链的估值体系。

1. 全球储能装机80\%增长的背后，真正驱动力量正在从政策转向算力

JEF专家预计2026年全球储能装机将超过150GWh，同比增长80\%，且2026年后仍将保持20-30\%的年增长率。但比总量更值得关注的是增长的结构：中国仍占全球60\%的份额，但美国AI数据中心储能需求正在成为新的增长引擎。

华为储能产品总监明确指出，美国AIDC储能需求预计2026年翻倍至10GWh，到2030年将增长至70-80GWh。这意味着未来五年，仅美国AI数据中心这一个场景，就将创造约7-8倍的储能增量。这与传统认知中“储能增长靠新能源配储”的逻辑完全不同——AI数据中心储能的核心诉求是电网稳定性、负载平滑和备用电源，而非新能源消纳。

这一判断的产业含义是：储能企业的技术路线选择将出现分化。面向AI数据中心场景，需要的是高可靠性（系统可用性99.999\%）、快速响应和与燃气轮机/柴油发电机的协同能力。那些在电力电子转换技术上有积累的企业，如阳光电源，将比纯粹的电池集成商更具竞争优势。

2. AI数据中心储能正在创造一个新的技术壁垒：系统可用性99.999\%

远景能源储能销售负责人提供了一个关键洞察：AI数据中心储能的核心价值不在于容量，而在于系统可用性。能够开发出支持99.999\%系统可用性的PCS（储能变流器）的企业，将在这一赛道占据先机。

这一要求远高于传统储能场景。风光配储的可用性要求通常在99\%左右，而AI数据中心需要的是接近不间断电源级别的可靠性。这意味着储能系统需要更复杂的冗余设计、更快的切换速度和更强的电网交互能力。

阳光电源的管理层在交流中透露，其固态变压器产品（涵盖10kV至800V、800V至50V）将于2026年下半年推出并立即销售。更重要的是，阳光电源拥有全栈自研能力，能够满足亚马逊、谷歌等超大规模客户对定制化产品的需求。相比之下，许多竞争对手仍依赖供应链采购组件，在定制化能力上存在明显差距。

这一技术壁垒的建立，将导致AI数据中心储能市场在未来2-3年内出现明显的“赢家通吃”特征。拥有完整技术栈和客户定制能力的头部企业，将获得显著的定价权溢价。

3. 工商业储能正在成为利润最丰厚的细分市场，但竞争格局尚未固化

JEF数据显示，2026年全球工商业储能装机预计同比增长66\%至32GWh，中国占60\%。增长驱动力包括美国AI数据中心扩张和中国零碳园区建设。

但更值得关注的是利润率结构。阳光电源管理层透露，其工商业和户用储能2025年出货4GWh，收入约80-100亿元，户用储能净利率约20\%。更重要的是，公司2026年将更加聚焦工商业和户用储能，该业务毛利率可达40-50\%，部分对冲公用事业级储能的利润率压力。

这意味着工商业储能可能是整个储能产业链中利润率最高的环节。原因在于：工商业客户对价格敏感度相对较低，更看重系统可靠性和服务能力；同时，这一市场尚未出现绝对的领导者，竞争格局仍在形成中。

Sigenergy的策略也印证了这一判断。该公司保持溢价定位，产品价格比同行

[... middle omitted ...]

026年全球户用储能出货量同比增长14\%至30GWh，其中欧盟约占50\%。欧盟市场预计2026年同比增长30\%，驱动力包括升级至更大容量系统、阳台储能的快速扩张等。

澳大利亚市场尽管补贴政策逐步退出，但专家仍预计2026/2027年户用储能装机同比增长20\%。Sigenergy管理层也确认，澳大利亚仍是2026年关键市场，尽管补贴自2026年5月起逐步退坡，但6-7月需求仍然强劲，预计全年澳大利亚市场仍将增长至8-10GWh。

这些数据指向一个反直觉的判断：户用储能正在从政策驱动转向经济性驱动。即使在补贴退出的情况下，消费者仍然愿意为储能付费，这暗示户用储能的商业模式已经初步跑通。

但结构变化值得警惕。Sigenergy预计欧盟市场2026年贡献收入60-70亿元，同比增长20-30\%，但这一增长可能更多来自产品升级和容量提升，而非新增用户数。阳台储能等新兴场景正在快速扩张，但单体价值量远低于传统户用储能系统。这意味着户用储能市场的“量增价稳”逻辑可能面临挑战。

\subsection{[R022] JEF：AI竞争已从模型性能转向Agent生态与成本效率的再定价}
{\small\color{refgray}归类：AI/算力：资本开支超越互联网泡沫，折旧冲击与融资结构成新焦点 / AI竞争格局：从模型性能转向Agent生态与成本效率，中国软件更具韧性}

JEF：AI竞争已从模型性能转向Agent生态与成本效率的再定价

这份报告最值得关注的核心判断是：中国AI模型正在从“追赶性能”切换到“定义成本效率”的新阶段，而真正驱动token消耗爆发式增长的不是聊天场景，而是Agent间协作（A2A）与行业专用工具链的落地。JEF在这期AI系列研报中提供了超过50个数据点，但最关键的信号隐藏在Token消耗的环比变化和模型排名的结构性位移中——中国模型在OpenRouter平台上的周Token消耗首次超越美国模型，而排名第一的DeepSeek V4 Flash单周消耗量达到3.69万亿Token，几乎是第二名腾讯Hy3 preview的1.25倍。

这不是一个简单的“中国大模型跑得更快”的故事。JEF的数据揭示出三个正在重塑行业格局的结构性变量：第一，Agent-to-Agent的生态化协作正在替代单一模型的能力竞赛，腾讯与美团、手机厂商的A2A合作是这一趋势的缩影；第二，中国模型在API成本上仅为美国模型的几分之一，这种成本优势在推理量持续放大的背景下，正在改写云服务商和AI应用层的利润模型；第三，OpenAI和Anthropic正在向金融和法律等垂直领域加速渗透，这意味着通用模型的红利期正在收窄，行业专用Agent将成为下一轮竞争的主战场。

1. Token消耗的结构性跃升背后，是Agent从“对话”到“工作流”的范式切换

JEF的数据显示，截至6月1日当周，OpenRouter平台周Token消耗环比增长13.5\%至36.1万亿。表面上看这是用量增长，但真正值得追问的是：什么场景在驱动这种增长？报告明确指出，2026年2月Token消耗的显著跃升有两个直接原因——Coding/Agent采用率的爆发，以及高性能低成本中国模型的密集发布。这意味着Token消耗的增长不是线性的“更多人聊天”，而是“机器替人执行任务”的指数级扩张。

当一个Agent需要调用另一个Agent的能力来完成用户的任务时，每一次交互的Token消耗不是单次对话的线性叠加，而是多Agent协同的多轮嵌套。这正是腾讯正在搭建的Agent-to-Agent生态的本质——用户通过元宝App发起本地服务请求，背后是元宝与美团AI助手小美的A2A通信；用户通过手机语音助手发送微信消息，背后是手机厂商的Agent与微信Agent的协作。每一次“隐形”的Agent间握手，都在消耗Token，但这些Token对用户来说是不可见的。

这对投资判断的含义是：Token消耗的增速本身已经不再是需求指标，而是生态复杂度的映射。那些能够建立Agent间协作网络的公司，将获得不成比例的Token消耗增量，而不仅仅是依赖自有模型推理量的线性增长。

2. 中国模型在OpenRouter上首次超越美国模型，但真正的差异不在性能，在成本结构

报告中最引人注目的数据点是：在6月1日当周，中国模型在OpenRouter上的Token消耗达到14.2万亿，环比增长27.5\%，而美国模型仅为3.2万亿。DeepSeek V4 Flash以3.69万亿的周消耗量排名第一，腾讯Hy3 preview以2.94万亿紧随其后。按公司维度，DeepSeek以18.7\%的市场份额位列第一，超过Anthropic的14.6\%和Google的11.9\%。

但JEF同时引用了Artificial Analysis的智能指数数据：在前20大模型中，中国模型有9个，且均为开放权重。Anthropic Claude Opus 4.8以61分排名第一，GPT-5.5以60分紧随其后，而Qwen3.7 Max以56分位列第五，MiniMax-M3和Kimi K2.6以54分并列第六。按2026年斯坦福AI指数报告，美国顶级模型领先中国模型仅2.7个百

[... middle omitted ...]

开这一推演，但这是投资者需要自行判断的关键风险。

3. 腾讯的A2A生态和阿里巴巴的Qwen Agent平台，正在定义两种不同的生态策略

JEF详细梳理了腾讯和阿里巴巴在Agent生态上的布局，但两者路径截然不同。腾讯的策略是“连接器”——通过元宝App与美团、手机厂商的A2A协作，腾讯不是在做自己的Agent应用，而是在做Agent间的通信协议层。用户通过手机语音助手发送微信消息，本质上是手机厂商的Agent调用了微信的Agent能力，而腾讯在这一过程中既没有获取用户数据，也没有改变微信的产品形态，但微信的Agent接口成为了整个生态的必经节点。

阿里巴巴的路径则更接近“平台型”——Qwen App向第三方Agent和技能开放，首批合作伙伴包括肯德基、瑞幸、蜜雪冰城和中国东方航空。用户可以在Qwen App内一站式获得从点餐到航旅的个性化服务。同时，阿里巴巴在6月1日发布了Qwen3.7-Plus，这是一个多模态交互式混合Agent，统一了GUI和CLI界面。

\subsection{[R023] 摩根斯坦利：AI资本开支已超越互联网泡沫，但市场尚未定价折旧的滞后冲击}
{\small\color{refgray}归类：AI/算力：资本开支超越互联网泡沫，折旧冲击与融资结构成新焦点}

摩根斯坦利：AI资本开支已超越互联网泡沫，但市场尚未定价折旧的滞后冲击

这份摩根斯坦利GVAT策略团队在2026年6月发布的研报，提供了一个关键但容易被忽视的视角：AI基础设施投资的规模不仅在总量上创下纪录，更在资本密集度、融资结构和资产负债表外承诺三个维度上，同时超越了市场此前的认知边界。

报告的核心判断可以概括为一句话：市场当前对AI投资的讨论，仍然停留在“资本开支增速”这个旧框架里，而真正需要被重新定价的，是这些投资在未来三年内如何转化为折旧、如何影响利润率、以及资产负债表外的租赁和采购承诺如何改变行业的风险分布。

这不是一份关于“AI是否过热”的报告。它是一份关于“当资本开支峰值已成定局，哪些财务指标将接力成为新的市场叙事”的路线图。对于产业决策者和机构投资者而言，理解这份报告的框架，比争论资本开支何时见顶更有价值。

1. 资本密集度超过互联网泡沫峰值，但融资结构的变化比数字本身更重要

摩根斯坦利预测，2026至2028年，超大规模云服务商的资本开支与销售比将分别达到36\%、44\%和42\%。作为对比，2000年互联网泡沫时期，光纤相关投资的峰值也不过32\%。这意味着，即便以最宽松的历史标准衡量，当前的投资强度也已经进入了一个前所未有的区间。

但真正值得关注的不是数字本身，而是数字背后的融资结构。报告明确指出，增量资本开支的修订现在往往伴随着资本募集或创新融资安排。融资租赁的使用正在推高名义资本开支数字。当资本开支中的相当一部分来自租赁而非自有现金时，企业的财务灵活性和风险敞口会发生本质变化。

这带来的第一个含义是：资本开支的“真实”规模比账面数字更大。因为融资租赁意味着未来数年内持续的现金流出义务，而这些义务在传统的资本开支统计中往往只体现为一次性投资。当租赁承诺和采购承诺合计超过1.8万亿美元时，企业的资产负债表外杠杆已经达到了一个需要被重新审视的量级。

2. 超过2万亿美元的剩余履约义务，正在将行业命运锁定在相互依赖的网络中

报告中最具冲击力的数据之一，是主要AI计算企业的剩余履约义务在过去一年内几乎翻了三倍，突破2万亿美元。与此同时，长期采购承诺接近1万亿美元，租赁承诺超过8000亿美元。这些数字意味着什么？

它们意味着，AI产业链的上下游已经通过长期合同深度锁定。客户承诺了未来数年的采购量，供应商则据此进行资本开支决策。这种锁定机制在短期内为投资提供了确定性，但也制造了一个潜在的系统性风险：如果需求增速低于预期，这些承诺将转化为巨大的沉没成本和减值压力。

报告展示的AI生态系统互联图揭示了一个高度集中的网络。OpenAI、微软、甲骨文、亚马逊、谷歌、英伟达、CoreWeave等核心参与者之间的资本、租赁和采购关系形成了一个闭环。这种结构的好处是效率极高，坏处是——当任何一环出现问题时，冲击的传导速度也会极快。

对于投资者而言，这意味着传统的“单公司分析”框架已经不够用了。理解一家AI基础设施公司的风险，必须同时理解其客户和供应商的承诺规模和财务状况。报告没有完全展开这一点，但这恰恰是下一阶段研究的关键方向。

3. 折旧是下一个需要关注的利润表变量，而建设中的资本将延长冲击的滞后效应

摩根斯坦利估计，微软、甲骨文、Meta和谷歌在未来三年内的折旧费用合计将超过5200亿美元。这是一个足以改变行业利润格局的数字。

折旧的逻辑很简单：今天投入的资本开支，明天会转化为固定资产，然后通过折旧逐年计入费用。当资本开支以如此规模增长时，折旧曲线的陡峭程度将是史无前例的。如果收入增长无法同步跟上，利润率将面临系统性压力。

但报告指出了一个重要的缓冲机制：在固定资产投入使用之前，大量资本会以“在建工程”的形式停留在资产负债表上。这意味着，折旧冲击不会立即发生，而是会滞后一到三年。这

[... middle omitted ...]

判地位会显著增强。英伟达、CoreWeave等关键供应商的客户集中度极高，而这些长期合同恰恰赋予了它们在定价和交付条件上的议价权。与此同时，采购方虽然获得了供应的确定性，但也牺牲了未来价格下行的灵活性。

这种动态意味着，AI基础设施行业的利润分配格局可能正在发生结构性变化。那些拥有稀缺技术和产能的供应商，正在通过长期合同将自身的竞争优势货币化。而采购方虽然规模庞大，但它们的议价空间可能比市场想象的要小。

报告没有直接讨论这种权力转移的定价含义，但它隐含了一个重要判断：在AI基础设施领域，规模本身并不等于议价权。真正决定利润分配的是技术的不可替代性和产能的稀缺性。

5. 报告尚未完全回答的关键问题：需求端的真实消化能力与折旧冲击的时点

尽管这份报告提供了大量关于供给端的量化数据，但它对需求端的分析相对克制。这并非报告的缺陷，而是反映了当前市场信息的一个天然不对称：我们能够相对准确地追踪资本开支、租赁承诺和采购订单，但很难提前预测AI应用的商业化速度和收入转化率。

\subsection{[R024] 摩根斯坦利：GB200/300机柜出货量环比下降7\%，但这轮调整的真正含义是供应链议价权的再分配}
{\small\color{refgray}归类：AI/算力：资本开支超越互联网泡沫，折旧冲击与融资结构成新焦点}

摩根斯坦利：GB200/300机柜出货量环比下降7\%，但这轮调整的真正含义是供应链议价权的再分配

2026年5月，全球GB200/300 NVL72机柜出货量约为7700台，环比下降7\%。如果只看这个数字，很容易得出“AI服务器需求见顶”的结论。但摩根斯坦利这份最新月度追踪报告传递的信号恰恰相反：2026年全年出货量预计仍将超过7-8万台，同比增长超过100\%。 月度波动不是需求问题，而是供应链节奏问题。

更重要的是，这份报告揭示了三个被市场低估的结构性变化：第一，鸿海精密（Hon Hai）的份额正在从2025年的51\%下降至2026年的41\%，而纬创（Wistron）和广达（Quanta）的份额相对稳定；第二，纬创的营收在5月环比增长2\%，主要来自其子公司纬颖（Wiwynn）的贡献，而非传统PC业务；第三，摩根斯坦利对三家ODM的偏好排序是“纬创 > 鸿海 > 广达”，这一排序本身就在告诉市场：在AI服务器供应链中，真正的价值创造正在从“规模”转向“议价权”。

1. 5月出货量环比下降是正常的季度节奏调整，不应被误读为需求拐点

5月GB200/300机柜出货量7700台，较4月的8300台下降7\%。广达从约2100台降至 台，鸿海从约3700台降至3300台，纬创则基本持平在 台。摩根斯坦利明确指出，广达的2Q整体出货量预计仍将环比增长40\%至约6700台，下降的部分只是被推迟到了下半年。

这意味着什么？AI服务器供应链的季度内节奏正在变得更有规律：1Q是爬坡期，2Q加速，3Q和4Q是出货高峰。2025年的数据也印证了这一模式——从1Q的900台到4Q的15500台，季度环比增速从未低于80\%。2026年的节奏虽然更平滑，但全年超过100\%的同比增长预期并未改变。

对于产业决策者而言，月度数据的意义不在于判断趋势方向，而在于验证供应链的执行力。 5月的下降更多反映了ODM的产能调配和客户提货节奏，而非终端需求的减弱。那些因为单月数据就调整全年预期的投资者，很可能低估了AI基础设施投资的惯性。

2. 鸿海份额下降不是竞争失利，而是供应链结构正在从“一家独大”走向“多极分化”

2025年，鸿海占据了GB200/300机柜出货量的51\%，广达21\%，纬创20\%。到了2026年，摩根斯坦利预计鸿海的份额将降至41\%，广达升至24\%，纬创维持在18\%左右，而“其他”供应商的份额将从6\%跃升至17\%。

这不是鸿海的失败，而是供应链的必然演化。当一项技术从早期量产进入规模化交付阶段，单一供应商的产能瓶颈会自然推动客户分散订单。鸿海的绝对出货量仍在增长——从2025年的约4.5万台增至2026年的约5.1万台——但增速慢于行业整体。

*更值得关注的是“其他”供应商的崛起。** 17\%的份额意味着约2.1万台的出货量，这部分产能正在被二线ODM和新兴服务器制造商瓜分。对于投资者而言，这意味着AI服务器供应链的投资机会正在从“押注龙头”转向“寻找具备差异化能力的二线玩家”。那些能够在散热、高速互连或系统集成上建立壁垒的公司，即使规模较小，也可能获得更高的估值溢价。

3. 纬创的营收结构正在发生质变：AI服务器正在取代PC成为增长引擎

5月，纬创营收2900亿新台币，环比增长2\%，同比增长39\%。在传统PC业务下滑的背景下——笔记本出货量环比下降6\%至170万台，台式机下降25\%至60万台——增长主要来自纬颖的营收增长2\%和AI服务器计算托盘出货量的小幅提升。

*这是一个重要的结构信号。** 纬创正在完成从“PC ODM”到“AI服务器核心供应商”的转型。摩根斯坦利维持其2Q机柜出货量预估约4200台，环比增长4\%，这意味着6月需要出货约1500台。虽然环比增速不高，但考虑到纬创的AI业务占比正在快速提升

[... middle omitted ...]

报告中承认了一个重要细节：“实际交付给终端客户的机柜数量可能低于这些数字，因为我们包含了纬创的计算托盘（L10）机柜当量，但没有考虑L11的机柜组装和测试时间。”

这意味着，ODM的出货量并不等于最终客户的安装量。 机柜从出厂到真正上线运行，中间还有运输、上架、网络配置、软件部署等多个环节。如果终端客户的消化速度跟不上ODM的出货节奏，供应链中就会出现“在途库存”或“渠道库存”的积累。

这个问题在2025年并不突出，因为当时的需求远大于供给。但到了2026年，随着供给端产能大幅释放，供需关系正在趋于平衡。如果终端客户的AI基础设施部署计划出现任何延迟——无论是由于电力、数据中心空间还是软件优化——ODM的出货数据就可能出现虚高。

*这是投资者需要密切跟踪的风险点。** 摩根斯坦利的全年70-80K出货量预测建立在终端需求持续强劲的假设之上，但报告并未提供终端客户的实际装机数据来验证这一假设。当前市场可能低估了“渠道库存调整”对2026年下半年出货节奏的潜在影响。

\subsection{[R025] NOM：人民币中间价模型正在发出一个被市场忽视的稳定信号}
{\small\color{refgray}归类：区域市场：新兴市场融资货币单一化，亚洲资金撤退信号值得警惕}

人民币汇率市场当前最值得关注的，或许不是每日的即期波动，而是那个看似技术性的、由模型驱动的每日中间价设定机制。NOM亚洲外汇策略团队近期发布的一份模型追踪报告，揭示了一个关键判断：基于当前模型估算，人民币兑美元中间价的预测值已出现显著下移，这表明定价机制本身正在释放一种被市场低估的稳定性信号。

这份报告的核心价值不在于提供一个具体的点位预测，而在于它提供了一套拆解中间价形成机制的量化框架。对于关注中国资产定价逻辑的投资者而言，理解这个框架，比追逐任何单一汇率目标都更重要。

我们需要回答的问题是：这个模型下移意味着什么？它背后的驱动力是什么？以及，这个信号对未来的市场博弈格局有什么含义？

NOM模型的最新预测显示，剔除逆周期因子后的中间价预测值为6.7775，较前次预测的6.8147大幅下移了372个基点。这是一个不容忽视的幅度。更值得注意的是，即便计入逆周期因子，模型预测值也达到6.7962，较前次调降了185个基点。

这个数字本身不是预言，而是一个信号。它表明，在NOM模型的框架下，推动中间价变动的核心变量组合——主要是隔夜一篮子货币的变动——正在指向一个更低的人民币定价中枢。这意味着，如果市场条件不发生剧烈变化，官方中间价可能在未来一段时间内，以更温和的方式向这个方向靠拢。

对于市场参与者而言，这相当于一个“隐性指引”。中间价设定机制本身，正在扮演一个比即期汇率更稳定的锚。投资者不应该只盯着即期市场每天几百点的波动，而应该关注这个锚的变动方向。

2. 驱动模型下移的四大货币：AUD、KRW、EUR、JPY的贡献并非随机

NOM报告明确指出，对模型预测变动贡献最大的四个货币依次是澳元、韩元、欧元和日元。这并非一个随机的排序，它揭示了人民币中间价定价机制中，全球风险偏好与贸易加权竞争格局的双重影响。

澳元和韩元的高贡献，反映了全球增长预期和风险偏好的变化。这两个货币对全球经济周期和贸易活动高度敏感，它们的走强往往意味着市场对全球需求的信心恢复。欧元和日元的贡献，则更多地体现了主要经济体货币政策差异和利差交易的影响。

把这些事实合在一起，结论是清晰的：人民币中间价模型的变动，并非孤立地反映中美双边关系，而是系统地捕捉了全球汇率市场的整体走势。 这意味着，任何试图单边押注人民币贬值的策略，都可能忽略了这个多元化的定价基础。人民币的稳定，正在被一个更广泛的全球货币篮子所支撑。

NOM报告提供的另一个关键图表，是模型误差的历史走势。数据显示，模型误差（即实际中间价与模型预测值之间的偏差）在2025年初曾达到-1800个基点的高位，随后迅速收窄，并在2026年初转为正值且维持在600个基点左右的水平。

这个模式传递的信息非常明确：中间价定价机制正在从过去两年的剧烈偏离状态，回归到一个更可预测、更透明的轨道。 误差的收敛意味着，官方在设定中间价时，正在更紧密地跟随模型所反映的市场力量，而非频繁地使用逆周期因子进行大幅干预。

对于投资者而言，这意味着中间价的可预测性正在提高。一个更可预测的定价机制，本身就是一种稳定性的来源。它降低了汇率政策的不确定性，为企业的汇率风险管理、以及境外投资者的资产定价，提供了一个更清晰的参考系。

4. 报告尚未完全回答的关键问题：逆周期因子的“隐形”干预边界

NOM模型虽然提供了清晰的量化框架，但它并未完全解答一个关键问题：当模型预测值与政策意图发生冲突时，逆周期因子的使用边界在哪里？

报告给出了计入逆周期因子后的预测值，但并未深入分析逆周期因子在什么条件下会被激活，以及其激活的力度如何变化。这是一个典型的“黑箱”问题。从历史数据看，逆周期因子曾在市场单边贬值预期强烈时被显著使用，以平滑中间价的波动。但在当前模型预测值本身已经下移的背景下，逆周期因子的角色可能正在从

[... middle omitted ...]

变化。

第一，关注模型误差的收窄或扩大。如果误差持续收窄并维持在低位，说明定价机制正在回归市场化，中间价的信号意义更强。如果误差再次显著扩大，则意味着政策干预力度加大，市场需要重新评估政策意图。

第二，关注AUD、KRW、EUR、JPY的权重变化。这些货币的权重并非固定不变，其变化反映了全球汇率市场的主要驱动力。例如，如果澳元和韩元的贡献权重上升，说明全球风险偏好是主导；如果欧元和日元的权重上升，则说明利差和货币政策差异是主导。理解这个权重的变化，有助于判断人民币汇率变动的根本原因，而非仅仅停留在对中美关系的讨论上。

这个观察框架的价值在于，它让投资者从“猜测官方意图”转向“理解市场与政策的互动”，从而做出更具前瞻性的判断。

6. 2026年下半年的关键事件：APEC会议与中美高层互动

NOM报告在附录中列举了2026年下半年的关键事件日历，其中两个事件尤为值得关注：11月在深圳举行的APEC峰会，以及年底美国总统特朗普表示将邀请中国国家主席习近平访问白宫。

\subsection{[R026] BARC：市场低估了供给侧的持续收紧，而非单纯押注地缘风险}
{\small\color{refgray}归类：宏观与利率：通胀但不加息，市场定价逻辑的重构 / 能源与大宗：供给硬约束重塑定价权，中国出口成金属市场泄压阀}

这份BARC研报最值得关注的核心判断，并非它预测2026年布伦特原油均价为100美元/桶——这个数字本身在市场中已有讨论——而是它揭示了一个更深层的结构性矛盾：即便假设霍尔木兹海峡在本月底恢复通行，2027年的油价仍将显著高于远期曲线和共识预期。报告给出的2027年布伦特均价预测为88美元/桶，高出远期曲线7美元/桶，高出共识13美元/桶。这意味着，BARC认为市场正在系统性地低估一个超出短期地缘冲突范畴的供给侧再定价过程。

当前市场的主流叙事是“地缘溢价正在消退”。管理基金在原油期货和期权上的净投机仓位已回落至战前水平，隐含波动率也同步下降。市场似乎在告诉投资者：最坏的时刻已经过去，基本面正在消化冲击。但BARC的数据指向了一个相反的结论：基本面不仅没有消化冲击，反而在被动地“长入”最初的恐慌——库存以每天450万桶的速度下降，美国商业原油库存（经管道填充和战略储备释放调整后）已逼近2008年金融危机的低点，库欣储油设施的利用率降至2010年以来最低。

这份报告的价值不在于预测一个精确的油价数字，而在于它拆解了一个关键问题：为什么市场与基本面之间出现了如此大的认知裂口？答案指向一个被广泛忽视的变量——中国需求的“战略性沉默”。

1. 中国需求的骤降并非价格弹性，而可能是一种审慎的战略储备行为

报告中最具颠覆性的观察来自中国。4月份中国石油需求同比下降12\%，市场普遍将其解读为“高油价下的需求破坏”——即价格弹性发挥了作用。但BARC提出了一个截然不同的假设：这种需求下降可能并非真实的消费萎缩，而是中国在主动抑制采购、暗中储备。

支持这一判断的证据来自库存数据。自伊朗战争爆发以来，中国的原油显性库存几乎没有变化。而同期全球其他地区的库存均在快速下降。如果真是高油价压垮了需求，库存应当同步消耗，而非纹丝不动。BARC指出，中国过去几年在电动车渗透率持续上升的背景下，石油需求仍保持稳步增长，因此用“能效突然提升”来解释4月12\%的同比降幅，逻辑上并不顺畅。

更合理的解释是：中国正在为霍尔木兹海峡长期封锁的极端情景做准备。如果这一假设成立，那么当前市场看到的“需求疲软”实际上是伪装的供给储备。一旦地缘局势明朗或中国认为储备目标达成，这部分被压抑的需求将集中释放，届时全球供需平衡表将面临比当前更剧烈的冲击。

这个判断如果正确，将对所有基于“需求见顶”叙事的长期油价预测构成挑战。市场可能正在用稳态条件下的需求弹性模型，去评估一个非稳态的战略行为。

BARC对美国库存的分析是这份报告另一处关键的技术贡献。表面上看，美国商业原油库存约为4.5亿桶，远高于2008年底的3.01亿桶。但BARC做了一个关键的“还原调整”：自2008年以来，战略石油储备（SPR）释放了约7亿桶，管道填充增加了约7.5亿桶，而有机的供需基本面因素实际上消耗了约4.5亿桶。调整后，美国商业原油库存的真实水平已经低于2008年。

更紧迫的信号来自库欣。作为WTI的定价枢纽，库欣的储油设施利用率已降至35\%，这是2010年以来的最低水平。如果当前库存消耗速度持续，美国商业原油库存（以炼油需求天数衡量）将在8到10周内跌破20年最低水平。

这意味着，即便霍尔木兹海峡问题得到解决，美国本土的库存缓冲也已极度脆弱。任何额外的供给中断或需求回升，都可能引发WTI的剧烈重新定价。BARC预测WTI在2026年Q2均价为102美元/桶，高出远期曲线5美元，但到了Q3和Q4，溢价分别扩大到12美元和14美元。这种溢价结构本身就在告诉市场：真正的紧张不在当下，而在未来。

3. 2027年的供给过剩叙事可能被高估，非OPEC增产无法填补缺口

市场普遍认为，一旦霍尔木兹海峡恢复通行，2027年将出现大规模供给过剩。BARC自己的平衡表也显示，基于本月

[... middle omitted ...]

本开支，仍是一个巨大的问号。BARC没有在报告中直接回答这个问题，但暗示了这种结构性变化可能使2027年的供给曲线比市场预期的更陡峭。

4. 报告尚未完全回答的关键问题：闲置产能的真实规模与释放意愿

尽管BARC的分析框架非常严谨，但有一个关键变量仍然悬而未决：OPEC+的闲置产能到底有多少，以及这些产能能否在需要时迅速释放。

报告提到，即使阿联酋完全动用闲置产能，也无法填补2027年的供需缺口。但这里隐含了一个假设：阿联酋愿意且能够将闲置产能全部转化为实际产量。这一假设在地缘政治格局下可能面临挑战。如果伊朗局势导致整个海湾地区的安全环境恶化，阿联酋的产能释放意愿和实际能力都可能打折扣。

此外，沙特阿拉伯的产能策略也值得推敲。过去几年，沙特一直在强调其“产能管理”而非“市场份额最大化”的立场。在油价突破100美元/桶的情况下，沙特是否会选择大规模增产来压制价格，还是维持高油价以支撑其财政支出和“2030愿景”投资？这个问题的答案将直接影响2027年的供需平衡。

\subsection{[R027] GS：澳大利亚消费者信心连续三个月低迷，但市场真正该关注的不是情绪本身}
{\small\color{refgray}归类：地产与消费：成交量回升但结构性分化，二手房成更可靠温度计}

GS：澳大利亚消费者信心连续三个月低迷，但市场真正该关注的不是情绪本身

澳大利亚消费者信心指数在6月进一步下滑至80.6，较历史均值低约20\%。这已经是连续第三个月处于极弱水平。GS最新研报披露了这一数据，但真正值得投资者警觉的，并非情绪本身有多差，而是情绪低迷正在从“可被忽略的噪音”转变为“可能自我实现的经济拖累”。

过去两年，澳大利亚央行一直认为消费者信心与消费行为之间的关联有限。但GS在报告中引用央行的最新表述时留下了一个关键的伏笔：当情绪低迷持续足够久，且伴随资产价格下跌时，它可能开始产生实质性的二阶效应。6月的数据恰好提供了这种转变的初步证据——家庭财务状况感知全面恶化，住房价格预期三年来首次跌破长期均值，悉尼和墨尔本的跌幅尤其显著。

这份报告的价值不在于告诉你“消费者不开心”，而在于它揭示了一个正在形成的反馈循环：信心疲弱正在抑制住房市场，住房市场走弱反过来进一步打击信心。这个循环是否会被打破，取决于接下来的劳动力市场走势和央行的政策路径。

整体信心指数下降2.9\%固然值得关注，但更值得细看的是分项数据。GS报告显示，家庭对当前财务状况的感知环比下降7.5\%，对未来一年的财务预期更是暴跌8.5\%。这两个分项同时恶化，意味着消费者不仅在抱怨当下，更在提前调低对未来的预期。

这里有一个重要的结构性差异：长期经济前景预期虽然也下降了3.2\%，但一年期经济预期反而回升了4.9\%。这种“短期经济预期改善、家庭财务预期恶化”的分化，实际上指向了一个更具体的担忧——消费者并不认为宏观经济会崩溃，但他们认为自己的钱包正在被挤压。换句话说，这是一种“我比经济更难过”的微观体感，这种体感一旦形成，往往比宏观悲观更难逆转。

对于关注澳大利亚消费类资产和零售行业的投资者而言，这意味着：即使GDP数据没有大幅走弱，消费行为也可能提前收紧。家庭财务感知是预测消费倾向的领先指标，而非滞后指标。

GS报告中最引人注目的数据点之一是：澳大利亚居民房价预期环比暴跌14.9\%，三年来首次低于长期均值。分州来看，悉尼和墨尔本的跌幅最为惨烈，分别达到19\%和18\%，而这两个城市恰好也是实际房价已经在下行的地区。

这里需要厘清一个因果关系：是实际房价下跌导致了预期恶化，还是预期恶化加速了实际房价下跌？GS没有给出明确的归因，但报告提供了另一个值得玩味的数据——“购房时机”指数在经历了5月的急剧下滑后，6月反弹了12.6\%。这个反弹与房价预期的大幅下跌形成鲜明对比，意味着部分潜在购房者可能认为价格回调创造了入场机会。

这种“想买但不敢买”的矛盾心理，恰恰是市场处于转折点的典型特征。对于房地产开发商、银行按揭部门和REIT投资者来说，这意味着：短期内交易量可能继续低迷，但价格调整速度可能快于市场共识。如果价格下跌足够快，反而可能提前触发行情企稳。

GS在报告中特意引用了澳大利亚央行的一个关键判断：过去的研究发现，情绪对消费的独立驱动作用有限，但“如果情绪的大幅变化持续存在，可能会产生更显著的影响”。这句话的潜台词值得反复咀嚼。

过去两年，央行一直用“情绪不等于行为”来安抚市场，认为消费者信心低迷不会自动转化为消费收缩。但连续三个月的极弱数据，加上住房市场信心的系统性恶化，正在挑战这一假设。GS报告虽然没有明确质疑央行的框架，但通过引用央行自身的条件性表述，实际上在提醒读者：这个假设的适用边界可能正在接近。

对于利率市场参与者来说，这意味着：如果后续消费数据出现超预期疲软，央行可能比市场预期的更快调整政策立场。当前市场定价中隐含的降息路径可能低估了这种风险。

4. 报告没有完全回答的关键问题：劳动力市场何时成为下一个多米诺骨牌

GS这份报告对消费者信心和住房市场的分析相当透彻，但它留下了一个尚未回答的关键问题：劳动力市场

[... middle omitted ...]

石。

基于GS这份报告的逻辑，投资者不应简单地将消费者信心指数视为一个独立的交易信号，而应建立一套从情绪到行为的跟踪框架：

第一，关注家庭财务感知分项的持续性。如果连续两个月以上维持在恶化区间，消费数据出现超预期下行的概率将显著上升。

第二，跟踪住房价格预期与实际房价的收敛速度。当前预期远低于实际价格，如果实际价格开始加速追赶预期，意味着调整正在深化；如果预期开始修复，则可能意味着市场底部正在形成。

第三，将悉尼和墨尔本的信心数据作为领先指标。这两个城市在房价预期跌幅上领先全国，它们的信心修复节奏将是判断全国趋势的重要参考。

第四，警惕“情绪-消费-就业”的负反馈循环。一旦失业预期开始上升，目前的信心低迷就可能从“可容忍的疲弱”升级为“需要政策回应”的系统性风险。

GS的这份报告提供了一个清晰的信号：澳大利亚消费者信心的恶化已经不再是简单的周期性波动，而是正在形成结构性的负面反馈。但真正的拐点何时到来，取决于劳动力市场能否继续吸收来自住房和消费领域的压力。

\subsection{[R028] JPM：中国出口正在成为全球金属市场结构性短缺的“泄压阀”}
{\small\color{refgray}归类：能源与大宗：供给硬约束重塑定价权，中国出口成金属市场泄压阀}

全球基本金属市场正在经历一个罕见的错位时刻。JPM最新发布的《基本金属供需追踪报告》揭示了一个被广泛讨论但未被充分定价的格局：中国国内需求正在放缓，但中国生产商却在高价区间继续扩产；与此同时，除中国以外的市场库存正在快速消耗，而中国自身的库存却处于多年高位。这两组矛盾现象的背后，隐藏着一个关键的结构性变量——中国出口正在成为全球金属供应链的“泄压阀”。

这份报告的核心判断并不复杂，但其对资产定价和产业策略的影响远未被市场完全消化。中国不是全球需求的“发动机熄火”，而是正在从需求方转变为供给的“调节器”。理解这一角色转换，是理解未来12个月基本金属价格走势的前提。

1. 中国铜消费的真实图景远比表面数字复杂，终端需求放缓与表观消费高企的背离蕴含关键信号

JPM的数据显示了一个令人困惑的现象：2026年4月，中国的铜表观消费同比增长了9\%，但该机构自建的终端消费加权指标却同比下降了4\%。即便将年初至今的数据拉平，终端需求也仅微增1.2\%。这种背离并非统计噪声，而是反映了供应链行为的深刻变化。

表观消费的高企主要来自两个因素。其一，2025年四季度至2026年一季度中国买家“罢工”期间积累的库存正在被消化，去库存行为在4月集中体现为表观消费的虚高。其二，电网投资和新能源汽车生产仍在支撑部分需求，但这不足以抵消可再生能源装机骤降和白电产量下滑带来的拖累。报告特别指出，4月太阳能装机同比暴跌79\%，风力装机下降7\%，白色家电中的空调产量同比下滑9\%。

这意味着什么？中国铜的真实需求并没有价格所反映的那么强劲。当前铜价在每吨13500至14000美元区间震荡，很大程度上是库存再分配和供给端故事在驱动，而非终端消费的全面复苏。对于投资者而言，需要警惕的是：一旦中国去库存周期结束，表观消费将迅速向真实需求收敛，届时铜价将面临一个缺乏基本面支撑的脆弱窗口。

2. 美国232条款铜关税的走向是未来6个月铜市场最大的政策变量，其影响可能被低估

6月30日，美国商务部长将向总统提交关于是否需要进一步修改232条款铜关税的评估报告。JPM的判断是，特朗普政府更有可能实施2025年7月行政命令中设定的阶梯式关税方案——2027年起征15\%，一年后升至30\%。其逻辑并非简单的贸易保护，而是战略意图：确保已经流入美国的铜留在境内，作为关键战略储备，而非鼓励其在未来几年快速去库存。

如果这一判断成立，未来几个季度美国将持续保持强劲的进口吸引力，这将持续抽紧美国以外的阴极铜供给，尤其是在中国进口窗口同时打开的时间段。这里的二阶效应值得深思：美国加征关税的初衷是保障自身供应安全，但实际效果可能是将全球铜市场的紧张程度推升至一个难以持续的水平，最终反噬美国本土的加工制造业成本。

报告没有完全展开的一个问题是：如果美国关税政策落地，COMEX与LME之间的价差将如何演变？价差的扩大是否会反过来刺激更多套利行为，从而进一步扭曲全球铜的贸易流向？这些问题的答案，将直接决定未来铜价曲线的形态。

3. 铝市场的核心叙事不是中国需求，而是中国出口如何为全球去库存提供缓冲

LME铝库存已降至33万吨的极低水平，而中国境内的显性库存仍接近140万吨的多年来高位。这种“冰火两重天”的格局，正在通过中国半成品和成品铝材的出口来弥合。JPM指出，中国对亚洲其他地区的半成品出口套利窗口已经打开，且处于有吸引力的水平。

然而，一个值得注意的数据点是：在截至5月中旬的过去6周内，中国上海期货交易所和区域仓库的铝库存仅下降了约9万吨。去库存的节奏并不快，原因在于中国产量仍在增长——4月精炼铝产量同比增长3.4\%，年化产量达到4500万吨，接近产能上限。与此同时，铝土矿进口依然坚挺。

这里有一个尚未被充分讨论的风险点：几内亚铝土矿出口的稳定性。中国75\%的

[... middle omitted ...]

近10\%，现货加工费自3月底以来已转为负值。这意味着矿山正在向冶炼厂收取加工费，而不是支付加工费——这在历史上是矿端极度紧张的信号。然而，冶炼厂的生产却依然强劲，4月产量同比增长5\%，因为它们在消耗一季度大量进口积累的库存。

问题的另一面是需求。中国的镀锌厂开工率持续低迷，建筑活动毫无起色，总显性库存稳定在26.4万吨附近。截至4月，表观需求同比下降12\%。这种“供给端收紧但需求端疲软”的格局，使得锌市场处于一种脆弱的平衡中。一旦冶炼厂库存消耗完毕，而需求仍无改善，锌价可能面临两种截然不同的路径：要么矿端紧张推动价格上行，要么需求崩溃导致价格暴跌。报告倾向于认为前者的概率更大，但并未给出明确的概率判断。

5. 镍市场的库存故事已经讲完，但印尼政策框架正在改写成本曲线

镍的显性库存（LME+SHFE）年初至今净增7万吨，总量约47万吨。虽然LME库存在2月见顶后小幅回落，但整体库存仍在积累。市场关注的焦点已经从“库存到底有多高”转向“印尼政策将如何改变供给格局”。

\subsection{[R029] JPM：美国军工复合体的真正危机不在于预算，而在于工业基础的结构性过时}
{\small\color{refgray}归类：风险偏好：CLO市场从beta转向alpha，军工复合体面临工业基础结构性过时 / 黄金与矿业：黄金股回调是表象，矿业主权供需重塑竞争格局}

JPM：美国军工复合体的真正危机不在于预算，而在于工业基础的结构性过时

美国仍然是全球最强大的军事力量。但JPM一份最新发布的深度研报提出了一个让产业决策者无法回避的问题：支撑这种力量的工业基础，是否已经跟不上它所面对的战争形态？

这份报告的核心判断是：美国国防工业基础（DIB）当前面临的挑战，不是钱不够，而是结构不对。过去三十年的整合与集中化，让美国军工体系走向了“精致但脆弱”的路径——以F-35为代表的高端旗舰平台在技术上无可匹敌，但开发周期长、成本高昂、难以规模化。而今天的战场，从乌克兰到中东，再到潜在的台海冲突，正在同时要求三样东西：可消耗的硬件、可迭代的软件、以及能够快速整合两者的工业能力。

这不是一份关于预算拨款的报告。这是一份关于工业逻辑如何被战争形态倒逼重写的分析框架。对于关注军工产业链、国防科技创业、以及中美竞争格局的读者，这份报告的参考价值不仅在于它梳理了历史，更在于它提出了一个尚未被市场充分定价的命题：美国军工股的估值逻辑，可能正站在一个结构性拐点上。

JPM研报追溯了美国DIB的演变轨迹，其核心叙事线非常清晰：从二战时期的全民动员，到冷战时期的专业化与垂直整合，再到冷战结束后的剧烈整合，每一次转变都在提升技术深度，但也在牺牲生产弹性。

二战期间，美国动员了超过80万家民用企业参与国防生产，飞机年产量从1939年的约6000架飙升至1944年的约23.7万架。那是一个“规模即优势”的时代。冷战时期，面对苏联的明确威胁，美国转向了专业化承包商体系，技术深度大幅提升，但供应商基础开始收窄。1993年，时任副国防部长威廉·佩里召集了著名的“最后的晚餐”会议，敦促军工企业合并。到2000年，主要承包商从超过50家锐减至5家——洛克希德·马丁、波音、诺斯罗普·格鲁曼、雷神、通用动力。这五家公司控制了超过60\%的合同价值。

这个结构的逻辑是：在一个超级大国对抗明确对手的时代，技术优势可以通过旗舰平台来锁定胜局。F-35、福特级航母、弗吉尼亚级潜艇——这些项目代表了美国军事技术的最前沿，也代表了“精品化”路线的极致。

但问题在于，当威胁从单一超级大国分散为多个层次——近邻竞争对手、中等强国、非国家行为体——时，这套体系暴露了四个相互关联的脆弱点：生产深度不足、采购周期过长、行业集中度过高、供应链韧性脆弱。研报特别指出，这些脆弱点不是理论上的，而是已经在现实中被验证的：乌克兰战争中弹药消耗的速度远超美国当前的产能，而中国在造船和导弹制造上的规模优势，正在从根本上改变“工业能力即战斗力”的方程。

2. 乌克兰和加沙的启示：可消耗系统与软件迭代正在改写“杀伤链”

研报在讨论现代战争形态变化时，给出了一个值得反复咀嚼的判断：“技术优势本身已不再足够；军事优势越来越取决于工业产能、成本纪律，以及能力部署的速度和适应性。”

这不是否定高端平台的价值。相反，报告明确强调，可消耗的无人机系统不应替代F-35这样的“精品”平台，而应与之互补。真正需要改变的是两者之间的整合方式。现代战争的杀伤链——从发现、定位、跟踪、瞄准、打击到评估——正在被软件和数据的实时整合彻底重塑。谁能在更短的时间内完成这个闭环，谁就掌握了战场主动权。

这一点在乌克兰战场上表现得尤为明显。低成本无人机与实时指挥控制架构的结合，让传统上需要昂贵平台完成的任务变得可以规模化执行。伊朗在加沙冲突中展示的无人机和导弹战术，同样指向了这个方向。JPM研报的洞察在于：这些冲突暴露的不仅是某种武器系统的有效性，更是一种新的工业逻辑——硬件必须可规模化、软件必须可迭代、系统必须可整合。

对于投资者而言，这意味着传统军工企业的护城河正在被侵蚀。过去，进入壁垒是技术复杂性和政府关系。现在，新的进入者——尤其是那些以软件为核心、以商业供应链为依托

[... middle omitted ...]

系。中国的造船能力和导弹产量已经远超美国，而在稀土、半导体等关键材料和元器件上的供应链控制力，更是直接渗透到了美国国防项目之中。JPM指出，中国的A2AD（反介入/区域拒止）战略不仅是一个军事概念，它背后是一整套工业能力的支撑：能够以低成本、大规模、快速迭代的方式生产武器系统。

这不是美国当前DIB所能轻易复制的。美国的问题在于，它的工业基础已经被优化为了“精品化”路线——擅长生产少量、高价、高性能的系统，但不擅长生产大量、低价、可消耗的系统。而后者，恰恰是当前和未来冲突中最需要的。

研报引用了其他国家的经验——韩国、瑞典、以色列——这些国家在国防工业整合上各有特色，但共同点是政府与产业之间的紧密协同、供应链的韧性设计、以及快速适应战场反馈的能力。这些经验对美国而言，既是参照，也是挑战。

JPM研报在分析政策推动时，给出了一个冷静的判断：“政策动力已经在过去几年启动，并建立在国防科技新公司快速崛起的基础上。但政策推动仍然面临着几个执行障碍，制度性障碍依然存在。”

\subsection{[R030] 摩根斯坦利：中国尿素出口底价的重新设定，是全球化肥定价逻辑的一次结构性重置}
{\small\color{refgray}归类：能源与大宗：供给硬约束重塑定价权，中国出口成金属市场泄压阀}

摩根斯坦利：中国尿素出口底价的重新设定，是全球化肥定价逻辑的一次结构性重置

全球化肥市场正在经历一个被多数投资者低估的定价机制转变。摩根斯坦利最新发布的这份报告揭示了一个关键信号：中国刚刚重新设定了尿素出口价格下限——对印度出口的大颗粒尿素不低于每吨500美元FOB，车用级不低于每吨510美元FOB。这一数字本身并不令人震惊，真正值得关注的是它背后的逻辑：中国正在用价格工具替代配额工具，从“限制出口数量”转向“锁定出口价格”。

这意味着什么？意味着全球尿素市场的价格锚点正在被重新定义。过去几年，市场习惯了用中国出口配额作为供给侧的调节器，但这份报告暗示，配额逻辑正在被价格逻辑取代。而这一转变的背景，是霍尔木兹海峡持续关闭、伊朗供给中断、以及沙特SABIC正在加速的产能转移。

这不是一个短期事件。它可能重塑未来12-18个月全球化肥贸易的竞争格局。

1. 价格下限不是简单的限价，而是中国化肥出口政策的范式转换

摩根斯坦利的分析显示，中国此次设定的价格下限与此前每吨660-680美元FOB的初始水平相比，已经显著下调。但更关键的是政策工具本身的变化。

过去几年，中国对尿素及DAP/MAP出口采取的是“配额+窗口期”管理。出口商需要在有限的配额内竞争，结果往往是价格被压低，利润被压缩。而现在，价格下限的引入意味着：只要出口价格不低于底线，出口商可以自由交易。这本质上是一种从“量控”到“价控”的转变。

报告提到，低于下限出口的出口商可能失去“自律性企业”资格。这不是一个软性约束，而是一个有实际惩罚机制的硬性门槛。对于中国化肥企业而言，这意味着它们的出口策略必须从“抢配额、拼价格”转向“保价格、赚利润”。

这个转变的含义是：中国不再追求全球市场份额的最大化，而是追求出口利润的最大化。对于全球买家而言，这意味着来自中国的廉价尿素供给将显著减少。

报告中最具操作意义的细节是：最新一轮印度尿素招标规模约170万吨，收到的报价低至每吨410美元CFR，折算回中国FOB价格不足400美元。而中国设定的下限是500美元FOB。摩根斯坦利的判断很直接：中国将不会参与这轮印度招标。

这是一个明确的信号：中国愿意放弃印度市场，以维护自己的价格底线。印度是全球最大的尿素进口国之一，这种“主动缺席”的行为，将直接推高其他供给方的议价能力。

报告没有展开的是，印度招标的低价本身反映了全球尿素市场近期的显著下跌：美国Nola价格已跌至每吨360-410美元区间，巴西跌至530-540美元CFR，伊朗更是暴跌100美元至620-630美元FOB。中国在这个时点设置价格下限，本质上是在全球价格下行周期中“筑底”。

这意味着什么？如果中国持续缺席印度招标，全球尿素价格的下行空间将被压缩。对于持有化肥资产的投资者而言，这是一个需要重新评估风险收益比的变化。

3. SABIC的产能转移正在改变中东供给格局，但市场尚未充分定价

报告另一个值得关注的信号来自沙特。摩根斯坦利指出，SABIC在5月首次实现了尿素散货出口，每日产量的40\%已从朱拜勒转移至延布。管理层目标是在未来通过替代路线转移70\%的日产量，并且所有工厂都在满负荷运转。

这个细节的重要性在于：它发生在霍尔木兹海峡持续关闭的背景下。SABIC的产能转移不是临时应对，而是结构性调整。如果SABIC能够通过红海港口维持甚至扩大出口，那么中东的尿素供给弹性将显著增强。

但报告也留下了一个未解问题：SABIC的产能转移能否完全弥补伊朗的供给缺口？伊朗的尿素价格暴跌100美元，说明其出口可能已经受到实质性影响。如果SABIC的增量不足以覆盖伊朗的减量，全球尿素供给仍将偏紧。

对于市场参与者而言，需要跟踪的关键变量是：SABIC的转移进度、伊朗出口的实际恢复时间、以及霍尔木兹海

[... middle omitted ...]

体水平仍在讨论中。这意味着价格下限可能不是一个单一数字，而是一个区间。如果中国出口商通过调整产品规格或目的地来规避限制，价格下限的实际约束力就会打折扣。

第二，DAP/MAP的出口恢复时间。报告提到，中国DAP/MAP出口可能在8月恢复，但摩根斯坦利对此持谨慎态度。原因是硫磺市场持续混乱——不仅仅是霍尔木兹海峡问题，俄罗斯近期也在减少出口。硫磺是生产磷肥的关键原料，如果硫磺供给无法改善，DAP/MAP的出口恢复可能推迟。

这两个问题直接关系到全球化肥市场的供给节奏和价格走势。投资者需要密切跟踪中国政策执行层面和硫磺市场的动态。

5. 投资者的观察框架：从“配额博弈”转向“价格锚定”和“替代供给”

过去，跟踪中国化肥出口的核心变量是配额总量和发放节奏。现在，核心变量变成了：价格下限水平、执行力度、以及印度等主要买家的反应。

同时，替代供给的弹性变得更加重要。SABIC的产能转移、伊朗出口的恢复速度、以及俄罗斯硫磺出口的变化，都将共同决定全球化肥市场的供需平衡。

\subsection{[R031] 摩根斯坦利：中国取消尿素出口底价，全球化肥市场的定价权正发生结构性转移}
{\small\color{refgray}归类：能源与大宗：供给硬约束重塑定价权，中国出口成金属市场泄压阀 / 化工与材料：合成橡胶供给缺口重塑定价权，松下投资者日揭示AI电源机遇}

摩根斯坦利：中国取消尿素出口底价，全球化肥市场的定价权正发生结构性转移

中国在5月27日刚为尿素出口设置了价格底线，不到两周就悄然撤除。这个动作本身并不意外，因为全球尿素价格已经跌到远低于中国设定的每吨660-680美元FOB水平，继续维持底价在商业上没有意义。但摩根斯坦利这份报告真正值得关注的判断，不在于“中国调整了政策”，而在于这个调整透露了中国化肥出口战略的深层逻辑转变——中国正在从“保国内供应、限出口”的防御姿态，转向“在保障国内低价的同时，最大化出口现金流”的平衡策略。这一转变的隐含含义，可能被全球农产品和化肥市场严重低估。

报告的核心洞察只有一个：中国化肥出口政策不再是简单的开关式管控，而是进入了精细化、选择性、与地缘风险高度耦合的新阶段。对于全球氮肥和磷肥市场而言，这意味着中国不再只是一个被动的价格接受者，而是正在成为供给侧的主动定价变量。

市场此前对中国设置尿素出口底价存在两种解读。一种认为这是中国在限制出口总量的同时，为国内生产商争取更高利润；另一种则认为这是中国在地缘不确定性下，战略性保留国内库存的手段。摩根斯坦利通过“底价取消”这个逆向验证，给出了清晰的判断：中国优先级的核心不是利润最大化，而是出口量的释放。

底价设置的初衷，是在3百万吨出口配额内，确保中国生产商不会因低价亏损而减少出口。但当全球价格跌至底价以下后，中国没有选择“等价格反弹再出口”，而是直接取消底价，允许出口在更低价格水平上继续。这个选择说明，对中国而言，维持出口通道的畅通、消化国内过剩产能，比短期价格谈判权更重要。

这一判断对中国化肥行业的竞争格局有直接含义。如果中国愿意在低于全球成本曲线的价格水平上出口，那么全球尿素市场的边际定价权将从高成本生产商（如欧洲的天然气路线）转向中国。中国生产商虽然单体盈利能力承压，但可以通过规模优势和政策灵活性，在全球市场获得结构性份额提升。

市场普遍认为，中国收紧化肥出口是为了优先保障国内春耕和粮食安全。摩根斯坦利提供了另一个重要视角：中国近年对尿素和磷酸二铵/磷酸一铵的出口选择性管控，与全球地缘供给冲击高度相关——先是俄乌冲突，现在是伊朗与霍尔木兹海峡的紧张局势。

报告明确指出，市场此前担心中国会等到霍尔木兹海峡重新开放后才恢复尿素出口，但事实是底价取消发生在海峡仍然关闭的背景下。这意味着，中国对化肥出口的管控逻辑，已经从“避险性关闭”转向了“在风险中寻找出口窗口”。只要国内库存和价格可控，中国愿意在地缘风险尚未解除的情况下释放出口。

这个观察对全球化肥贸易流有深远影响。如果中国能够在霍尔木兹海峡关闭期间维持甚至增加尿素出口，那么中东生产商（特别是依赖海峡出口的沙特、卡塔尔、阿曼）将面临更大的市场份额压力。中国正在利用这个窗口期，建立从亚洲到美洲的新贸易路线和客户关系。

3. 下一个关键节点是磷肥出口，但硫磺供应链给出了更复杂的信号

报告将目光投向了8月，届时中国磷酸二铵/磷酸一铵的出口政策将迎来重要判断节点。摩根斯坦利的分析师明确表示，这个判断“目前更难预测”，原因在于硫磺市场的持续紊乱——不仅是因为霍尔木兹海峡，还因为俄罗斯近几周进一步减少了硫磺出口。

这里有一个容易被忽视的产业链传导关系。中国的磷肥生产高度依赖进口硫磺作为原料。如果硫磺供给不足或价格过高，即使政策允许磷肥出口，中国生产商也可能在成本上缺乏竞争力。换言之，磷肥出口的恢复，不仅取决于政策意愿，更取决于硫磺供应链的修复程度。

对于全球磷肥市场而言，这意味着8月的供给增量可能低于预期。如果中国磷肥出口延迟或规模受限，全球磷肥价格可能在北半球秋季施肥季前获得支撑。反之，如果硫磺问题缓解且中国恢复出口，则磷肥市场可能面临类似于尿素的下行压力。

报告提到，中国取消底价的时间点正好在印度最新一轮尿素招标

[... middle omitted ...]

中国卖多少量”。月度海关出口数据将成为比现货报价更重要的先行指标。

第二，硫磺市场的两个来源。霍尔木兹海峡的通行状况和俄罗斯的硫磺出口量，将共同决定中国磷肥出口的实际可行性。这两个变量中任何一个出现改善，都可能触发磷肥市场的重新定价。

第三，中国化肥企业的利润与现金流的背离。报告指出，过去几年的低出口政策导致国内化肥企业“利润和现金流受限”。如果出口量释放但价格承压，企业可能面临“量增利薄”的局面。这对中国化肥行业的资本开支和产能扩张计划有抑制作用，进而影响全球中长期供给曲线。

这些观察框架可以帮助读者在纷繁的短期新闻中，识别出真正具有结构性意义的信号。

中国取消尿素出口底价，表面上是一个战术调整，实则揭示了中国在全球化肥市场中的角色正在发生根本性转变。从被动的供给限制者，到主动的边际定价者，这一转变的完整影响需要时间显现。而磷肥出口能否在8月恢复、印度招标中中国企业的参与深度、硫磺供应链的修复节奏，这些尚未完全展开的变量，将决定这一轮定价权转移的速度和幅度。

\subsection{[R032] BARC：企业级服务器的意外走强，正在改写服务器市场的增长叙事}
{\small\color{refgray}归类：AI产业链：从GPU短缺到存储与计算架构再定价，企业级需求加速爆发}

BARC：企业级服务器的意外走强，正在改写服务器市场的增长叙事

市场对AI服务器的关注从未减弱，但真正让这份BARC研报值得反复阅读的，不是它确认了超大规模云厂商的资本开支仍在攀升，而是它在企业级服务器这个被普遍认为“平稳甚至萎缩”的细分市场，捕捉到了一个出乎意料的信号。

在最近的财报季中，多家服务器厂商的企业级业务表现显著超出预期。BARC团队在更新其服务器模型后，给出了一个值得产业决策者和投资者深思的判断：企业级服务器收入的增长，并非仅仅来自价格驱动的短期波动，其背后正在浮现由代理式AI和推理需求催生的结构性机会。当然，这并不意味着企业级市场将迎来长期繁荣，但至少在未来一年内，它正在改写整个服务器市场的增长曲线。

这份报告的核心价值不在于提供了多精确的数字，而在于它揭示了一个正在发生的结构性转变：当所有人都盯着云端的千亿级资本开支时，企业端的采购逻辑和需求形态正在悄然变化。理解这个变化，比记住任何营收预测都更重要。

1. 云服务器的高增长已是明牌，真正的增量来自企业级市场的“意外跃升”

BARC的模型调整，最引人注目的部分并非云服务器。云端的增长逻辑非常清晰：超大规模云厂商的资本开支预期已从2025年的约4230亿美元跃升至2026年的超过7000亿美元。BARC据此将云服务器的五年复合年增长率预期从约59\%上调至约62\%。这是一个巨大的数字，但市场对此已有充分预期，甚至已经部分定价。

真正的意外发生在企业级服务器市场。BARC此前预计企业级服务器收入在2026年仅有约2\%的增长，但在最新的财报季中，戴尔、慧与等OEM厂商的数据彻底颠覆了这一判断。戴尔预计其传统服务器业务在本财年将增长超过60\%，慧与则报告传统服务器订单实现了三位数增长。BARC因此将2026年企业级服务器收入增长预期大幅上调至约33\%。

这个数字的冲击力在于，它发生在市场普遍认为企业级IT支出趋于保守、云化迁移持续进行的背景下。企业级市场的“意外走强”，说明存在一股被低估的需求力量。BARC认为，这种力量主要来自两个方面：一是代理式AI和推理应用对传统服务器的新需求，二是采购方在价格持续上涨的情况下，仍然维持了比预期更好的购买行为。

2. 价格与数量的博弈：ASP推高了营收，但单位销量才是真正的信心指标

BARC在报告中反复强调一个关键矛盾：企业级服务器营收的大幅增长，究竟有多少来自价格上涨，又有多少来自真实的销量增长？这个问题的答案，决定了当前的增长是否可持续。

从报告披露的细节看，平均销售价格（ASP）无疑是主要驱动力。戴尔和慧与都提到了“丰富的配置”和“更高的组件成本”带来的价格上升。BARC估计，ASP的涨幅可能接近40\%-50\%。但真正让分析师感到意外的是，在如此大幅度的价格上涨面前，单位销量并未如预期般大幅下滑。戴尔明确提到，尽管ASP上升，其传统服务器在当季仍然实现了单位销量的增长。

这意味着，企业客户对服务器性能的需求正在发生质变。代理式AI和推理任务对算力的要求，不再是简单的“更多CPU核心”，而是需要更复杂的配置、更大的内存、更强的互联能力。这种需求变化，使得企业愿意为更高性能的服务器支付溢价。BARC的判断是，当前的增长中，ASP贡献了大部分，但单位销量的韧性表明，需求端确实存在结构性的支撑，而非仅仅是采购周期的提前。

不过，BARC也谨慎地指出，这种“量价齐升”的局面可能难以长期维持。他们预计，随着价格继续上涨，需求侵蚀最终会显现。因此，对于2027年，他们仅给出了约1\%的增长预期，且明确表示这一预期存在下修风险。

3. 云服务器市场的赢家格局已定，但企业级市场的竞争仍在胶着

将云服务器和企业级服务器分开来看，BARC的研报描绘了两幅截然不同的竞争图景。

在云服务器市场，白牌厂商（包括

[... middle omitted ...]

服务的依赖度更高，对白牌服务器的接受度相对较低。BARC在报告中也指出，一旦企业级推理和主权AI需求真正起量，品牌厂商的份额有望获得再分配。

慧与的战略转型尤其值得关注。BARC认为，在收购瞻博网络之后，慧与正在从一家服务器公司转型为一家网络公司。其网络业务预计将占全年总运营利润的50\%以上。这意味着，即使企业级服务器市场出现意外增长，慧与的核心价值驱动因素可能已经改变。

4. 报告尚未完全回答的关键问题：价格持续上涨何时会反噬需求？

BARC的这份报告数据扎实、逻辑清晰，但它也留下了一个所有读者都关心、却尚未完全回答的问题：当ASP持续以接近50\%的幅度上涨时，企业客户的承受极限在哪里？

报告指出，当前的需求韧性部分来自“采购提前”，即客户为了锁定当前价格而提前下单。但这一行为本身就是不可持续的。此外，报告也承认，如果价格继续上涨，需求侵蚀“最终会赶上来”。但“最终”是多久？是2026年下半年，还是2027年？不同的时间节点，对应着完全不同的投资和产业决策。

\subsection{[R033] Citi：市场低估的不是新能源渗透率，而是出口定价权的结构性重构}
{\small\color{refgray}归类：新能源与汽车：中国车企欧洲优势不在价格，储能市场供给端结构性收紧}

Citi：市场低估的不是新能源渗透率，而是出口定价权的结构性重构

这份Citi于2026年6月8日发布的CPCA 5月销量追踪报告，表面上是每月一次的数据更新，但若只看到“5月新能源乘用车批发135万辆，同比+10.6\%”这个数字，就错过了报告真正有价值的信息。

真正值得关注的判断是：中国汽车产业正在经历从“国内渗透率竞赛”到“全球定价权争夺”的转折点。 5月的数据揭示了一个结构性信号——新能源出口占比首次突破54\%，而国内零售端却出现了新能源销量同比下滑7.5\%的反直觉现象。这两个方向同时发生，意味着市场惯用的“渗透率单边上升”叙事已经不足以解释当下的产业逻辑。

Citi这份报告提供的不是乐观或悲观的情绪，而是一个需要重新校准的观察框架。以下是我们从数据中提炼出的四个核心洞察，以及一个报告尚未完全回答的关键问题。

5月新能源乘用车批发135万辆，同比+10.6\%，但零售仅95万辆，同比-7.5\%。批发与零售之间40万辆的差距，远超正常季节波动。Citi估算5月整个乘用车行业去库存约8.2万辆，而新能源板块的批发-零售差值意味着，新能源渠道库存仍在累积而非消化。

这背后不单是需求疲软，更可能是车企在国七排放标准切换预期下的主动铺货行为。传统燃油车批发同比下滑21\%，零售更是暴跌39\%，说明经销商对燃油车库存极度谨慎。而新能源批发保持正增长，更多是车企为抢占下半年市场份额而提前铺货——这是一种“以库存换份额”的策略性博弈。

对于投资者而言，批发数据的高增长需要打折扣。真正应该关注的不是月度批发量的绝对值，而是零售端能否在未来2-3个月内跟上批发节奏。如果零售持续落后，下半年价格战的烈度可能超出预期。

2. 出口成为唯一没有“杂音”的增长极，但结构性风险正在转移

5月乘用车出口78.4万辆，同比+75.1\%；其中新能源出口42.4万辆，同比+112.6\%。前5个月累计出口338万辆，同比+69\%。这些数字的含金量在于：出口增长几乎没有受到国内价格战和库存压力的干扰。

更值得关注的是新能源出口占比从去年同期的45\%跃升至54\%。这意味着中国汽车出口的驱动力正在从燃油车切换为新能源车。Citi报告没有直接讨论的是，这种切换将改变中国汽车出口的定价结构——新能源车平均单价高于燃油车，出口结构升级会直接改善车企的海外收入质量。

但风险也在转移。5月新能源出口同比翻倍，基数效应将在下半年逐步显现。更重要的是，欧盟、美国等主要市场的贸易政策不确定性正在上升。Citi报告未展开讨论的是：如果出口增速从100\%+回落至30-40\%，国内产能压力将如何重新分配。出口不再是“无限缓冲垫”，而是需要更精细化的区域和产品组合管理。

5月A00级纯电动车批发量同比暴跌44\%，而B级纯电动车同比增长42\%。这两个数字放在一起，揭示了一个清晰的消费结构迁移：低端电动车正在被市场淘汰，中高端电动车成为增长主力。

这不是短期波动，而是中国新能源市场从“政策驱动”转向“产品力驱动”的必然结果。A00级电动车过去依赖补贴和限购城市的牌照红利，当补贴退坡、消费者对续航和智能化的要求提升，这类产品的生存空间快速收窄。B级电动车增长则受益于家庭增购需求和换购需求的双重拉动。

对车企而言，这意味着产品线策略必须做出调整。继续在A00级市场打价格战的企业将面临量价齐跌的困境，而能够在中高端市场建立品牌溢价的企业，将在下一轮竞争中占据有利位置。Citi报告中特斯拉5月国内零售环比+80\%的数据，也间接印证了品牌力在消费升级周期中的重要性。

5月PHEV（不含增程）批发37.2万辆，同比+10.5\%；而EREV（增程）仅9.5万辆，同比-24.9\%。这是近年来少见的增程式同比负增长。

增程路线曾是 年增长最快的细分赛道，但5月的数据表明

[... middle omitted ...]

价值将重新凸显。

报告尚未完全回答的关键问题：出口的“真实利润率”能否支撑估值重估

Citi报告提供了详尽的出口数量数据，但一个关键变量没有被充分讨论：出口的利润率结构。中国新能源车出口到欧洲的均价约为国内售价的1.5-2倍，但扣除关税、物流、渠道费用后，实际利润率是否显著高于国内？不同车企的海外定价策略差异很大——有的以品牌建设为主，牺牲短期利润；有的则直接复制国内价格战模式。

这个问题的答案，将直接决定出口增长能否转化为可持续的盈利改善。如果出口利润率并不比国内高出太多，那么市场对“海外业务重估”的叙事就需要重新审视。这正是完整报告中最值得深入阅读的图表和假设部分——Citi分析师很可能在附录中提供了分区域、分车型的利润率敏感性分析，而这些细节是月报摘要无法涵盖的。

出口数据的繁荣背后，利润率、区域风险、政策不确定性这三个变量还没有被市场充分定价。Citi这份报告的真正价值，不在于告诉你5月卖了135万辆新能源车，而在于它提供了一套拆解这些变量的数据框架。

\subsection{[R034] JPM：市场误读了SOCAMM降配，真正的机会藏在供给硬约束和NVIDIA的生态锁定里}
{\small\color{refgray}归类：AI产业链：从GPU短缺到存储与计算架构再定价，企业级需求加速爆发}

JPM：市场误读了SOCAMM降配，真正的机会藏在供给硬约束和NVIDIA的生态锁定里

过去一周，全球存储器板块经历了一次幅度约11\%的急跌，同期费城半导体指数仅下跌6\%。触发因素是一则看似利空的消息：NVIDIA下一代Vera Rubin超级计算系统的CPU内存模组SOCAMM，单颗CPU配置可能从192GB降至96GB。市场的第一反应是“CPU驱动的高阶内存需求逻辑被削弱了”，随即抛售了三星电子、SK海力士等头部厂商的股票。

但JPM在最新发布的这份研报中，提出了一个与市场情绪截然相反的判断。这份报告的核心主张是：SOCAMM降配并非需求疲软的信号，而是供应极度紧张的副产品。更关键的是，NVIDIA与SK海力士刚刚签署的多年技术合作协议，以及Computex 2026上浮现的AI存储架构变革，正在将存储器产业从“周期博弈”推入“结构锁定”的新阶段。

对于关注AI基础设施投资的决策者而言，现在需要区分两类信息：一类是短期噪音，另一类是正在重塑行业竞争格局的长期变量。SOCAMM降配属于前者，而NVIDIA的生态圈层化、eSSD在AGI架构中的基础设施化，以及HBM4向HBM4E的加速迭代，都属于后者。JPM认为，过去一周的调整恰恰提供了一个“买入回调”的窗口。

以下是我们从这份研报中提炼出的四个核心洞察和一个尚未完全解答的关键问题。

1. SOCAMM降配不是需求萎缩，而是供应瓶颈下的被动重组

市场之所以对SOCAMM从192GB降至96GB反应剧烈，是因为它触及了投资者对AI内存需求增速的核心假设。如果单颗CPU的内存容量腰斩，那么即便Vera Rubin出货量增长，整体内存比特需求也可能达不到预期。

但JPM的供应链核查得出了一个相反的结论：降配并非因为NVIDIA认为96GB已经足够，而是因为高容量SOCAMM模组的供应根本跟不上需求。 报告明确指出，96GB规格的SOCAMM需求正在上升，而192GB版本的短缺才是配置调整的真正原因。更关键的是，从整体比特采购量来看，需求并未减少——这意味着降配将被更高数量的低容量模组所弥补。

JPM在报告中提供了一张敏感性分析表格。假设2026年Vera Rubin出货1050个节点，无论96GB与192GB的配置比例如何变化，NVIDIA对SOCAMM的总比特需求始终稳定在约33.88亿Gb当量（以8Gb为单位）。换句话说，内容量降配与出货量上行之间存在一个近乎完美的对冲关系。 市场只看到了前者，忽略了后者。

这一节的核心含义是：投资者不应将“单机内容量下降”等同于“总需求下降”。在供应受限的市场中，客户被迫接受更低的配置，但这恰恰说明需求本身并未减弱。如果供应瓶颈在未来得到缓解，内容量还有进一步上行的空间。

2. NVIDIA与SK海力士的“全栈”合作，标志着存储器供应商从供应商升级为生态伙伴

在Computex 2026期间，NVIDIA与SK海力士宣布了一项“多年技术合作”协议，SK海力士被正式指定为NVIDIA的“最大内存合作伙伴”，协议期两年以上且附带续约选项。但这份合作的价值远不止于HBM。

JPM指出，这项合作覆盖了NVIDIA的整个产品线：Vera Rubin超级计算机（HBM/DRAM/NAND）、Vera CPU（SOCAMM2）、RTX Spark PC（LPDDR5X）、以及Jetson Thor机器人平台。这意味着SK海力士实际上被嵌入到了NVIDIA正在构建的三个AI市场——AI基础设施、个人AI、以及物理AI——的每一个核心环节。

这一节的核心洞察在于：存储器供应商的竞争维度正在从“谁能提供更好的产品”转向“谁能与NVIDIA的路线图实现更深度的协同”。 过去，存储器是一个高度标准化的商品市场，客户可以根据价

[... middle omitted ...]

的展示内容给JPM留下了深刻印象。其中最具颠覆性的信号是：NAND SSD正在从传统的“数据仓库”角色，演变为AI推理流水线中的主动计算层。

Solidigm提出的“G3/G3.5/G4”存储分层架构清晰地说明了这一点。在AI推理过程中，KV Cache（键值缓存）的存储和检索效率直接决定了推理延迟。Solidigm的数据显示，将KV Cache卸载到NAND SSD上，相比重新计算上下文，可以将“首次令牌生成时间”缩短最多27倍。这意味着，在推理规模持续扩大的背景下，高性能SSD不再是可选项，而是必要的基础设施。

更值得关注的是Kioxia的进展。这家NAND厂商在Computex上展示了与NVIDIA的深度合作，其AiSAQ软件解决方案被用于超大规模RAG（检索增强生成）服务器。JPM特别指出，一家硬件厂商开发出具有粘性的软件解决方案，这在NAND行业中极为罕见。 这种“硬件+软件”的绑定模式，可能会成为Kioxia在与西部数据、美光等对手竞争时的差异化优势。

\subsection{[R035] JEF：市场对AI颠覆软件的恐惧过度，中国软件公司反而更具韧性}
{\small\color{refgray}归类：AI竞争格局：从模型性能转向Agent生态与成本效率，中国软件更具韧性}

JEF：市场对AI颠覆软件的恐惧过度，中国软件公司反而更具韧性

这份研报的核心判断，与过去半年市场的主流叙事截然不同。当全球投资者因为Claude Cowork、GPT-5.5等AI Agent的涌现而抛售软件股时，JEF的分析团队提出了一个反直觉的结论：AI对软件行业不是“末日审判”，而是一场“自然选择”。更关键的是，中国软件公司在这场选择中，可能比美国同行拥有更强的生存韧性。

这份报告发布于市场情绪最为悲观的时刻。自2025年10月以来，美国软件板块因AI颠覆SaaS的担忧大幅下跌，Anthropic在2026年1月12日发布的Claude Cowork更是将恐慌推至顶峰，导致板块在3月31日前累计下跌约23\%。此后市场虽有所修复，但中国软件板块的处境更为艰难——年初至今下跌13\%，远超美国软件6.4\%的跌幅。然而，JEF认为，这种下跌更多源于盈利下修和估值压缩，而非AI颠覆本身。

为什么中国软件公司反而可能更安全？为什么JEF将金蝶列为中国软件板块的首选？这些判断背后，是对软件行业底层商业逻辑的重新审视。

1. AI对软件的影响是“自然选择”而非“末日”，关键在于护城河的厚度

JEF将AI对软件行业的影响定义为“自然选择机制”，而非市场的普遍认知——“SaaS末日”。这一判断的核心在于，AI并不会让软件变得无关紧要，而是会重塑软件的架构和经济模型。

报告指出，AI Agent能够自主执行任务，但它们仍然依赖软件的API接口来获取工作流、业务逻辑和数据。这意味着，拥有强大护城河的软件公司——那些深度嵌入客户核心业务流程、掌握专有数据、拥有严格合规资质的公司——不仅不会被颠覆，反而可能进化为AI原生平台。而那些护城河薄弱、对AI反应迟缓的公司，则面临被替代的风险。

这一框架将软件公司分为两类：一类是“系统记录”的持有者，另一类是“界面工具”的提供者。前者如ERP、财务IT、工业软件，它们是企业运营的中枢神经，替换成本极高；后者如简单的CRM或项目管理工具，更容易被AI直接替代。这种分化意味着，投资者不能简单地用“软件板块”这个标签来概括AI的影响，而必须深入到每家公司的业务本质。

这是JEF研报中最具洞察力的部分。报告与7家中国软件公司进行了深入交流，发现中国软件市场的商业模式与美国存在根本性差异，这种差异恰好构成了对AI颠覆的天然防御。

美国SaaS公司普遍采用按席位收费的模式。当AI Agent开始替代人类操作员时，每用户平均收入直接承压——企业不再需要为每个员工购买软件许可，因为AI可以完成原本需要多人协作的任务。这是市场恐慌的核心逻辑。

但中国软件公司的定价模式完全不同。报告指出，许多中国软件公司采用项目制或模块化定价，而非按席位收费。它们主要为大型国有企业服务，需要进行大量的系统集成和定制开发。数据安全顾虑和重资产思维，导致公有云SaaS在国企中的渗透率有限。这意味着，即使AI替代了部分人力，中国软件公司的收入也不会直接受到冲击——因为它们的收入来源是项目交付和功能模块，而不是人头数。

这一判断的意义在于：市场将美国SaaS的恐惧情绪直接复制到了中国软件板块，但两者的商业基础完全不同。中国软件公司的收入结构决定了它们对AI替代的敏感度天然更低。

JEF的数据分析揭示了一个关键事实：中国软件板块的下跌与美国软件板块的下跌，驱动因素截然不同。

美国软件板块年初至今下跌6.4\%，但市场对2026年收入的盈利预期实际上在1月12日后上调了2\%。这意味着下跌纯粹由估值压缩驱动——市场在情绪上过度反应，但基本面并未恶化。而中国软件板块年初至今下跌13\%，市场对2026年收入的盈利预期在同期却下调了5\%。这说明中国软件的下跌是“双杀”：盈利下修叠加估值压缩。

盈利下修的根源，并非AI颠

[... middle omitted ...]

选。这一判断基于六个维度，每个维度都指向同一个结论：金蝶拥有中国软件公司中最强的AI防御力。

首先，金蝶的ERP系统深度嵌入客户的核心业务流程，作为“系统记录”的持有者，替换成本极高。其次，金蝶积累了大量的专有数据和行业知识，AI公司很难复制这些资产。第三，模块化定价模式使得收入不受客户裁员影响。第四，金蝶已经推出了自己的AI原生ERP套件和Lingee AIOS，并给出了2026年10亿元人民币的AI收入指引。第五，SaaS收入占比已超过50\%，驱动两位数的收入增长和毛利率扩张。第六，当前估值仅为2026年预期市盈率的0.7倍，具有吸引力。

更重要的是，报告中的AI防御力排名显示，金蝶在15家中国软件公司中得分34分（满分40分），位居前列。这一评分综合了工作流集成深度、专有数据、对幻觉的零容忍度、供应商资质等多个维度。金蝶在这些维度上的表现，使其成为AI时代最有可能幸存并壮大的中国软件公司之一。

5. 报告尚未完全回答的关键问题：AI带来的成本压力如何消化？

\subsection{[R036] 摩根斯坦利：市场正在用不同估值逻辑为互联网公司定价，但真正的分歧不在AI硬件与软件之间}
{\small\color{refgray}归类：AI竞争格局：从模型性能转向Agent生态与成本效率，中国软件更具韧性}

摩根斯坦利：市场正在用不同估值逻辑为互联网公司定价，但真正的分歧不在AI硬件与软件之间

这周美国互联网板块经历了一场令人瞩目的分化。整体互联网公司下跌了5\%，但跌幅分布极不均匀。亚马逊下跌了9\%，Meta下跌了6\%，而Pinterest却逆势上涨了7\%。在同一周内，同一板块内出现如此巨大的差异，市场到底在交易什么？

一个常见的叙事是：市场正在从AI硬件转向AI软件。但这个判断过于笼统，甚至可能误导决策。摩根斯坦利这份6月5日发布的研报，通过一张详尽的估值表和一组长达一周的板块表现数据，揭示了一个更精确的信号：市场不是在简单地切换板块偏好，而是在对每个公司重新评估其商业模式在AI时代下的护城河宽度。

这份报告最值得关注的判断是：互联网公司的估值离散度正在急剧扩大，传统倍数框架已经无法解释股价表现。真正决定资金流向的，是企业能否将AI能力转化为可定价的竞争优势，而非仅仅作为成本节约的工具。

1. 电商与数字广告的估值逻辑正在分道扬镳，但原因不是市场风格切换

从摩根斯坦利提供的估值表来看，数字媒体板块的2026年EV/EBITDA中位数是8.5倍，而电商/市场板块是8.7倍，两者几乎相同。但过去一周的表现却天差地别：电商板块市值加权平均下跌了8.9\%，而数字广告板块仅下跌3.8\%。

这不是市场在简单地“卖出电商、买入广告”。如果看个股，Pinterest上涨了7\%，而同样属于数字广告的Meta却下跌了6\%。如果是板块轮动，不应该出现同一个子板块内部如此大的分化。

真正的解释在于：市场正在对每个公司独立评估其AI投资的回报周期。那些AI投入尚未转化为可见收入增长的公司，即使属于热门赛道，也被重新定价。而那些能够讲清楚AI如何直接驱动广告变现效率提升的公司，即使规模较小，也获得了溢价。

这意味着，对于投资者而言，与其关注“AI软件还是硬件”这样的宏观判断，不如回到每个公司的商业模式本身，追问一个更具体的问题：这家公司的AI投资，是在扩大其与竞争对手之间的差距，还是在缩小差距。

亚马逊是本周跌幅最大的大型互联网公司之一，市值加权下跌了9.1\%。但它的估值并不贵：2026年P/E仅25.9倍，低于其TTM平均的约32倍，也低于谷歌的25倍和Meta的18倍。

如果只看估值倍数，亚马逊似乎处于历史低位，应该是买入机会。但摩根斯坦利的数据揭示了一个值得警惕的信号：亚马逊的自由现金流收益率在2026年和2027年均为“NM”（无意义），这意味着其FCF相对于市值的比例极低甚至为负。

亚马逊的困境在于：它同时在进行三场大规模资本开支竞赛——AWS的AI基础设施、电商物流网络、以及新业务扩张。每一场都是必要的，但三场同时进行，导致自由现金流的释放被持续推迟。市场正在问一个尖锐的问题：这些投资的回报率，什么时候才能超过资本成本？

这与Pinterest形成了鲜明对比。Pinterest本周上涨了7\%，其2026年FCF收益率高达8.8\%，且2027年预计进一步改善至10.1\%。对于一家仍在增长的数字广告公司而言，这种自由现金流生成能力是稀缺的。市场愿意为此支付溢价。

这里的含义是：在利率环境依然存在不确定性的背景下，市场对资本回报率的重视程度正在上升。那些能够同时实现增长和现金流改善的公司，将获得估值重估的机会。

3. 共享经济与旅游板块的韧性，暗示消费者行为正在发生结构性变化

与电商板块的剧烈下跌相比，共享经济和旅游板块表现出明显的韧性。共享经济板块市值加权平均仅下跌0.1\%，旅游板块下跌0.3\%。Uber甚至微涨0.4\%，Instacart上涨3.7\%。

这种分化不是偶然的。共享经济和旅游板块的共同特征是：它们提供的是“体验型”消费，而非“商品型”消费。当消费者对经济前景感到不确定时，他们倾向于削减实

[... middle omitted ...]

要么意味着风险未被充分定价。报告没有直接回答这个问题，但估值差异本身就是一个值得深挖的信号。

4. 报告没有完全回答的关键问题：AI投资的分化何时会反映在利润表中

摩根斯坦利这份报告的核心价值在于它提供了一幅详尽的市场估值地图，但它留下了一个尚未回答的关键问题：当前市场对AI投资的定价分歧，是基于理性预期还是情绪驱动？

从P/E倍数来看，AppLovin的2026年P/E高达36.6倍，而Criteo仅为4.2倍。两家公司都从事数字广告技术，但市场对它们的AI能力赋予了截然不同的估值。AppLovin的AI驱动广告平台被市场视为颠覆性技术，而Criteo的传统重定向广告业务则被市场视为被颠覆的对象。

这种估值差异是否合理？报告没有给出明确的判断。但一个值得注意的数据是：AppLovin的短期利率为4.7\%，相对较低，说明做空力量不强。而Criteo的短期利率仅为2.1\%，同样不高。这意味着市场对这两家公司的分歧，更多体现在多头对未来的预期上，而非空头的质疑上。

\subsection{[R037] NOM：2万亿投资计划不仅是短期催化剂，更是中国AI供应链国产化进程的加速器}
{\small\color{refgray}归类：未被主题摘要引用}

NOM：2万亿投资计划不仅是短期催化剂，更是中国AI供应链国产化进程的加速器

市场对于中国AI领域的投资预期，往往停留在“政策刺激”的短期逻辑上。但NOM最新的一份报告揭示了一个更深层的信号：中国正在以国家意志，系统性地构建一个自主可控的AI基础设施体系。这份报告的核心判断并非简单的“政府要花钱了”，而是“这笔钱将如何重塑中国AI供应链的竞争格局”。

这份报告发布于2026年6月9日，针对的是当日彭博社报道的中国计划在未来五年内投资约2950亿美元（约合人民币2万亿元）建设全国数据中心网络的消息。NOM的分析师团队迅速给出了他们的专业解读：这并非一个全新的消息，但它的意义在于确认了政府长期、坚定的决心，并指明了最直接的受益方向。

为什么这个判断值得决策者高度关注？因为它回答了一个核心问题：在先进AI芯片获取受限的背景下，中国AI产业的增长路径在哪里？NOM的答案是，路径不再单纯依赖算力的绝对领先，而是转向了基础设施的网络化、国产化，以及将规模优势转化为议价权。

1. 2万亿投资的核心不是“花钱”，而是“建网”，这改变了AI产业的竞争维度

NOM报告的核心洞察在于，它没有将这笔投资仅仅视为一次性的财政刺激，而是将其解读为一次基础设施的“网络化”重构。根据报告，中国发改委等关键部门正在起草一份蓝图，旨在建设一个全国性的互联计算枢纽网络。这意味着，投资的重点并非简单地建设一个个孤立的数据中心，而是要将它们连接成一个高效、协同的“算力网”。

这种“建网”逻辑，与过去几年中国在特高压、5G等领域的思路一脉相承。其战略意图非常清晰：通过国家主导的顶层设计，解决算力资源分布不均、利用率低下的问题。对于产业而言，这意味着未来的竞争将从“谁拥有更多算力”转向“谁能更高效地调度和利用网络内的算力”。这将深刻影响云服务商、IDC运营商以及网络设备供应商的商业模式和估值逻辑。

NOM特别指出，中国移动和中国电信等国有企业将运营大部分数据中心，并确保这些中心之间的互联。这暗示了一个关键变化：国有资本将在AI基础设施的“主干网”中扮演核心角色，而民营资本则可能在边缘计算、特定行业应用等“毛细血管”层面找到机会。这种分工格局，是理解未来投资机会的重要前提。

2. 报告明确点出的三大受益链条，其核心驱动力是“国产替代”

NOM在报告中清晰地列出了三大受益板块：AI计算/服务器、AI网络（包括光模块、交换机、光纤等）、以及IDC/云。但这份报告的价值远不止于罗列受益方。它更深层的判断是，这些板块的增长驱动力并非来自市场自然需求，而是来自“国产替代”这一强制性政策导向。

报告明确提到，政府的策略涉及依赖包括华为在内的本土供应商提供关键技术，如AI芯片。这直接点出了整个投资计划的核心矛盾与机遇：在无法获得最先进制程AI芯片的情况下，如何支撑起一个大规模、高算力的数据中心网络？答案必然是：通过系统架构创新、异构计算、先进封装等手段，最大化现有国产芯片的效率。

因此，对于投资者而言，仅仅关注“服务器出货量”是不够的。NOM的分析暗示，真正的投资机会可能在于那些能够帮助实现“国产替代”的细分领域。例如，为国产AI芯片提供配套散热、电源、高速互联解决方案的公司，其价值可能比单纯的服务器组装商更高。光模块和光纤的需求，也将因为网络内数据传输量的激增而获得长期增长动力。

3. 短期催化剂的背后，是“自主可控”这一长期战略的再次确认

NOM报告虽然将此消息定性为“短期催化剂”，但其分析的落脚点却是长期的战略定力。报告指出，尽管“2万亿投资计划”并非全新概念，但它的再次出现和具体化，表明“中国政府建立有竞争力且自给自足的AI价值链的长期雄心依然坚定”。

这提醒我们，不能将这份报告仅仅视为一次事件驱动的交易信号。它更像是一份对中国AI产业

[... middle omitted ...]

全回答的关键问题：先进芯片产能的瓶颈如何被突破？

NOM在报告中坦诚地指出了关键风险：“缺乏先进AI芯片产能”。这是整个投资计划中最核心、也最不确定的变量。报告本身并未给出这个问题的答案，但这恰恰是它最有价值的地方——它揭示了市场下一步必须关注的焦点。

2万亿的投资，如果没有足够的先进芯片来“喂饱”这些数据中心，那么投资的回报率将大打折扣。这里需要进一步思考的问题包括：国产芯片的良率和性能能否在五年内实现质的飞跃？通过Chiplet（芯粒）技术和先进封装，能否绕过单芯片的制程限制，构建出有竞争力的算力集群？政府的投资是否会直接流向芯片制造环节，以加速产能建设？

NOM的报告触及了这个问题，但并未深入展开。这恰恰是后续跟踪和研究的关键。任何关于芯片产能突破的正面消息，都可能成为引爆整个AI供应链估值的导火索。反之，如果瓶颈无法突破，那么庞大的基础设施投资可能会面临利用率不足的风险。

5. 对于高净值读者和产业决策者的观察框架：从“看需求”转向“看供给的国产化率”

\subsection{[R038] Bernstein：存储行业的结构性重定价远未完成，真正的机会在NAND和HDD的后半场}
{\small\color{refgray}归类：AI产业链：从GPU短缺到存储与计算架构再定价，企业级需求加速爆发}

Bernstein：存储行业的结构性重定价远未完成，真正的机会在NAND和HDD的后半场

当一家公司股价在12个月内上涨近4000\%，绝大多数投资者的第一反应是“太贵了”。但Bernstein在最近的第42届战略决策会议上，恰恰把这家公司——SanDisk——作为短期首选推荐。这不是追高，而是基于一个从2013年就开始论证、至今仍在展开的结构性判断：存储行业正在经历一次“新范式”的重定价，而市场目前只定价了前半段。

这份研报的核心主张是：AI正在驱动存储行业有史以来最大规模的需求浪潮，但市场对供给侧的长期结构性约束仍然定价不足。DRAM已经验证了“新范式”的逻辑，NAND正在复制它，而HDD则出现了历史上第一次“既集中又增长”的格局。这三个细分市场合在一起，意味着存储行业的周期性正在被系统性削弱，盈利中枢正在上移。

大多数投资者对AI存储需求的理解停留在HBM和DRAM上。这是对的，但不够。Bernstein将AI工作负载分为四个阶段：训练、基础推理、高级推理（长上下文、RAG、推理）、以及最终Agentic AI。每个阶段对存储的需求类型和强度都不同。

训练阶段对HBM的需求是“极端的”，对NAND和HDD的需求是“高的”。但到了Agentic AI阶段，HBM的需求回落至“高”，而常规DRAM、NAND和HDD的需求全部升至“极端”或“高”。这意味着，随着AI从“生成答案”进化到“执行任务”，存储需求的重心正在从最顶层的SRAM/HBM向下迁移到NAND和HDD。

这个判断的含义很直接：如果市场对存储的定价主要基于HBM和DRAM的紧缺程度，那么当NAND和HDD的需求加速时，相关公司的盈利弹性可能会超出当前预期。Bernstein对SanDisk的短期看多，核心逻辑正是NAND目前还处于周期早期，盈利上修的潜力尚未被完全计入。

存储行业历史上最痛苦的周期，往往来自消费电子需求的剧烈波动和客户的极端价格敏感。但这一次，需求结构发生了根本性变化。

Bernstein的数据显示，到2026年，数据中心在NAND、DRAM和HDD中的需求占比将分别达到52\%、50\%和89\%。而在2016年，这三个数字分别是14\%、18\%和35\%。更关键的是，数据中心客户——主要是超大规模云厂商——对价格敏感度远低于历史上的PC和手机厂商。

这个变化带来了一个直接后果：在过去六个月里，存储价格已经上涨了3到4倍。在传统周期中，如此剧烈的价格上涨会迅速刺激供给扩张，然后价格崩塌。但这一次，供给端受到技术瓶颈的约束，而需求端又由“不计代价买算力”的客户主导，价格在高位停留的时间可能远超历史经验。

这意味着，存储公司的盈利周期顶部可能比市场预期的更高、更持久。对投资者而言，用历史周期估值模型来给当前存储公司定价，可能系统性低估了它们的盈利中枢。

3. DRAM已经验证的“新范式”，NAND正在复制，但市场尚未完全定价

Bernstein在2013年和2014年的黑皮书中首次提出了“新存储范式”的概念。核心逻辑是：随着技术迁移带来的每晶圆产出增长持续放缓，供给增长将系统性低于需求增长，导致上行周期更长、下行周期更浅，盈利呈现“更高高点和更高低点”的特征。

DRAM是这一范式最完美的验证案例。从2012年到2027年，DRAM的每晶圆产出增长率从43.5\%持续下降到7.5\%，而需求增长率始终在15\%-22\%之间波动。供给增长持续低于需求增长的结果是，DRAM的上行周期越来越长，下行周期越来越温和。行业集中度也在提高，HHI指数从2001年的1500上升到目前的接近3000，进一步强化了供给纪律。

NAND的情况更复杂一些。平面NAND原本也在走DRAM的老路，但3D NAND的转型带来了一次性的每晶圆产出跃升，暂时

[... middle omitted ...]

AND的故事还有一定市场共识，那么HDD可能是最被忽视的结构性机会。

自2012年行业整合以来，HDD一直处于萎缩状态，因为NAND的持续侵蚀导致需求不断下降。但Bernstein认为，这种侵蚀已经基本耗尽。剩下的市场——主要是近线存储（Nearline）——本质上是一个容量市场，NAND在成本上无法与之竞争。

更关键的是，这是HDD历史上第一次同时满足两个条件：行业高度集中（HHI指数接近5000，是存储细分中最高的），且市场在增长。在过去的任何周期中，HDD要么是增长但分散的，要么是集中但萎缩的。当“集中”和“增长”同时出现，供给纪律和定价权都会显著改善。

Bernstein对Seagate的长期看好，正是基于这一判断。作为HAMR（热辅助磁记录）技术的领导者，Seagate不仅享受量价齐升，还能通过技术升级持续降低成本。在传统周期中，HDD公司的估值往往被压制，因为市场将其视为夕阳产业。但如果“既集中又增长”的格局成立，HDD的估值体系可能需要系统性重估。

\subsection{[R039] Citi：特高压订单翻倍背后的关键信号——电网设备竞争已从“订单量”转向“议价权”}
{\small\color{refgray}归类：特高压与电网设备：订单翻倍背后，竞争从订单量转向议价权}

Citi：特高压订单翻倍背后的关键信号——电网设备竞争已从“订单量”转向“议价权”

国家电网2026年第二批特高压设备招标结果出炉，总金额94亿元，是第一批的两倍有余，已接近2025年全年五批招标总额的43\%。这份来自Citi研报的数据，在多数市场参与者还在关注“谁拿到了最大份额”时，揭示了一个更值得追问的结构性变化：随着招标量放大、铜价上行、交付节奏加快，这一轮电网投资的真正赢家，可能不是订单最多的企业，而是能在规模扩张中守住甚至提升毛利率的企业。

换句话说，市场对特高压板块的定价逻辑，正在从“订单驱动”切换为“盈利质量驱动”。Citi在这份研报中，通过平高、思源、特变电工三家公司的订单结构、产品交付细节和估值对比，提供了这一判断的支撑证据。

2026年第二批特高压设备招标94亿元，放在历史坐标中看，意义远超数字本身。2025年全年五批合计221亿元，而2026年仅前两批就已达到184亿元。如果加上3月31日第一批的约50亿元（Citi在研报中标注了第一批为98.34亿元，但此处以正文数据为准），2026年上半年累计招标已接近去年全年。

更关键的是节奏的变化。2025年的招标集中在第四批（165亿元），呈现出“前低后高、年底集中”的特征。而2026年第二批就达到94亿元，且Citi明确指出，平高预计2026年还将有至少2个新的交流特高压项目获批，叠加2025年已获批但尚未招标的3个交流项目（达拉特-蒙西、攀枝花-西昌、浙江环网），后续招标量仍有支撑。

这意味着，电网投资不再是一次性的“脉冲式”放量，而是进入了持续高强度的建设周期。Citi在研报中援引的数据也佐证了这一点：2026年一季度全国电网资本开支同比增长40\%至1675亿元，其中土建工程是主要增量——政府有意通过电网基建拉动固定资产投资，而土建支出对FAI的拉动效果快于设备采购。

所以，这一轮招标翻倍的本质，不是简单的“订单多了”，而是电网投资从“项目储备”阶段正式进入“施工高峰”阶段。设备企业的订单确认节奏、产能利用率、原材料成本传导能力，都将在这个阶段接受真正的压力测试。

2. 平高拿下22\%份额的背后，是特高压GIS市场的寡头格局正在固化

平高以20.92亿元中标额、22.2\%的份额位居第一。但比份额数字更值得关注的是产品结构：94亿元招标中，1000kV GIS（气体绝缘开关设备）占了57.4\%，即54亿元。平高的核心竞争力恰好集中在这一产品线。

Citi在研报中提供了平高的交付细节：2026年预计交付15个间隔的1100kV GIS（同比增长36\%），以及超过100个间隔的750kV GIS（同比增长7.5\%）。其中750kV GIS在2025年的市场份额已达45\%，且招标量同比增长77\%。

这些数据合在一起，指向一个判断：特高压GIS市场已形成寡头格局。平高、中国西电、山东电工日立三家分走了前三位，合计份额超过50\%。对于新进入者而言，特高压GIS的技术壁垒、产能门槛、国网认证周期，构成了极高的进入成本。这意味着，现有龙头企业的市场份额不仅稳定，而且有进一步集中的趋势。

但这里有一个尚未被充分讨论的问题：份额集中是否意味着定价权增强？从平高自身预期来看，公司已明确表示第二批招标价格可能因铜价上涨而上调。如果这一预期兑现，那么平高的利润率弹性将超出市场当前模型的假设。反之，如果招标价格未能充分传导成本，则份额扩张反而可能侵蚀利润。

3. 思源电气的“2\%份额”被低估了——出口增速才是它的核心叙事

在94亿元招标中，思源电气仅中标1.96亿元，份额2.1\%，排名第五。如果只看国内特高压招标，思源似乎并不突出。但Citi将思源列为行业首选之一，核心理由是“出口增速超过80\%”。

这一判断背后的逻辑是：思源的竞争优势不在

[... middle omitted ...]

至29.6\%，且账上净现金。相比于国内SOE企业（如平高PE 19.8倍、ROE 10.8\%），思源的估值溢价背后，是市场对其海外增长持续性的定价。

这里有一个值得追问的盲点：出口业务的毛利率是否高于国内？如果出口订单的盈利能力更强，那么思源的实际利润弹性可能比市场预期的更大。Citi在研报中没有展开这一分析，但这是投资者需要自行验证的关键假设。

Citi在研报中提及，平高预计第二批招标价格可能因铜价上涨而上调。这是一个容易被忽略但可能影响整个板块定价逻辑的信号。

特高压设备的核心原材料之一是铜，而铜价在 年持续走高。如果企业能够将成本上涨传导给下游（国网），那么毛利率将保持稳定甚至改善；如果不能，则订单增长反而会带来利润压力。

从Citi的估值表格看，平高2026年预期ROE仅为10.8\%，在行业内处于偏低水平。如果平高能够在铜价上涨环境下维持或提升利润率，其ROE有显著上行空间，估值也将获得重估。反之，如果成本传导不畅，则当前19.8倍的PE可能并不便宜。

\subsection{[R040] 0005-41-BARC-U.S. Equity-Linked Strategies AI Capex Funding- Why Equity- Why Converts--260609}
{\small\color{refgray}归类：AI/算力：资本开支超越互联网泡沫，折旧冲击与融资结构成新焦点}

DeepSeek 生成 WeChat article 失败：Response ended prematurely

Personal reading notes and learning share only. Not investment advice.

最近投行研报拆了AI巨头们的融资策略，核心逻辑很清晰：AI投资周期太长、体量太大，光靠发债已经不够用了。

1️⃣ 为什么需要股权融资？ AI资本支出正在快速吞噬经营现金流。研报数据显示，2028年五大云厂商总资本支出预计达1.16万亿美元，而现金流覆盖比例可能逼近97\%。债务市场虽然仍开放，但无限度发债会挤压信用空间。股权融资的核心作用不是“缺钱”，而是： 拓宽融资渠道，减少对单一债务市场的依赖 降低杠杆率，为未来继续发债留出空间 给债券持有人一个“安全垫”，稳定市场信心

2️⃣ 可转债：债务和股权的中间地带 可转债兼具两者优势：不需要立即发行全部普通股，又能获得类股权资本。美国可转债市场今年已发行约780亿美元，创2003年以来同期最高。AI相关发行占比高达71\%。对于成熟大公司，可转债还能做到“隔夜定价”，速度极快。

3️⃣ 强制可转债：最干净的股权信用结构 谷歌和甲骨文都采用了这种结构。它能在资产负债表上获得近乎完全的股权处理，三年后强制转换为普通股，给市场一个明确的资本强化路径。不是最便宜的资金，但信号最强。

4️⃣ 谁可能是下一个？ Meta、微软、亚马逊现金流充裕，但研报认为它们也可能通过股权挂钩工具来主动管理资产负债表。一些AI基础设施公司已经在广泛使用可转债。

*核心观点**：这不是融资危机，而是管理团队在利用有利市场条件，提前为多年度AI建设储备资本。股权和可转债是战略工具，不是最后手段。

AI Capex Funding: Why Equity? Why Converts?

As AI capex becomes larger and longer duration, hyperscalers are broadening their funding mix beyond IG debt. Equity and equity-linked capital can preserve debt capacity, strengthen credit confidence, and support continued investment.

A reminder that we are entering the late innings of the 2026 Extel All-America Research Survey! If you’ve already submitted your ballot, thank you – we truly appreciate it. If not, our team would be grateful for your consideration for a ★★★★★ 5 Star Vote in the [Macro] → [Equity-Linked Strategies \& Portfolio Strategy] categories.

Equity capital raise does not signal "last resort" financing. Debt markets remain open to AI hyperscalers, but equity broadens the funding base, reduces reliance on any single channel, and gives issuers more flexibility 

[... middle omitted ...]

o IG debt rather than a replacement. In addition, the US convertibles market provides speed of access (one-day or overnight issuance) and structural flexibility, including dilution management tools, for issuers.

Mandatories offer the cleanest equity-credit structure. GOOGL and ORCL show how mandatories can provide equity-credit treatment on the balance sheet, conversion certainty, and balance-sheet signaling within a defined timeframe.

\subsection{[R041] 摩根斯坦利：CLO市场真正的分化不在区域，而在管理人的议价权}
{\small\color{refgray}归类：风险偏好：CLO市场从beta转向alpha，军工复合体面临工业基础结构性过时}

摩根斯坦利：CLO市场真正的分化不在区域，而在管理人的议价权

当大多数市场参与者仍在争论“美国还是欧洲”作为CLO配置的主战场时，摩根斯坦利这份6月CLO追踪报告揭示了一个更本质的信号：市场正在从“买什么区域”转向“买谁管理的资产”。这不是一个温和的风格切换，而是信用周期进入后半段时，市场定价逻辑的一次结构性重组。

这份报告基于摩根斯坦利团队近期跨大西洋与大量CLO投资者的直接对话，结合5月发行数据与信用基本面研判，给出了一组值得产业决策者与资产配置者认真对待的观察。5月新发行量从4月的60亿美元反弹至约170亿美元，而再融资与重置活动更达到年内最强的390亿美元——但在这组看似乐观的数字背后，投资者的叙事已经发生了微妙但深刻的转变。

报告的核心判断可以浓缩为一句话：CLO市场的超额回报正从beta（市场方向）转向alpha（管理人能力），而这一转移对资产定价、管理人格局乃至私人信贷市场的结构都将产生持续影响。

1. 管理人选股能力正在取代宏观判断，成为CLO定价的首要变量

在与几乎每一位投资者的对话中，摩根斯坦利团队都听到了同一个主题：管理人的选择已成为CLO投资中最重要的决策。这不是一句空洞的行业共识，而是基于一组正在加速变化的市场事实。

报告明确指出，当前信用利差在资本结构的多数层级上仍处于偏紧状态，而AI颠覆风险正在导致抵押品表现出现更大的离散度。这意味着，过去那种“只要市场方向判断对了，组合就能赚钱”的模式正在失效。投资者越来越关注管理人特有的alpha，而非广泛的市场beta。

这一判断有坚实的底层逻辑支撑。摩根斯坦利在报告中提到，当前周期的一个关键特征是：市场走势的方向性减弱，而不同组合之间的表现差异在拉大。这直接考验的是管理人的承销纪律、组合构建能力和交易执行水平。

具体到偏好，投资者明显倾向于那些拥有长期跟踪记录、经历过多个信用周期考验的成熟管理人。这一偏好在欧洲市场尤为强烈——许多欧洲投资者认为，对于新兴管理人而言，利差上的微小溢价并不足以补偿其风险。而在权益投资者群体中，对管理人控制权的要求更高，尤其是在私人信贷交易中。

这里有一个值得关注的信号：管理人之间的整合正在加速。新管理人不断涌现的同时，收购兼并也在频繁发生。大型管理人通过有机或非有机方式扩张规模，以拓宽能力范围、提升效率并深化与投资者和融资提供方的关系。报告指出，规模更大的管理人通常拥有更便捷的仓库融资渠道、更强的吸引大型锚定投资者的能力，以及更专业的资产处置和重组资源。

*这意味着什么：** 未来18个月，CLO管理人市场将出现明显的“头部集中”趋势。规模本身正在成为一种定价因子——不仅仅是效率优势，更是风险缓释能力的背书。对于资产配置者而言，这意味着“选管理人”的颗粒度需要从品牌层面下探到具体的组合构建方法论和AI风险评估框架。

2. 私人信贷CLO的风险溢价远未反映真实成本，40个基点是入场门槛而非机会信号

私人信贷CLO（PC CLO）是本次投资者对话中最受关注的话题之一，但情绪的转向值得警惕。报告描述了一个耐人寻味的画面：尽管近期利差已经走阔，但多数投资者仍认为当前的风险溢价不足以补偿私人信贷资产固有的复杂性、非流动性和估值不确定性。

摩根斯坦利在中期展望报告中已经给出了一个具体的违约预测：到2027年中，BSL（银团贷款）违约率将升至5.5\%，而PC（私人信贷）违约率将升至8\%。AI暴露的资产将被迫面对其短期到期墙。基于这一判断，报告认为PC与BSL之间的利差将进一步走阔，并明确指出40个基点的利差溢价是PC AAA级证券值得关注的入场点。

有趣的是，投资者的反馈与这一判断高度吻合。虽然多数投资者并不预期PC与BSL AAA级证券的利差溢价会回到几年前60-70个基点的水平——因为投资者基

[... middle omitted ...]

驱动型交易与套利驱动型交易，以及区分底层借款人的规模。那些由对资产和重组流程有更强控制权的管理人发起的融资交易更受青睐。借款人的EBITDA规模是另一个重要考量因素。

*这里有一个关键警示：** 多位投资者对中上市场（upper middle market）段位表达了更大的谨慎，因为该领域的契约文件和担保保护正在弱化。报告提到，在与BSL市场的激烈竞争中，中上市场段位的“cov-lite”交易正在增多。相比之下，中下市场（lower middle market）段位因拥有更强的担保保护、更小的贷款人群体和更大的重组影响力，被更积极地看待。

*这意味着什么：** 私人信贷CLO的“质量分化”正在加速。当前市场的错误定价不在于“是否应该配置私人信贷”，而在于“用同一把尺子衡量所有私人信贷交易”。那些能够区分融资驱动与套利驱动、能够识别中下市场中真实担保保护的投资者，将拥有显著的信息优势。40个基点的利差溢价不是机会信号，而是入场门槛——低于这个门槛，风险回报就不成立。

\subsection{[R042] 摩根斯坦利：全球矿业正在经历一场“主权供需”的结构性重塑}
{\small\color{refgray}归类：能源与大宗：供给硬约束重塑定价权，中国出口成金属市场泄压阀 / 黄金与矿业：黄金股回调是表象，矿业主权供需重塑竞争格局}

这份摩根斯坦利最新发布的澳大利亚材料行业研报，表面上是一份常规的行业数据汇编与个股偏好排序。但真正值得产业决策者和长期投资者关注的，是报告中一个正在加速形成的结构性趋势——全球关键矿产的供需两端，都在被主权力量重新组织。这不再是周期性的价格波动，而是贸易规则和竞争格局的底层逻辑变化。

报告提出的核心观察是：以中国新成立的中央采购机构（CMRG）为代表的需求侧主权集中，与印度尼西亚通过国有出口机构（DSI）对煤炭、棕榈油和铁合金进行中央化出口的供给侧主权集中，正在同时发生。过去，市场习惯将矿业视为自由市场定价的典范；但这份研报揭示，未来的游戏规则可能截然不同——国家正在成为最重要的交易对手。

为什么这个判断现在至关重要？因为澳大利亚的矿业公司正处于这场结构性变化的十字路口。澳大利亚是全球铁矿石、冶金煤和锂的关键供应国，但其矿业公司目前受到竞争法的严格约束，无法像主权买家或主权卖家那样进行协调定价或分配策略。摩根斯坦利提出了一个大胆的假设性框架：如果澳大利亚也建立类似的国家出口平台，会怎样？这种“主权对主权”的格局，将彻底改变当前基于边际成本和供需平衡的定价模型。

以下，我们将沿着报告的逻辑链条，逐一拆解这一趋势对行业、公司以及投资者观察框架的具体含义。

1. 主权买家与主权卖家的同步崛起，正在重新定义矿业的交易对手

摩根斯坦利研报最重要的贡献，是将两个看似独立的政策动向——中国的集中采购和印尼的集中出口——放在同一个分析框架下。这不仅仅是巧合，而是全球资源民族主义在贸易层面的具体表现。

从需求侧看，中国成立CMRG的目标是增强议价能力。对于一个占全球铁矿石进口量超过70\%的国家来说，分散的采购结构意味着巨大的价格劣势。通过一个中央机构统一对外谈判，中国试图将庞大的进口规模转化为实际的价格折扣。这并非没有先例，中国在稀土、焦煤等领域的集中采购已有实践。

从供给侧看，印尼通过DSI对煤炭、棕榈油和铁合金出口进行中央化，其官方目标是提高出口透明度、减少低开票和转移定价行为，并将更多收益留在国内。但实际效果是，印尼作为全球约55\%的棕榈油出口和约45\%的贸易动力煤供应国，正在获得对出口价格和流向的直接控制权。

这两股力量合在一起意味着什么？意味着矿业的定价机制正在从“市场均衡”向“双边谈判”演化。过去，一个铁矿石生产商面对的是数百个钢厂买家；未来，它可能面对的是一个国家级的采购机构。同样，一个下游冶炼厂过去可以向多个独立矿商询价；未来，它可能只能与一个国家级的出口平台交易。这种变化对定价透明度、合同条款和风险分配的影响，是深远的。

2. 澳大利亚的供应主导地位，使其成为这场博弈的核心玩家而非旁观者

摩根斯坦利在报告中用一张清晰的图表展示了澳大利亚在全球海运矿产供应中的份额：铁矿石约59\%，冶金煤约38\%，动力煤约20\%，锂约24\%。这些数字意味着，澳大利亚不是这场“主权供需”博弈的旁观者，而是核心参与者。

但报告同时指出一个关键矛盾：澳大利亚的矿业公司目前“协调能力有限，无法在定价、条款、分配、产量或营销策略上达成一致，否则将面临竞争法风险”。换句话说，在主权买家（中国CMRG）和主权卖家（印尼DSI）都在集中力量的时候，澳大利亚的供应端仍然是高度碎片化的。

这种碎片化对澳大利亚矿商意味着什么？意味着它们在谈判中处于结构性劣势。当你的交易对手是一个可以协调全国需求的国家机构时，你作为一家独立的矿业公司，即使拥有优质资产，也很难在价格谈判中获得有利地位。这正是摩根斯坦利提出那个假设性问题的背景——如果澳大利亚也建立国家出口平台，情况会怎样？

报告谨慎地指出，这种做法会带来贸易和主权风险，但同时也为澳大利亚生产商提供了一种应对“政府集中购买”的手段。这并非空穴来风。在铁矿石领域，澳大利亚与中国的贸易

[... middle omitted ...]

70\%以上。这是最极端的“双边垄断”结构，也是最容易催生主权干预的领域。报告中对BHP、RIO和FMG的不同评级，本质上反映了对这些公司在不同铁矿石价格假设下的估值敏感性。

动力煤则是另一个方向。印尼占贸易动力煤供应的45\%，且已经通过DSI实现了出口中央化。这意味着动力煤的供给侧已经完成了主权整合。对于澳大利亚的动力煤出口商（如Whitehaven Coal），它们面临的竞争环境正在变化——不是与印尼的独立矿商竞争，而是与一个拥有国家背书的出口机构竞争。

相比之下，锂和铜的供应集中度较低，且下游产业链更为分散，短期内受到主权干预的可能性较小。但这并不意味着可以忽略这一趋势。一旦铁矿石和煤炭的“主权化”模式被验证有效，其他矿种很可能被纳入同样的政策框架。

4. 报告隐含的个股偏好，本质上是基于不同“主权风险敞口”的排序

摩根斯坦利在这份报告中给出了明确的个股评级和偏好排序，但如果我们从“主权供需”的角度重新审视，会发现这些评级的背后逻辑远超传统的估值倍数比较。

\subsection{[R043] JPM：合成橡胶的供给缺口正在重塑化工定价权，而市场仍在关注需求侧}
{\small\color{refgray}归类：化工与材料：合成橡胶供给缺口重塑定价权，松下投资者日揭示AI电源机遇}

JPM：合成橡胶的供给缺口正在重塑化工定价权，而市场仍在关注需求侧

一份来自JPM亚洲能源与化工研究团队的报告，在2026年6月初释放了一个值得产业链决策者高度重视的信号：韩国NBL（丁腈胶乳）价格在5月攀升至每吨1633美元，4-5月均价创下五年新高。更关键的是，NBL与SBR（丁苯橡胶）的价差结构发生了根本性逆转——NBL的加工利润（Spread）自2022年以来首次超越SBR，达到每吨916美元，为2021年以来的最高水平。

这不是一个简单的价格波动。这是一次由供给端结构性短缺驱动的再定价，而市场的主流叙事仍然停留在需求复苏的节奏上。

这份报告的核心判断可以概括为一句话：当前合成橡胶市场的超额利润并非来自需求端的强劲反弹，而是来自原料石脑油/LPG短缺引发的供给刚性收缩，以及全球丁二烯体系与下游橡胶产品之间的传导断裂。 理解这一判断，比猜测下一个季度的油价走势更重要。

为什么现在必须重视这个信号？因为化工行业过去两年的核心叙事是“产能过剩、利润压缩”。几乎所有投资者都在关注需求何时恢复、库存何时去化。但JPM的这份报告揭示了一个相反的局部：在合成橡胶的特定细分领域，供给侧的收缩已经硬到足以让价格脱离成本曲线的牵引。如果这种供给刚性在未来2-3年持续，那么部分化工子行业的定价逻辑将从“需求驱动”切换为“供给约束驱动”。

1. 原料短缺不是暂时的扰动，而是正在改变合成橡胶的竞争格局

报告最引人注目的数据之一，是韩国NBL的出口量在2026年4-5月平均仅为4.4万吨，比2026年第一季度的月均水平下降了33\%。这不是需求疲软导致的减产。JPM明确将原因指向了中东局势后的石脑油供应短缺，并据此预测锦湖石化（Kumho Petchem）2026年第二季度的NBL产量将环比下降超过20\%。

这里需要拆解一个关键链条：石脑油短缺 -> 裂解装置开工率下降 -> 丁二烯供应紧张 -> NBL产量受限。但有趣的是，报告同时指出，亚洲丁二烯价格在5月环比下跌了27\%，从4月的四年高点每吨1945美元回落。这意味着丁二烯的下跌并未传导到NBL价格上——NBL价格反而环比上涨了3.4\%。

这个背离至关重要。它说明NBL的供给约束已经不仅仅来自上游原料的紧缺，而是叠加了更底层的结构性因素：韩国作为全球NBL核心供应国，其出口能力的下降是刚性的，无法通过其他地区的增产或原料价格下跌来快速弥补。报告提到，马来西亚手套制造商64\%的NBL需求依赖进口，主要来自韩国。当韩国出口量骤降时，马来西亚厂商面临的是“有钱也买不到货”的局面。Bloomberg报道的WRP Asia Pacific因原料短缺导致流动性危机而停止运营，正是这一传导链条的终端验证。

2. 利润结构的历史性逆转：NBL正在从“跟随者”变成“定价者”

JPM测算的NBL加工利润达到每吨916美元，同期SBR加工利润约为每吨900美元。这是2022年以来NBL利润首次超过SBR。表面上看，这只是一个数字的超越，但其背后的行业含义是：NBL正在从过去跟随丁二烯和SBR价格波动的“被动品种”，转变为具有独立定价能力的“主动品种”。

为什么这个转变重要？因为化工行业的大多数产品，其利润水平在长期内会被产能扩张和成本竞争压缩到接近边际成本。NBL过去五年平均价格约为每吨786美元，而JPM预测 年NBL价格将分别维持在每吨940美元和890美元——显著高于 年均值，即使假设油价回落至类似水平。

这意味着什么？报告给出的核心逻辑是“全球NBL产能增量微不足道，供需紧张将持续到2028年”。如果这一判断成立，那么NBL将成为化工板块中少有的、能够在中长期维持超额利润的品种。这对于持有相关资产的投资者而言，意味着估值框架需要从“周期股”切换到“结构性稀缺资产”。


[... middle omitted ...]

主导。当供给端的约束足够强时，原料成本下降带来的利润空间会被下游生产商保留，而不是传导给客户。

对于产业决策者来说，这意味着传统的“成本加成”定价模型需要修正。过去，化工产品的价格下限由原料成本决定；现在，在供给约束的子行业中，价格下限正在由“可获得的供给量”决定。这种定价逻辑的切换，会改变库存管理策略、长协谈判框架，乃至产能投资的决策依据。

4. 报告尚未完全回答的关键问题：丁二烯与NBL的背离能持续多久

尽管JPM的分析框架令人信服，但报告中仍有一个关键假设值得进一步验证：丁二烯价格下跌与NBL价格上涨之间的背离，其持续性的边界在哪里？

报告指出，丁二烯价格下跌的主要原因是油价回落、石脑油裂解装置检修后重启，以及ABS/SBR需求疲软导致的开工率下降。但问题在于，如果丁二烯价格持续下行，NBL生产商的原料成本优势将不断扩大，这会激励韩国以外的NBL产能加速复产或扩产。报告虽然强调“全球NBL产能增量微不足道”，但并未详细拆解现有产能的闲置率和潜在复产弹性。

\subsection{[R044] 摩根斯坦利：AI融资的“夏季攻势”已提前到来，市场定价权正从基本面转向供给节奏}
{\small\color{refgray}归类：AI/算力：资本开支超越互联网泡沫，折旧冲击与融资结构成新焦点}

摩根斯坦利：AI融资的“夏季攻势”已提前到来，市场定价权正从基本面转向供给节奏

市场参与者正在经历一个罕见的时刻：基本面依然强劲，但价格行为的主导力量已经悄然切换。摩根斯坦利在最新发布的《AI债务融资追踪》报告中给出了一个清晰的信号——截至5月31日，2026年全球AI相关债务发行量已达2360亿美元，是2025年同期的4倍以上。但真正值得关注的不是这个数字本身，而是它背后的结构性变化：四月的发行热潮之后，五月美国市场迅速降温，但超大规模云服务商（hyperscalers）并未停止融资脚步，而是通过欧元、加元、瑞士法郎和日元等非美元市场发行了约240亿美元债券。

这意味着什么？AI基础设施的资本密集型属性，正在催生一个前所未有的、跨币种、跨资产类别的融资周期。摩根斯坦利预测，2026年全年全球AI相关债务供应将达到约5700亿美元，较2025年增长162\%。但这份报告最核心的判断并不在于供应量的增长，而在于一个被市场低估的变量：融资节奏的波动本身，正在成为比需求基本面更重要的定价因子。

报告明确指出，当前市场定价更多由供应预期驱动，而非基本面。这听起来像是一个短期技术面现象，但如果我们深入拆解，会发现这背后隐藏着一个更深层的叙事转换——AI基础设施建设正在从“概念验证”阶段进入“资本运作”阶段，而资本市场对此的定价逻辑，也正在经历一次范式转移。

1. 四月发行高峰揭示了AI融资的结构性特征，而非季节性波动

四月单月发行量超过740亿美元，创下年内新高。这个数字本身令人印象深刻，但摩根斯坦利的数据揭示了更关键的细节：在AI相关的高收益债发行中，85\%用于数据中心外壳建设；在投资级债券中，这一比例也达到了40\%。

这意味着，当前AI融资的核心用途并非购买GPU或研发投入，而是最“重”的物理资产——数据中心。这些项目融资结构（project finance）的兴起，标志着AI基础设施投资正在从科技公司的资产负债表内支出，转向独立的、有明确资产抵押的融资工具。这种转变的意义在于：它打开了更广泛的投资者基础，同时也将AI投资的信用风险与科技公司的运营风险进行了部分隔离。

对于信用投资者而言，这既是机会也是挑战。机会在于，数据中心作为实物资产，其现金流可预测性可能高于科技公司的整体业务；挑战在于，这些项目融资结构的信用评估框架尚未成熟，市场仍在学习如何定价这类资产的风险。

摩根斯坦利用了一个巧妙的措辞“May Be Catching Its Breath”来描述五月美国市场的发行放缓。从数据看，五月美国市场AI相关发行量骤降至约42亿美元，与四月的744亿美元形成鲜明对比。

但仔细观察会发现，同期超大规模云服务商在非美元市场发行了约240亿美元债务。这不是市场降温，而是融资渠道的主动多元化。这些发行主体在欧元、英镑基准中的代表性远低于美元市场，这意味着非美元市场仍有充足的容纳空间。换句话说，五月美国市场的“安静”更像是一次战术性转移，而非需求不足。

对全球信用投资者而言，这意味着AI融资正在成为一个真正的全球性资产类别。欧元、英镑、加元、瑞士法郎和日元市场的AI债发行，不仅拓宽了投资者的选择范围，也为不同货币的资产管理人提供了参与AI基础设施投资的机会。这是一个值得持续跟踪的趋势，因为币种多样化本身也会影响定价效率和流动性分布。

这是摩根斯坦利报告中最具洞察力的判断之一。报告指出，尽管基本面背景依然强劲，但当前价格行为主要由供应预期驱动。一个反直觉的现象是：尽管供应量大幅攀升，自一季度末以来，各资产类别的利差反而普遍收窄。

这听起来像是市场对供应的消化能力超预期，但摩根斯坦利的分析暗示，真正的原因可能更微妙——市场已经将“供应将持续增加”这一预期充分定价，而实际发行节奏的波动，反而成为短期交易的核心

[... middle omitted ...]

定价的“二阶效应”

摩根斯坦利预测2026年全球AI相关债务供应将达到约5700亿美元，而截至5月底已发行2360亿美元。这意味着下半年还有约3340亿美元的供应等待释放。

这个数字本身已经足够庞大，但真正值得关注的是它背后的“二阶效应”。报告提到，摩根斯坦利的股权研究团队预计，到2027年超大规模云服务商的现金资本支出将超过1万亿美元。这意味着债务融资只是冰山一角，更大的股权融资需求可能正在酝酿。如果股权市场出现大规模融资，将对资本结构产生深远影响——可能改变现有债务的信用质量，也可能稀释现有股东的权益。

此外，芯片融资（chip financing）正在成为新的热点。报告指出，芯片融资交易期限更短、采用完全摊销结构，这与数据中心建设的长期项目融资形成鲜明对比。这种差异化的融资结构，意味着AI产业链的不同环节正在形成各自的信用定价体系。投资者需要区分的是：你是在为物理资产（数据中心）提供融资，还是在为技术迭代（芯片）提供融资？两者的风险特征和定价逻辑截然不同。

\subsection{[R045] JPM：AI基础设施的下一轮投资机会，藏在松下投资者日被低估的细节里}
{\small\color{refgray}归类：化工与材料：合成橡胶供给缺口重塑定价权，松下投资者日揭示AI电源机遇}

JPM：AI基础设施的下一轮投资机会，藏在松下投资者日被低估的细节里

当市场还在为AI芯片的算力竞赛兴奋时，一份来自JPM日本团队的报告，将目光投向了AI基础设施中一个容易被忽视的环节——电源与材料。这份报告的核心判断不是关于GPU的迭代，也不是关于云资本开支的增速，而是指向一个更具体、更可追踪的叙事：AI基础设施的物理层正在经历从“够用”到“必须重构”的转折，而松下等传统电子巨头，正在这个转折点上占据一个被低估的位置。

这份报告发布于松下控股投资者日之后，聚焦于能源与产业两大板块的AI基础设施业务。报告中最引人注目的信号，并非松下已经公布的FY2028销售目标——这些在5月的业绩会上已经披露——而是管理层在投资者日上展现出的两个关键态度：第一，他们对FY2028目标的实现有高度信心，甚至暗示实际数字可能更高；第二，他们详细披露了实现这些目标的具体路径和管道。对于习惯于日本企业保守指引的投资者来说，这种程度的明确性本身就是一种信号。

但真正值得深入挖掘的，是报告在字里行间透露出的几个结构性变化。这些变化不仅关乎松下这一家公司，也关乎整个AI基础设施供应链的竞争格局和投资逻辑。

1. 数据中心储能市场正在从“备用”走向“主动”，CBU将成为下一个增长奇点

松下对数据中心储能系统的目标非常具体：FY2028年实现约1万亿日元的销售收入，而FY2026年的指引仅为5500亿日元。这意味着两年内接近翻倍的增长。但更值得注意的是，松下认为这个目标已经包含了“已确认的产品开发和订单协议”，以及“部分延迟风险”。换句话说，这个目标不是乐观预测，而是基于已有订单的保守估计。

报告中真正的新信息，是松下对电容器备用单元（CBU）市场的展望。松下预计CBU市场在FY2028年将达到接近100亿日元，而到FY2030年将膨胀至1000亿日元。这个数字本身并不大——相对于万亿级别的整体储能市场——但它的增长斜率值得关注。CBU从2028年到2030年要增长10倍，这暗示的不仅是技术替代，而是整个数据中心电源架构的范式转变。

传统的数据中心备用电源以电池备用单元（BBU）为主，但CBU使用超级电容器，其优势在于响应速度更快、寿命更长、维护成本更低。随着AI训练集群的功率密度急剧上升，数据中心对电源的瞬时响应能力和可靠性要求也在提升。CBU不是要取代BBU，而是在“毫秒级响应”这个维度上填补BBU无法覆盖的空缺。

这意味着什么？对于投资者而言，CBU市场的爆发将直接利好上游的超级电容器供应商和材料厂商。松下计划在FY2026年开始与能源部门合作量产用于CBU的超级电容器，这比市场普遍预期的时间点更早。如果CBU市场真的按照松下的预期轨迹增长，那么相关的材料供应链——尤其是导电聚合物和先进电解质——将面临一次需求的结构性跃升。

2. 松下的AI材料业务利润率与BBU相当，这是一个被低估的盈利引擎

报告中的一个关键数据点：松下表示AI相关业务的调整后营业利润率约为20\%，与BBU业务相当。这个数字值得反复咀嚼。BBU业务作为成熟的数据中心电源产品，其利润率水平已经被市场充分认知。但AI材料业务——主要包括先进基板材料和导电聚合物电容器——也能达到同样的利润率，说明这个业务并非简单的“卖材料”，而是具有技术壁垒和定价权的。

更具体地说，松下在先进基板材料方面已经进入了M7和M8级别，而客户已经开始询问M9和M10。这揭示了一个重要的行业趋势：AI芯片的算力提升正在倒逼封装基板材料的升级。每一代M系列材料的升级，都意味着更高的信号传输速度、更低的损耗和更好的热管理能力。松下目前的产品结构中，M6及以下占60\%，M7和M8占40\%，但随着AI芯片进入下一个制程节点，M7和M8的占比将快速上升，而M9和M10的研发竞赛已经开始。

[... middle omitted ...]

另一句话来看，就变得意味深长了：“我们认为这一产能扩张表明，松下正在与电源设备制造商建立合作关系，同时在供应链中进行沟通。”这句话的潜台词是：松下的产能扩张不是盲目的，而是基于已确认的客户需求和合作框架。在AI基础设施供应链中，能够提前锁定产能分配的供应商，往往拥有更高的议价能力和更稳定的订单能见度。

这与其他日本电子零部件厂商的策略形成了对比。许多日本厂商在AI浪潮中采取了相对保守的产能扩张策略，担心过度投资导致产能利用率不足。但松下这次展现出的进取姿态，暗示管理层对AI基础设施需求的持续性有很强信心。这种信心本身就是一个值得关注的信号。

报告在讨论超级电容器时，提到松下“没有提及外部销售的可能性”。这是一个值得深究的细节。松下计划在FY2026年开始量产用于CBU的超级电容器，但这一产品目前似乎主要面向内部供应——即松下能源部门自己的CBU产品。如果松下未来决定向外部客户——比如其他电源系统集成商——销售超级电容器，那么其市场空间将远大于当前指引所暗示的规模。

\subsection{[R046] JPM：市场误读了英伟达进入基站射频单元的影响}
{\small\color{refgray}归类：未被主题摘要引用}

过去一周，诺基亚和爱立信的股价双双下跌，市场给出的解释高度一致：英伟达正在将AI芯片的触角延伸至基站中最核心、最定制化的射频单元（Radio Unit），这对传统设备商构成生存威胁。

但JPM在最新发布的研报中给出了一个直接对冲市场情绪的结论：这轮下跌是“事后合理化的价格反应”，而非基本面逻辑的必然结果。报告的核心判断是——英伟达的介入不会颠覆设备商的竞争格局，反而可能加速一个更值得关注的结构性变化：电信设备的价值链正在从“硬件定制决定性能”转向“软件差异决定溢价”。

这一判断之所以重要，是因为它触及了通信行业未来五年最根本的路线之争。传统上，基站设备商靠自研ASIC芯片实现性能和功耗的最优解，这是爱立信和诺基亚护城河的核心组成。但英伟达的方案提出了一个替代路径：用标准化的GPU平台承载更多基站功能，让设备商的差异化从芯片设计转移到软件算法。JPM认为，两条路径各有优劣，市场目前高估了英伟达对设备商的威胁，同时低估了诺基亚率先做出选择后可能获得的先发优势。

以下是我们从这份研报中提炼出的四个关键洞察，以及对一个尚未被充分讨论的风险的拆解。

JPM的分析起点是对技术事实的澄清。英伟达正在开发的是面向6G射频单元的GPU芯片，用于替代传统Massive MIMO天线中负责波束赋形等底层物理层功能的定制ASIC。技术驱动力是明确的：随着天线配置从4T4R向128T128R甚至1024T演进，射频单元所需的计算量呈32倍以上增长，可编程的GPU在处理灵活性和规模经济上开始具备理论优势。

但报告明确指出，这并不意味着设备商失去了选择权。诺基亚已经宣布将在6G基带单元（CU/DU）中标准化使用英伟达的AI加速器，并正在评估是否完全转向标准化半导体平台——即放弃自研ASIC，将差异化全部押注在软件层。而爱立信则明确表示，考虑到ASIC在性能上的固有优势，目前不计划放弃自研路线。

JPM认为，这两条路径各有合理的商业逻辑。标准化平台带来硬件升级成本的降低和软件更新的独立性，但永远无法达到定制芯片的极限性能；ASIC路线则相反。最终，电信运营商需要在“灵活性与性能”之间做出选择，而设备商的核心任务是押注并执行其中一条路线，而非被两条路线的拉扯所困。

这意味着，英伟达的进入只是为设备商增加了一个技术选项，而非剥夺了它们的战略自主权。市场将诺基亚和爱立信的股价下跌归因于英伟达，在JPM看来，是一种缺乏基本面支撑的后验归因。

报告中的一个关键细节值得放大：诺基亚与英伟达在基带单元上的合作已经进入执行阶段，这意味着诺基亚团队对英伟达CUDA架构和GPU编程模型的熟悉程度，远高于尚未启动类似评估的爱立信。

JPM认为，如果诺基亚最终决定将英伟达方案延伸至射频单元，它的实施速度将显著快于爱立信，因为后者需要对其现有解决方案进行更大幅度的架构调整。考虑到6G商用至少还有三年窗口期，这个时间差虽然不足以构成永久壁垒，但足以让诺基亚在6G标准制定初期就占据更有利的产业话语权。

这一点对投资者的含义是微妙的。市场当前将“诺基亚与英伟达合作”等同于“诺基亚放弃自研能力”，从而将其视为利空。但JPM的分析框架提示了一个相反的可能性：率先拥抱标准化平台，可能让诺基亚在软件差异化上积累更早的经验曲线，尤其是在AI原生网络功能（如智能资源调度、自动化运维）与通信功能的融合方面。

当然，这一判断的前提是诺基亚的软件团队确实能够兑现差异化价值。报告没有回避这一不确定性，但至少为市场提供了一个不同于主流叙事的观察视角。

英伟达GPU方案面临的两个核心挑战——功耗和性能——在报告中得到了克制但诚实的讨论。射频单元消耗了基站约90\%的网络能耗，任何新增计算单元都必须在这个约束条件下设计。JPM指出，英伟达有可能借鉴其汽车级嵌入式GPU的设计经验

[... middle omitted ...]

不确定的变量之一。报告没有给出明确的量化判断，这可能是投资者需要自行跟踪的关键验证点。

4. 真正需要关注的不是英伟达，而是设备商“软件溢价”能否被市场定价

如果我们将视线从英伟达进入射频单元的新闻本身移开，JPM这份报告指向了一个更深层的结构性问题：电信设备行业的估值框架是否需要重置。

在ASIC主导的时代，设备商的护城河是硬件性能+系统集成能力，估值更多反映的是市场份额和订单能见度。但如果行业转向标准化硬件+软件差异化的模式，设备商的竞争壁垒将从“别人造不出来的芯片”转变为“别人写不出来的算法”。后者的可复制性更低，但可验证性也更难——投资者需要判断的不再是芯片流片是否成功，而是软件迭代速度和组织能力。

这对诺基亚和爱立信的估值逻辑意味着什么？报告没有直接回答，但给出了一个重要的分析线索：如果诺基亚成功转向标准化平台，其硬件成本结构可能显著优化，而软件收入的占比和可预测性将提升。这理论上应该带来估值倍数的重估，但前提是市场能够看到软件差异化的实际落地案例。

\subsection{[R047] 摩根斯坦利：移动广告技术市场最被低估的增长引擎，不是需求，而是转化效率的杠杆效应}
{\small\color{refgray}归类：未被主题摘要引用}

摩根斯坦利：移动广告技术市场最被低估的增长引擎，不是需求，而是转化效率的杠杆效应

移动广告技术行业正处于一个关键拐点。多数投资者将注意力集中在广告主预算的周期性波动上，关注宏观环境如何影响总支出。但摩根斯坦利最新发布的移动广告技术入门报告揭示了一个更深层的结构性机会：真正的增长驱动力并非来自广告主多花多少钱，而是来自广告技术平台如何通过提升转化率来撬动自身收入。

这份报告的核心判断是：独立移动广告技术平台正站在一个“效率杠杆”的起点上——转化率每提升一个百分点，平台净收入增长的弹性远高于广告主支出同等幅度的增长。 当前市场对这一机制的理解仍不充分，尤其是当我们将头部平台（如Meta）与独立平台（如AppLovin、Unity）的转化率进行对比时，差距的规模本身就暗示了巨大的潜在上升空间。

为什么这个判断现在重要？因为市场正在两个方向上产生分歧。一方面，游戏内广告的存量增长正在放缓，独立平台需要寻找新的垂直领域。另一方面，电商、零售等非游戏领域的广告支出正在快速流入移动端。摩根斯坦利估算，独立应用内广告技术平台的可寻址市场在2025年约为790亿美元，到2030年将增长至1360亿美元，年复合增长率约为11\%。但真正有趣的问题不是这个市场有多大，而是哪些平台能够最有效地将流量转化为收入，以及这一转化效率本身如何成为竞争的护城河。

1. 转化效率的差异，定义了独立平台与围墙花园之间的估值鸿沟

理解移动广告技术行业，首先需要理解一个核心指标链条：点击率乘以转化率，决定了广告平台为广告主带来一次安装或购买所需消耗的展示次数。而这一效率，直接决定了平台能够从广告主支出中“截留”多少作为自己的净收入——也就是所谓的“抽成率”。

摩根斯坦利的分析提供了一个令人震惊的对比。报告估算，Meta的“真实转化率”大约是AppLovin的10倍。这个数字是怎么来的？对于AppLovin，其广告产生应用安装的转化率约为1.3\%，而游戏内付费用户的占比约为2.0\%，两者相乘得到的“真实转化率”仅为0.026\%。相比之下，Meta的转化率要高出一个数量级。

这一差距意味着什么？如果AppLovin的转化率能够向Meta的水平靠近，那么在不增加任何广告主支出的情况下，其净收入将出现数倍的增长。报告用一个简化的例子说明了这一机制：假设广告主为每次应用安装支付10美元，平台目前需要花费8美元购买20次展示才能产生一次安装，净收入为2美元，抽成率为20\%。如果转化效率提升，使平台只需15次展示就能产生一次安装，那么媒体成本降至6美元，净收入翻倍至4美元，抽成率升至40\%。

这是一个极其强大的杠杆。它意味着，对于独立广告技术平台而言，提升转化率比争夺更多广告主预算更具战略价值。而当前市场的定价，显然还没有充分反映这一杠杆的潜在影响。

2. 790亿美元的可寻址市场背后，隐藏着垂直扩张的结构性机会

摩根斯坦利将独立应用内广告技术平台的市场机会定位为790亿美元，这个数字仅仅是移动广告生态中的一部分。更宏观的数据是：全球移动广告支出约为5540亿美元，其中应用内广告支出约为3320亿美元。独立平台目前占据的份额仍然较小，但增长空间巨大。

报告指出，游戏仍然是独立平台最大的收入来源，2025年预计贡献约460亿美元。但真正值得关注的是非游戏领域的扩张。零售、非游戏娱乐（如流媒体）、教育、金融等垂直领域正在成为新的增长引擎。尤其值得注意的是，报告专门列出了“AI应用”这一类别，尽管2025年其规模为零，但到2030年预计将成为一个独立的、有意义的收入来源。

这种垂直扩张并非易事。报告在“增长可持续性”部分明确提出了一个核心辩论：电商和其他非游戏领域的广告支出能否规模化？这一问题目前尚无定论。但摩根斯坦利的分析框架表明，那些能够在多个垂

[... middle omitted ...]

学分配图：在典型的3美元CPM（千次展示成本）中，广告主和代理商支出3美元，其中需求方平台和供应方平台各收取约0.45美元，聚合平台收取0.06美元，最终出版商获得2.04美元。这意味着，平台方的抽成合计约为0.96美元，占总支出的32\%。但这一比例并非固定不变。

真正能够改变分配格局的，是平台能否证明自己不仅仅是流量通道，而是效果放大器。如果一家平台能够持续提升广告主的投资回报率，它就有能力提高抽成率，或者在不提高抽成率的情况下吸引更多预算。报告在“竞争”部分明确指出，新进入者（尤其是大型平台进入游戏内广告领域）可能对现有玩家的转化率和抽成率构成压力。这意味着，未来的竞争将不再是比谁拥有更多的流量，而是比谁能够更有效地将流量转化为广告主想要的结果。

摩根斯坦利的报告专门用了一个章节讨论信号、隐私和平台风险。这是移动广告技术行业最不确定的外部变量之一。随着苹果和谷歌持续推进隐私保护措施，第三方数据的使用受到越来越严格的限制，广告平台的定向能力和归因能力都在被削弱。

\subsection{[R048] 摩根斯坦利：市场低估的不是半导体设备周期性，而是结构性瓶颈的再定价}
{\small\color{refgray}归类：AI产业链：从GPU短缺到存储与计算架构再定价，企业级需求加速爆发}

摩根斯坦利：市场低估的不是半导体设备周期性，而是结构性瓶颈的再定价

当台北的Computex和AI Summit刚刚落幕，许多参会者带回的共识是：半导体设备板块的热度似乎从2025年底至2026年初的亢奋中冷却下来。投资者态度变得“更加平衡”。但这份来自摩根斯坦利日本团队的调研笔记，真正想传递的信号恰好相反——市场正在犯一个典型的“半杯水”错误。

会议期间，投资者对存储周期的担忧在现货价格回升后明显消退，对2027年前的存储景气度几乎没有悲观声音。但真正值得关注的，不是周期是否延续，而是三个正在重塑半导体设备行业定价权的结构性变量：Agentic AI对DRAM需求的量级跳跃、铜墙(Cu Wall)迫使光互连提前成为技术瓶颈，以及日本设备商在定价能力上的沉默但剧烈变化。

这份报告最核心的判断是：半导体设备行业正在从“周期成长”切换为“结构成长”，而市场对这一切换的定价仍不充分。

1. 一颗Vera CPU的需求量级，正在改写DRAM的供需方程

摩根斯坦利台湾半导体分析师Charlie Chan从供应链确认的数据，值得每一个关注AI资本开支的人重新审视。NVIDIA的Vera CPU单颗可搭载1.5TB的LPDDR5X，是上一代Grace CPU的三倍。以200-400万颗产量估算，仅Vera一个产品线，就需要消耗2025年全球DRAM bit总需求的10\%-16\%。

这不是一个边际增量。这是在一个原本被认为供需趋于平衡的市场中，突然出现的、由单一产品驱动的结构性缺口。换算成产能，这相当于10-20万片/月的1cnm晶圆产能需求。考虑到先进DRAM产能的扩产周期和良率爬坡速度，这个数字意味着未来2-3年内，存储设备订单的可预见性远高于市场共识。

摩根斯坦利的调研显示，投资者对存储周期持续性的信心已因现货价格回升而增强，但真正支撑信心的应该是这类终端需求的量级变化，而非价格信号本身。后者只是结果，前者才是原因。

Marvell在Computex的主题演讲中详细阐释了“铜墙”概念，这是理解未来2-3年后端设备需求的关键框架。铜线传输在100G/lane速率下可以覆盖5米，足以连接机架内部设备；但当速率提升至200G时，传输距离缩短至2.5米；到了400G——预计2028年左右开始规模部署——这一距离将进一步压缩至1.25米。

标准机架高度约为2米。这意味着，当数据传输速率进入400G时代，铜线将无法完整覆盖机架内部的互联需求。光互连不再是性能升级的选项，而是物理约束下的唯一解。

这个转变的直接含义是：CPO（共封装光学）将从 年起成为后端设备市场的主要主题。摩根斯坦利引用了Teradyne的测算，CPO测试设备TAM将从2026年的1亿美元，跃升至2028年的3-7亿美元，此后随着规模部署将突破10亿美元。

报告特别强调了Advantest在CPO测试中的潜在受益，以及Disco在混合键合技术中的角色。但更重要的是，这一趋势意味着后端设备市场的增长逻辑正在从“跟随先进封装”转向“驱动下一代互联架构”。设备商的议价能力和估值框架需要重新评估。

在摩根斯坦利日本峰会上，Kokusai Electric给出了一个看似保守的指引：下半财年营收环比下降17\%。但公司同时承认，这一指引存在上行空间——如果下半年环比持平，全年营收将从2800亿日元升至3000亿日元。

报告没有展开但值得深思的是，Kokusai的DRAM占比已从F3/23的18\%升至F3/26的44\%，而NAND占比从40\%降至12\%。这家公司正在从“NAND设备股”转型为“DRAM设备股”。而DRAM正是Agentic AI需求最直接受益的领域。

更重要的是，Kokusai在指引中对中国新兴公司和DRAM制造商的需求“仅谨慎纳入”。考

[... middle omitted ...]

率目标——当前摩根斯坦利预测为33\%——这将是一个新的估值锚点。

报告没有回答的一个关键问题是：这种定价权是东京电子独有的，还是整个日本设备行业的趋势？从Kokusai、Disco到Advantest，日本设备商在各自细分领域的技术壁垒正在加深。如果行业整体定价能力回归，那么当前基于周期估值框架的定价模型可能全面低估了这些公司的结构性盈利能力。

报告提到了一个值得关注的细节：在GPU测试领域，Advantest的优势“依然稳固”，且摩根斯坦利“预计不会出现100-0的局面”。这意味着市场此前对竞争格局的担忧可能过度了。

但报告也承认，在CPO测试领域，“竞争格局尚未确定”。这是整份报告中最需要后续跟踪的变量。如果CPO测试市场在2028年达到3-7亿美元，谁将主导这一市场？Advantest的技术领先能否转化为可持续的份额优势？Teradyne的1亿美元TAM预测是否过于保守？这些问题在报告中只是点到为止，但恰恰是决定2027年后设备板块估值分化程度的核心变量。

\subsection{[R049] 摩根斯坦利：欧洲猎头行业的真正风险不在需求，而在永久招聘的结构性萎缩}
{\small\color{refgray}归类：风险偏好：CLO市场从beta转向alpha，军工复合体面临工业基础结构性过时}

摩根斯坦利：欧洲猎头行业的真正风险不在需求，而在永久招聘的结构性萎缩

市场正在为欧洲人力资源服务公司的有机增长改善而松一口气，但这份摩根斯坦利的最新研报提出了一个更值得警惕的判断：临时招聘的回暖并不能掩盖永久招聘（perm recruitment）正在发生的结构性恶化，后者才是决定这些公司毛利率、经营杠杆和估值能否修复的真正变量。

这份报告发布于2026年6月，正值欧洲宏观环境因地缘政治不确定性再度升温之际。摩根斯坦利的分析师团队对欧洲人力资源板块进行了评级调整，其核心逻辑不是基于对经济周期的简单判断，而是基于对招聘业务结构本身的风险重估。

报告明确指出，尽管行业有机增长在过去几个季度有所改善，但改善主要来自临时招聘，而永久招聘的持续恶化正在侵蚀毛利率和经营利润。更关键的是，市场共识尚未充分反映这一风险。

这意味着什么？如果永久招聘的萎缩不是周期性的，而是结构性的——比如企业用人模式正在从“雇佣”转向“灵活用工”——那么当前估值所隐含的复苏预期可能过于乐观。

1. 临时招聘回暖是假象，永久招聘的恶化才是决定毛利率的关键

过去几个季度，欧洲人力资源公司的有机增长数据确实出现了改善迹象。但摩根斯坦利指出，股价对这些改善几乎没有反应。原因在于，投资者真正在等待的是毛利率的恢复和经营杠杆的改善，而这两者恰恰与永久招聘的占比高度相关。

永久招聘的毛利率远高于临时招聘。当一家公司从永久招聘中获得的收入占比下降时，整体毛利率就会承压。报告中的数据显示，Adecco的永久招聘占毛利润比例约为15-17\%，Randstad约为16\%，而Hays高达38\%，Page Group更是达到72\%。

这意味着，Page Group和Hays对永久招聘的依赖度远高于Adecco和Randstad。一旦永久招聘市场继续恶化，前者的毛利率和利润将受到更严重的冲击。

这就是为什么摩根斯坦利将Page Group的评级维持为Underweight，并将其目标价从195便士大幅下调至110便士。这不是一个简单的周期判断，而是对业务结构脆弱性的识别。

在宏观环境不确定的背景下，财务杠杆和股本稀释风险成为区分公司质量的关键维度。

报告特别指出了Adecco的风险。其净债务/EBITDA比率在2024年达到2.8倍，而Randstad同期仅为1.5倍。更值得关注的是，Adecco通过股票股息（scrip dividend）来维持分红，这实际上是在稀释现有股东的权益。在盈利下行周期中，这种操作虽然能暂时缓解现金流压力，但长期会损害每股收益。

相比之下，Randstad的杠杆水平更低，且没有显著的稀释风险。这就是为什么摩根斯坦利将Randstad的评级从Underweight上调至Equal-weight，同时将Adecco从Equal-weight下调至Underweight。

这一调整的逻辑非常清晰：在不确定的环境中，低杠杆、低稀释风险的公司拥有更大的安全边际和战略灵活性。当行业整体面临压力时，资产负债表的质量将成为决定谁能在下一轮周期中率先反弹的关键。

3. 报告尚未完全回答的关键问题：AI对招聘行业的长期影响到底有多大

摩根斯坦利在报告中提到了AI风险，但并未深入展开。报告指出，对AI威胁的担忧与短期宏观逆风一起，继续压制估值。

这里有一个尚未被充分讨论的问题：AI对人力资源服务行业的影响，可能不是简单的“替代”，而是重塑整个招聘价值链。

一方面，AI可以自动化简历筛选、初步面试和候选人匹配，这直接威胁到传统猎头公司的核心价值主张。另一方面，AI也可能创造新的需求——比如帮助企业评估员工的技能缺口、设计灵活用工方案等。

但问题在于，这些新需求能否弥补传统业务被侵蚀的损失？目前看来，答案并不明确。Page Group

[... middle omitted ...]

聘还是永久招聘。如果改善主要来自临时招聘，毛利率和经营杠杆的修复将非常有限。

第二，评估业务结构脆弱性。永久招聘占比越高的公司，在宏观不确定性下的下行风险越大。这不是一个周期判断，而是一个结构判断。

第三，关注资产负债表质量。在盈利下行周期中，高杠杆和股权稀释风险会放大下行空间。低杠杆、低稀释风险的公司拥有更强的抗压能力和战略灵活性。

第四，不要忽视AI的长期影响。AI对招聘行业的影响可能比市场预期的更深远。它不会均匀地影响所有公司，而是会加速行业内部的分化。

5. 一个尚未被市场定价的风险：永久招聘的复苏可能比预期更慢

当前市场共识可能隐含了一个假设：随着宏观经济改善，企业会重新增加永久招聘。但摩根斯坦利的数据和分析暗示，情况可能并非如此简单。

Indeed的职位发布数据在多个欧洲关键市场持续恶化，法国、德国、英国和美国的职位空缺指数都在下降。同时，欧洲的宏观和地缘政治不确定性仍在上升。这意味着，企业可能更倾向于保持灵活用工模式，而非大规模增加正式雇员。

\section{Disclaimer}
{\scriptsize\color{refgray}
This community is a paid learning and information-sharing community created for general educational purposes only. The membership fee is charged solely for access to community discussions, educational content, organization, and operational costs. It is not an advisory fee, management fee, performance fee, brokerage fee, or compensation for personalized financial, investment, legal, tax, or professional advice.\par
I participate and share information solely as an individual and for educational discussion among community members. I am not acting as a financial adviser, investment adviser, broker, dealer, portfolio manager, legal adviser, tax adviser, accountant, fiduciary, or any other licensed professional adviser. Nothing shared in this community constitutes, and should not be interpreted as, financial advice, investment advice, legal advice, tax advice, a recommendation, solicitation, offer, endorsement, instruction, or guarantee to buy, sell, hold, short, trade, or invest in any security, cryptocurrency, fund, derivative, strategy, or other financial product.\par
All content shared in this community, including messages, discussions, links, files, charts, examples, market commentary, watchlists, AI-generated or AI-polished summaries, and third-party materials, is general, impersonal, and educational in nature. It does not take into account any individual's financial situation, investment objectives, risk tolerance, investment horizon, liquidity needs, tax status, legal restrictions, or suitability requirements. No content should be relied upon as suitable for any specific person, account, portfolio, or financial decision.\par
Information shared in this community may be incomplete, inaccurate, outdated, speculative, AI-generated, AI-edited, or based on third-party internet sources that have not been independently verified. I make no representation or warranty as to the accuracy, completeness, timeliness, reliability, or suitability of any information shared.\par
Financial markets involve substantial risk, including the possible loss of principal. Past performance, historical data, hypothetical examples, simulated results, backtests, personal opinions, or educational examples do not guarantee future results. Each member is solely responsible for conducting their own research, exercising independent judgment, and making their own decisions. Before making any financial, investment, legal, or tax decision, members should consult qualified and licensed professionals.\par
Participation in this community, payment of any membership fee, or communication with me or other members does not create any adviser-client, fiduciary, professional, contractual, agency, partnership, employment, or other advisory relationship. By joining or participating in this community, you acknowledge and agree that you are solely responsible for your own decisions, actions, transactions, profits, losses, risks, and consequences, and that I shall not be liable for any direct, indirect, incidental, consequential, financial, legal, tax, or other loss or damage arising from or related to any content, discussion, information, or material shared in this community.\par
}
\end{document}
